% ---- preamble ----
\documentclass[12pt,reqno]{amsart}

\usepackage[utf8]{inputenc}
\usepackage{amsmath,amssymb,amsthm,mathtools,mathrsfs,esint}
\usepackage{mathalpha}
\usepackage{enumitem}
\usepackage{libertinus}
\usepackage{microtype}
\usepackage{booktabs,longtable,array}
\usepackage[a4paper,top=25mm,bottom=26mm,left=29mm,right=29mm,headsep=7mm]{geometry}
\usepackage[dvipsnames]{xcolor}
\usepackage[
  breaklinks=true,
  pagebackref=true,
  hypertexnames=false
]{hyperref}
\hypersetup{
  colorlinks=true,
  pdfpagemode=UseNone,
  citecolor=ForestGreen,
  linkcolor=MidnightBlue,
  urlcolor=RoyalBlue,
  pdfstartview=FitH,
  pdftitle={The lower-class escape rate of planar favorite sites},
  pdfsubject={Planar simple random walk, favorite sites, and lower-class escape rate}
}
\usepackage[nameinlink,noabbrev]{cleveref}

% Page geometry and paragraph spacing from the template.
\linespread{1.05}
\normalsize
\linespacing=\baselineskip
\setlength{\parindent}{1em}
\setlength{\parskip}{3pt}

\makeatletter

% First-page title, author block, abstract, and separate AI usage statement.
\def\@settitle{%
  \begin{center}%
    \uppercasenonmath\@title
    \Large\@title
  \end{center}}

\def\@setauthors{%
  \vskip21pt plus6pt minus3pt
  \begin{center}%
    {\tiny BY\par}\medskip
    {\normalsize\scshape\authors\par}\medskip
    \def\address##1##2{\par ##2\par}%
    \def\email##1##2{\par e-mail:~##2\par}%
    {\small\itshape\addresses\par}%
  \end{center}}

\def\enddoc@text{}

\renewenvironment{abstract}{%
  \global\setbox\abstractbox=\vtop\bgroup
  \vskip10pt
  \normalfont\tiny\centering\MakeUppercase{\abstractname}\par
  \Small
  \list{}{%
    \leftmargin3pc
    \rightmargin3pc
    \labelwidth0pt
    \listparindent\normalparindent
    \itemindent0pt
    \parsep0pt plus1pt
    \let\fullwidthdisplay\relax
  }%
  \item[]
}{%
  \vspace*{-10pt}
  \endlist\egroup
}

\newbox\aiusagebox
\newenvironment{aiusage}{%
  \global\setbox\aiusagebox=\vtop\bgroup
  \normalfont\Small
  \list{}{%
    \leftmargin3pc
    \rightmargin3pc
    \topsep0pt
    \partopsep0pt
    \parsep0pt
    \itemsep0pt
    \listparindent0pt
    \itemindent0pt
  }%
  \item[]%
  \nobreak\vspace{4pt}%
  {\normalfont\tiny\centering\MakeUppercase{AI usage}\par}%
  \nobreak\vspace{2pt}%
  \noindent\ignorespaces
}{%
  \par\endlist\egroup
}
\let\original@setabstracta\@setabstracta
\def\@setabstracta{%
  \original@setabstracta
  \ifvoid\aiusagebox\else
    \vskip10pt
    \box\aiusagebox
    \prevdepth\z@
  \fi
}

% Section headings.
\let\@secnumfont\@empty
\def\section{\@startsection{section}{1}{0pt}%
  {20pt plus4pt minus4pt}
  {.5\linespacing}
  {\normalfont\bfseries}}
\def\subsection{\@startsection{subsection}{2}{0pt}%
  {.5\linespacing plus.7\linespacing}
  {-.5em}
  {\normalfont\scshape}}
\def\subsubsection{\@startsection{subsubsection}{3}{0pt}%
  {.5\linespacing plus.7\linespacing}
  {-.5em}
  {\normalfont\itshape}}
\let\@afterindenttrue\@afterindentfalse
\@afterindentfalse

% Theorem and proof styles.
\newtheoremstyle{plain}
  {.5\linespacing}{.5\linespacing}
  {\itshape}{1em}{\scshape}{.}{.5em}
  {\thmname{#1}\thmnumber{ #2}\thmnote{\normalfont{} (#3)}}
\theoremstyle{plain}
\newtheorem{theorem}{Theorem}[section]
\newtheorem{lemma}[theorem]{Lemma}
\newtheorem{proposition}[theorem]{Proposition}
\newtheorem{corollary}[theorem]{Corollary}

\newtheoremstyle{definition}
  {.5\linespacing}{.5\linespacing}
  {\upshape}{1em}{\scshape}{.}{.5em}
  {\thmname{#1}\thmnumber{ #2}\thmnote{\normalfont{} (#3)}}
\theoremstyle{definition}
\newtheorem{definition}[theorem]{Definition}
\newtheorem{assumption}[theorem]{Assumption}

\newtheoremstyle{remark}
  {.5\linespacing}{.5\linespacing}
  {\upshape}{1em}{\scshape}{.}{.5em}
  {\thmname{#1}\thmnumber{ #2}\thmnote{\normalfont{} (#3)}}
\theoremstyle{remark}
\newtheorem{remark}[theorem]{Remark}

\renewenvironment{proof}[1][\proofname]{%
  \par
  \pushQED{\qed}%
  \normalfont
  \topsep6pt plus6pt\relax
  \trivlist
  \item[\hskip\labelsep\slshape #1\@addpunct{.}]%
  \ignorespaces
}{%
  \popQED
  \endtrivlist
  \@endpefalse
}
\DeclareRobustCommand{\qed}{%
  \leavevmode\unskip\penalty9999\hbox{}\nobreak
  \ifmmode
    \makebox[0pt][l]{\qquad\hbox{\qedsymbol}}%
  \else
    \qquad\hbox{\qedsymbol}%
  \fi
}
\renewcommand{\qedsymbol}{\rule[.47pt]{4pt}{6pt}}

\makeatother

\setlength{\emergencystretch}{1.8em}
\numberwithin{equation}{section}
\allowdisplaybreaks[1]

\DeclareMathOperator{\supp}{supp}
\DeclareMathOperator{\Var}{Var}
\DeclareMathOperator{\Cov}{Cov}
\DeclareMathOperator{\osc}{osc}
\DeclareMathOperator{\sech}{sech}

\title[Escape of planar favorite sites]{The lower-class escape rate of planar favorite sites}
\author{ GPT-5.6 Sol \and Heng Ma}
\thanks{Only Heng Ma take responsebility of the work.}
\date{}

\begin{document}

\begin{abstract}
We study the distance from the origin of the set of favorite sites of planar
simple random walk. The logarithmic-power boundary for its almost-sure lower
class is $\sqrt{n/\log n}$: above exponent $1/2$, every favorite eventually
escapes each fixed multiple of $\sqrt n/(\log n)^\gamma$, whereas at exponent
$1/2$ the entire favorite set enters arbitrarily small multiples of this
scale infinitely often. The escape argument combines a localized upper
bound for thick points with a lower bound for the global maximum at
$J\log J$ precision. Traversal barriers, relative lattice comparisons and
conditional terminal occupation estimates preserve the required logarithmic
power. Critical recurrence follows from a stopped-resolvent estimate,
reversal at a priority-selected favorite, and an injective insertion that
removes ties.
\end{abstract}

\begin{aiusage}
This manuscript was substantially generated and subsequently checked by AI
systems. At the time of this version, the arguments and citations have not yet been
independently verified in full by the human contributor. \textbf{A digested and more streamlined version   will replace this preliminary version on both Hexagon and arXiv in the coming months.}
\end{aiusage}

\maketitle

\begingroup
  \raggedbottom
  \vspace*{-20pt}
  \setcounter{tocdepth}{1}
  \tableofcontents
  \clearpage
\endgroup


% ---- introduction ----
\section{Introduction}\label{sec:introduction}

A favorite site of a random walk is a site at which the accumulated local time
is maximal. Although the planar walk returns to the origin infinitely often,
its favorite sites need not remain near the origin. The question considered
here is how small a neighborhood of the origin can still contain a favorite
at arbitrarily large times. Because several sites can tie for the maximum,
there are two distinct observables: the distance of the nearest favorite and
the distance of the farthest favorite.

Let $(S_n)_{n\ge0}$ be simple symmetric random walk on $\mathbb Z^2$, started
at the origin. Define
\begin{equation}\label{eq:local-time-definition}
 L_n(x)=\sum_{j=0}^n\mathbf1_{\{S_j=x\}},\qquad
 M_n=\max_{x\in\mathbb Z^2}L_n(x),\qquad
 \mathcal F_n=\{x\in\mathbb Z^2:L_n(x)=M_n\}.
\end{equation}
The set $\mathcal F_n$ is finite and nonempty. Put
\begin{equation}\label{eq:radii-definition}
 R_n^- =\min_{x\in\mathcal F_n}\|x\|_2,\qquad
 R_n^+ =\max_{x\in\mathcal F_n}\|x\|_2,
 \qquad
 a_\gamma(n)=\frac{\sqrt n}{(\log n)^\gamma},\quad n\ge3.
\end{equation}
For example, along the path $(0,e_1,0,e_1)$ the two sites tie, so
$R_3^-=0$ and $R_3^+=1$. Thus localizing an arbitrarily chosen favorite
does not by itself localize the whole favorite set.

\begin{theorem}[Lower-class escape rate]\label{main:precise-escape-rate}
Almost surely, simultaneously for every real $\gamma$,
\begin{align}
 \liminf_{n\to\infty}\frac{R_n^-}{a_\gamma(n)}
 &=\infty, &&\gamma>\tfrac12,\label{eq:main-escape}\\
 \liminf_{n\to\infty}\frac{R_n^+}{a_\gamma(n)}
 &=0, &&\gamma\le\tfrac12.\label{eq:main-recurrence}
\end{align}
\end{theorem}

In particular, for each $\gamma>1/2$ and each fixed $A>0$, eventually
every favorite lies outside $B(0,Aa_\gamma(n))$. At the critical exponent,
for every $\varepsilon>0$ there are infinitely many times at which
\emph{all} favorites lie in $B(0,\varepsilon\sqrt{n/\log n})$.
The conclusion identifies the boundary among logarithmic powers of the
diffusive scale. It does not assert an integral test for arbitrary slowly
varying perturbations of the critical gauge.

The corresponding one-dimensional theorem has critical exponent $1$;
see Bass~\cite[Theorem~2.2]{Bass2023}. The planar escape-rate question is
posed in Hao, Li, Okada and Zheng~\cite[Remark~1.4(1) of the first arXiv
version]{HLOZ2025}. Our focus is the location of the full set of maximizers.
The two assertions in Theorem~\ref{main:precise-escape-rate} therefore retain
the distinction between its nearest and farthest members throughout the proof.

\subsection{Notation and the two local-time estimates}

All logarithms are natural. We write $P_x$ and $E_x$ for the law and
expectation of a walk started at $x$, and omit the subscript when $x=0$.
Its transition operator is
\[
 Pf(x)=\frac14\sum_{y:\|y-x\|_2=1}f(y).
\]
By symmetry the same notation acts on lattice measures. For $r>0$ let
$D(x,r)=\{y\in\mathbb Z^2:\|y-x\|_2<r\}$ and
$\tau_r=\inf\{n\ge0:S_n\notin D(0,r)\}$. Closed balls are denoted by
$B(x,r)$. We use the same geometric notation in the plane when the ambient
space is specified. A local time at a stopping time includes that time;
inside a killed domain it agrees with occupation strictly before exit.

The Green-function normalization and the local-time levels are
\begin{equation}\label{eq:level-definition}
 g=\frac2\pi,\qquad
 u_{J,c}=2gJ^2-cgJ\log J.
\end{equation}
The logarithmic radius is $J$, so the corresponding time scale is $e^{2J}$.
The coefficient of $J\log J$ will matter: knowing only the leading
$2gJ^2$ term cannot distinguish the powers in
Theorem~\ref{main:precise-escape-rate}.

The escape argument uses the following two estimates, proved in
Propositions~\ref{prop:localized-disk-upper} and
\ref{prop:disk-maximum-lower}. For fixed $\beta>0$, $c>2$ and $A>0$,
\begin{equation}\label{eq:intro-localized-upper}
 P\left(\max_{\|x\|_2\le e^JJ^{-\beta}}
 L_{\tau_{e^J}}(x)\ge u_{J,c}\right)
 \le J^{-2\beta+c-2+o(1)}+O(J^{-A}).
\end{equation}
For every fixed $\eta>0$, there is $p_\eta>0$ such that
\begin{equation}\label{eq:intro-disk-lower}
 P(M_{\tau_{e^J}}\ge u_{J,2+\eta})\ge p_\eta
 \quad\text{for all sufficiently large }J.
\end{equation}
The latter estimate is amplified over independent segments to give an
eventual almost-sure lower bound at deterministic exponential times.
The former then excludes a favorite near the origin throughout each time
block. Recurrence at the critical exponent uses a separate geometric-time
localization argument and does not depend on either disk estimate.

Constants denoted by $c,C$ may change between occurrences. Dependence on a
fixed parameter is indicated when relevant. Integer parts and fixed changes
of radius are harmless in estimates with a $J\log J$ margin: replacing $J$
by $J+O(1)$ changes $u_{J,c}$ by $O(J)$. Directed lattice crossings and
conditional bridge laws are defined exactly when they are introduced.



% ---- proof_ideas ----
\section{Proof ideas: obstacles and the estimates that overcome them}
\label{sec:proof-ideas}

The two halves of the theorem require different information. To prove escape,
one must show that a small disk cannot compete with the global local-time
maximum. To prove critical recurrence, one must bring the entire set of
maximizers into a small disk despite all local time accumulated in the past.
We explain the quantitative mechanisms before developing their proofs.

\subsection{Why a site-by-site bound misses the escape threshold}

The local time at a site, after the first hit and before leaving a disk of
radius $e^J$, has a geometric tail with mean $gJ+O(1)$. At level $u_{J,c}$
this gives
\[
 P_x(L_{\tau_{e^J}}(x)\ge u_{J,c})
 \le C\exp\{-u_{J,c}/(gJ)+O(1)\}
 =Ce^{-2J}J^c
\]
for sites in a fixed inner fraction of the disk. A union bound over
$e^{2J}J^{-2\beta}$ sites therefore gives only $CJ^{c-2\beta}$.
Even incorporating the hitting probability does not give the needed
estimate. In a logarithmic shell at distance about $e^{J-k}$, that
probability is at most $C(k+1)/J$, and the number of sites is
$O(e^{2J-2k})$. Summing over $k\ge\beta\log J+O(1)$ yields
\begin{equation}\label{eq:ideas-naive-upper}
 C J^{c-1}\sum_{k\ge\beta\log J+O(1)}(k+1)e^{-2k}
 \le J^{-2\beta+c-1+o(1)}.
\end{equation}
At the relevant level $c\downarrow2$, summability of this bound over
exponential time blocks would require $\beta>1$. The desired conclusion
requires the extra factor $J^{-1}$ in
\eqref{eq:intro-localized-upper}.

The improvement comes from recording how local time is built across nested
annuli. Let $T_j$ count inward traversals at logarithmic level $j$.
For the ideal concentric Brownian model, descendants of completed source
excursions form a critical Galton--Watson process with geometric offspring.
The square-root coordinate $\sqrt{2T_j}$ has an approximately linear
profile on an extreme traversal event. We restrict this profile below
\[
 B_j=\rho_Jj+z+(j\wedge(J-j))^{1/4},\qquad
 \rho_J=2-\frac{\log J}{J},
\]
where $z>0$ is fixed. Early crossings are controlled separately by a
coarser net of representatives. On the event with no crossing, the
Galton--Watson ballot estimate retains the cost of staying below the barrier,
rather than merely reaching a high endpoint.

To see where that cost enters, stop the traversal description at
$n=J-h$, with terminal radius $M=e^h$. Consider a cell in source shell
$\beta\log J+O(1)\le k\le C_0\log J$. The joint event that the prefix
stays below the barrier and $\sqrt{2T_n}\in[q,q+1]$ has probability at most
\begin{equation}\label{eq:ideas-ballot}
 C(k+1)^C(1+p)n^{-2}e^{-q^2/(2n)},\qquad
 p=B_n+1-q.
\end{equation}
This statement applies to the central range $\epsilon_0n\le q\le B_n+1$;
it is not a probability conditional on the prefix being good.
Extinction has substantial probability and is treated separately. The
factor $n^{-2}$ in \eqref{eq:ideas-ballot} follows by combining the
source law with the square-root prefactor in the ballot theorem. The
application-specific identity is
\[
 \frac{x^2}{m}+\frac{(q-x)^2}{n-m}
 =\frac{q^2}{n}+\frac{n(x-mq/n)^2}{m(n-m)},\qquad m=k+1.
\]
It combines the cost of the source height $x$ with the cost of its
descendants without losing an exponential factor.

\subsection{Keeping the terminal local times on the correct conditional law}

The traversal endpoint does not determine the local times inside the
terminal cell. We erase each completed outward path from radius $M$ to
radius $eM$, retaining its entrance and directed exit edge. Given the
exterior path and these endpoints, the erased paths are independent killed
bridges. If $b$ is their number and $X_i(x)$ their occupation at $x$, the
remaining quantity is $\max_x\sum_{i\le b}X_i(x)$.

Conditioning on a complete radial word also constrains the future of the
endpoint sequence. Consequently, a one-step harmonic-measure calculation
cannot simply be declared invariant under that conditioning. The proof
factors the entire word likelihood, including the last suffix and the
connectors between source intervals. It compares the true law with a law
obtained by normalizing the successive kernels in chronological order.
Uniform contraction controls dependence within each centered interval;
a fixed number of intervening connector kernels is allowed. This gives
both concentration of bridge moments and the exponential estimate needed
for the upper bound.

The central exact identity is a rank-one Green-function formula. For one
bridge with endpoint pair $\omega$, set $m_x(\omega)=E^\omega X(x)$ and
write $G$ for the Green function of its killing disk. Then
\[
 E^\omega e^{\theta X(x)}=1+r_\theta(x)m_x(\omega),\qquad
 r_\theta(x)=\frac{e^\theta-1}
 {1-(G(x,x)-1)(e^\theta-1)}.
\]
Averaging over the full endpoint law gives, to a controlled error,
$\log E[\exp\{\theta\sum_iX_i(x)\}\mid w]=bg r_\theta$.
For the upper bound choose $h=(C_h\log J)^4$ and initially
$r_0=(gn)^{-1}$. Since $G(x,x)=gh+O(1)$, the corresponding tilt satisfies
$\theta(r_0)u_{J,c}=2J-c\log J+O(1)$. A union bound over the $O(e^{2h})$
sites in one cell gives
\[
 P\left(\max_x\sum_iX_i(x)\ge u_{J,c}\mid w\right)
 \le Ce^{-2n}J^c e^{b/n}.
\]
Because $b=q^2/2+O(n)$, the factor $e^{b/n}$ cancels the exponential
in \eqref{eq:ideas-ballot}. This cancellation explains why the traversal
endpoint and the terminal maximum must be estimated together.

The displayed bound alone is insufficient: summing all endpoint bands
would lose powers of $J$. Varying the tilt around $r_0$ supplies the
additional factor
\[
 \exp\left\{-\kappa\min\left(
 \frac{(p-p_*)^2}{h},\frac{J^2}{h}\right)\right\},
 \qquad p_*=O(\log J).
\]
Its central sum, including the factor $1+p$, is
$O((1+|p_*|)\sqrt h+h)=J^{o(1)}$; the saturated part is negligible.
Thus one original cell contributes at most
$(k+1)^Ce^{-2n}J^{c-2+o(1)}$. There are $O(e^{2(n-k)})$ such cells in
shell $k$. Summing from $k=\beta\log J+O(1)$ produces exactly
$J^{-2\beta+c-2+o(1)}$. The remaining shells $k>C_0\log J$ contain
few enough sites that the elementary hitting and geometric-tail bounds
suffice, after choosing $C_0$ large.

\subsection{Producing a sufficiently large global maximum}

For the lower bound, the traversal construction must produce a cell with
\begin{equation}\label{eq:ideas-coarse-count}
 b=2L^2-2L\log L+O(L),\qquad L=J-h.
\end{equation}
Rosen's barrier construction provides such a root with positive probability
in the spherical Brownian model. Its angular restrictions make the
two-point sum finite. The passage to the lattice has to preserve this
existence event: an additive coupling error cannot be summed over
exponentially many rare candidate roots. Accordingly, the strong coupling
is applied to the union-existence event in one direction. The Brownian
pair estimate is used before the transfer. The lattice proof then uses a
first-moment bound and a separate conditional estimate for each root.

The last scales supply a contribution that cannot be discarded.
Conditionally on regular endpoint moments, the normalized sum of centered
bridge occupations has covariance $2gG(x,x')+O(1)$. A Gaussian
approximation and the maximum estimate for the discrete Gaussian free field
give
\begin{equation}\label{eq:ideas-terminal-lower}
 \max_x\sum_{i\le b}X_i(x)
 =gb+2\sqrt2\,g\sqrt b\,h
 +O_P\left(J\sqrt h+\frac{J\log J}{\log\log J}\right).
\end{equation}
Inserting \eqref{eq:ideas-coarse-count} shows the role of the second term:
\begin{align*}
 gb&=2gJ^2-4gJh-2gJ\log J
       +O(h^2+h\log J+J),\\
 2\sqrt2\,g\sqrt b\,h&=4gJh+O(h^2+h\log J).
\end{align*}
The terminal fluctuation restores the $4gJh$ lost by stopping the
traversal description early. The result is
$2gJ^2-2gJ\log J+o_P(J\log J)$.

Here the terminal scale is $h=(\log J)^p$ with $1<p<2$.
The inequality $p>1$ makes $M=e^h$ large enough for the relative
discretization errors. The inequality $p<2$ gives
$J\sqrt h=o(J\log J)$ and makes simultaneous moment concentration over
$O(M^4)$ observables possible. This is a different choice from the upper
bound, where no simultaneous screen of all these moments is needed.

Positive-probability disk success is enough. Partition a time interval of
length $e^{2j}$ into $j^a$ independent segments. Their effective logarithmic
radius is $j-(a/2)\log j+O(1)$, which costs $2a$ in the coefficient of
$gj\log j$. Choose $a>0$ within the available margin. Failure in every
segment has probability at most $e^{-cj^a}$, a summable bound. This yields
the deterministic-time lower estimate used for escape.

\subsection{Localizing the entire favorite set at the critical scale}

For critical recurrence, work at an independent geometric time $K$ with
$P(K=n)=(1-z)z^n$ and put $N=(1-z)^{-1}$. Choose a favorite by independent
continuous site priorities and reverse the path up to its last visit.
The reversed path consists of completed excursions from the selected site
and one terminal arm. Once the excursions, the forward suffix and the
priorities are fixed, being the selected favorite is exactly a collection
of upper bounds on the terminal arm's occupation at the other sites.
The arm is therefore a walk stopped when it returns to its root or exceeds
one of these capacities.

Its endpoint law is a positive mixture of stopped occupation measures.
Such a measure $\mu$, of mass $\beta$, satisfies
\[
 (I-zP)\mu=\alpha\delta_0-\nu,\qquad \nu\ge0,\qquad
 \alpha-\nu(\mathbb Z^2)=(1-z)\beta.
\]
Testing against a resolvent of the autocorrelation of a square gives
\begin{equation}\label{eq:ideas-resolvent-gain}
 \mu(B(0,C\sqrt N/f))
 \ge c\beta\frac{1+\log f}{f^2}.
\end{equation}
The Fourier transform of that test function is nonnegative, so its
resolvent is maximal at the origin. This makes the unknown nonnegative
exit measure $\nu$ usable. The two-dimensional integral is
$R^2(1+\log(N/R^2))$, which gives the logarithmic gain in
\eqref{eq:ideas-resolvent-gain} when $R\asymp\sqrt N/f$.
A restriction depending only on the exposed pieces makes the radius
proportional to $\sqrt K/f$ rather than $\sqrt N/f$.

Two remaining obstacles concern maximizers themselves. First, the selected
favorite can have distant ties. Insert the fixed loop
$(x,x+e_1,x,x+e_2,x)$ after its last visit. It raises the local time at $x$
by two and at each affected neighbor by one, making $x$ unique. The loop
is recovered from the last visit to the new unique favorite, so this is
an injection with a fixed positive probability cost. Second, a fresh
segment does not erase the old local-time profile $H$. If its own profile
$\ell$ has unique maximizer $y$, then
\[
 \operatorname*{arg\,max}(H+\ell)
 \subset\operatorname{supp}H\cup\{y\}.
\]
Thus it suffices to make the old range small on the new spatial scale;
the old local-time heights need not be bounded.

Choose block starting times $s_k=e^{J_k}$, with
$J_{k+1}-J_k=10\log J_k$, and take $f_k\asymp\sqrt{J_{k+1}}$.
The probability of a localized fresh singleton is at least
\[
 d_k\asymp\frac{\log J_k}{J_k},\qquad J_k\sim10k\log k.
\]
Hence $\sum_kd_k=\infty$. The logarithmic factor improves the available
success bound to order $1/k$; the weaker bound of order $1/J_k$ would
also have divergent sum for this particular schedule. The role of the
block choice is to make the old range negligible while preserving
divergence. More precisely, take
$N_k=e^{J_{k+1}}/J_{k+1}^2$ and retain clocks
$K_k\ge(d_k/8)N_k$. The old range is eventually at most
$\sqrt{s_k}J_k$, while $s_{k+1}/s_k\asymp J_k^{10}$ gives
\[
 \frac{\sqrt{s_k}J_k}{\sqrt{K_k/J_{k+1}}}
 \le\frac{C}{J_k^2\sqrt{\log J_k}}\longrightarrow0.
\]
At a smaller scale with $f_k\asymp J_k^\gamma$ and
$\gamma>1/2$, the probabilities furnished by the same estimate would
instead have summable lower bounds. This observation describes the
reach of this construction; the separate upper estimate proves escape.

Finally, the escape estimate with $c=2+2\delta$ gives a block probability
$j^{-2\gamma+2\delta+o(1)}$. For $\gamma>1/2$ choose $\delta>0$ so that
$2\gamma-2\delta>1$ and apply Borel--Cantelli. At $\gamma=1/2$, the
independent recurrence construction gives arbitrarily small normalized
favorite-set radii. These are the two summability calculations that place
the boundary at $1/2$.


% ---- elementary ----
\section{Elementary estimates}\label{sec:elementary}
For a finite lattice domain $D$, write
\[
 G_D(x,y)=E_x\sum_{j<\tau_{D^c}}\mathbf1_{\{S_j=y\}},
 \qquad \tau_{D^c}=\inf\{j\ge0:S_j\notin D\}.
\]
The following estimates fix the normalization used throughout the paper.
\begin{lemma}[Elementary walk estimates]\label{lem:elementary-walk-estimates}
There are constants \(c,C\in(0,\infty)\) with the following properties.

\begin{enumerate}\item If \(\tau_R=\inf\{n:\lVert S_n\rVert_2\ge R\}\), then
   \[
       \sup_{\lVert x\rVert<R}E_x\tau_R\le C R^2,\qquad
       \sup_{\lVert x\rVert<R}P_x(\tau_R>mCR^2)\le 2^{-m}.
   \]
\item For all \(N,R\ge1\),
   \[
       P\left(\max_{k\le N}\lVert S_k\rVert_2\ge R\right)
       \le C\exp\{-cR^2/N\}.
   \]
\item Uniformly for \(1\ll r<\lVert x\rVert<R\),
   \[
   P_x(\tau_{\partial D(0,R)}<\tau_{\partial D(0,r)})
     =\frac{\log(\lVert x\rVert/r)+O(r^{-1})}{\log(R/r)}
   \]
   and the complementary inner-hitting probability is obtained by replacing
   the numerator by \(\log(R/\lVert x\rVert)+O(r^{-1})\).
\item For each fixed $\kappa\in(0,1)$, uniformly for $\|x\|\le(1-\kappa)R$ in
   \(D(0,R)\),
   \[
       G_{D(0,R)}(x,x)=g\log R+O_\kappa(1).
   \]
   Conditional on first hitting \(x\), the number of visits to \(x\) before
   leaving the disk is geometric with this mean.
\end{enumerate}
\end{lemma}

\begin{proof}
For (1), optional stopping of
\(\lVert S_{n\wedge\tau_R}\rVert_2^2-(n\wedge\tau_R)\), followed by the
one-step overshoot bound, gives the expectation estimate.  Markov's
inequality at time \(2CR^2\), iterated with the strong Markov property,
gives the geometric tail.  Part (2) follows by applying the usual
exponential martingale separately to both coordinates and to both signs.

For the annulus estimate, we use \cite[Lemma~2.1, (2.5)]{HLOZ2025}:
uniformly for $r<\|x\|<R$,
\[
 P_x(H_{\partial D(0,R)}<H_{\partial D(0,r)})
 =\frac{\log(\|x\|/r)+O(r^{-1})}{\log(R/r)}.
\]
Here \(H\) is the positive-time hitting time and the boundaries are the
digital outer vertex boundaries defined in that paper.  The later
directed-edge convention can differ from this vertex-boundary convention
only within a one-unit radial collar.  Applying the estimate with \(r,R\)
shifted by at most two units, equivalently repeating its potential-kernel
proof at the directed crossing time, changes the logarithmic terms by
\(O(r^{-1})\).  Hence the same estimate holds for the directed-edge rows
below.  The potential-kernel proof of that result also gives (4).  The
geometric assertion in (4) is the strong Markov
property at successive positive returns to \(x\).
\end{proof}


% ---- recurrence ----
\section{Recurrence at the critical scale}\label{sec:recurrence}
The goal of this section is to place every favorite in a small disk at infinitely
many times. We first localize one favorite at an independent geometric time,
then make that favorite unique by an injective four-step insertion. Uniqueness
allows the resulting fresh segment to be combined with an arbitrary past.

\subsection{A resolvent estimate and reversal at a favorite}

\begin{lemma}[A stopped resolvent bound]\label{lem:stopped-resolvent}
Let \(0<z<1\), \(N=(1-z)^{-1}\), and let \(\mu\) be a finite positive
measure on \(\mathbb Z^2\) of mass \(\beta\).  Suppose that
\begin{equation}
(I-zP)\mu=\alpha\delta_0-\nu,\qquad \nu\ge0,
       \qquad \alpha-\nu(\mathbb Z^2)=(1-z)\beta.                 \label{eq:rec-1}
\end{equation}
There is a universal \(c>0\) such that, whenever
\(2\le f\le N^{1/8}\),
\begin{equation}
\mu\left(B\left(0,C\frac{\sqrt N}{f}\right)\right)
       \ge c\beta\,\frac{1+\log f}{f^2}.                         \label{eq:rec-2}
\end{equation}
\end{lemma}

\begin{proof}
Let \(Q_R=[-R,R]^2\cap\mathbb Z^2\) and set
\[
 w_R:=|Q_R|^{-1}{\bf1}_{Q_R}*{\bf1}_{-Q_R}.
\]
Then \(0\le w_R\le1\), \(w_R(0)=1\), and \(w_R\) is supported in
\(B(0,3R)\).  Put
\[
 G_z:=\sum_{n\ge0}z^nP^n,\qquad \phi_R:=G_z*w_R.
\]
Fourier inversion gives
\[
 \widehat{\phi_R}(\theta)=
 \frac{|\widehat{{\bf1}_{Q_R}}(\theta)|^2/|Q_R|}
      {1-z\widehat P(\theta)}\ge0.
\]
Consequently \(\phi_R(y)\le\phi_R(0)\) for every \(y\).  Pairing \eqref{eq:rec-1} with
\(\phi_R\) and using \((I-zP)\phi_R=w_R\) yields
\begin{equation}
\begin{split}
 \mu(w_R)
 &=\alpha\phi_R(0)-\nu(\phi_R)\\
 &\ge(\alpha-\nu(\mathbb Z^2))\phi_R(0)
   =(1-z)\beta\phi_R(0).                                      \label{eq:rec-3}
\end{split}
\end{equation}
The elementary Fourier integral over
\(\{\lVert\theta\rVert\le c/R\}\), together with
\(1-z\widehat P(\theta)\asymp (1-z)+\lVert\theta\rVert^2\),
gives, for \(R^2\le N/4\),
\begin{equation}
\phi_R(0)\ge cR^2\left(1+\log\frac N{R^2}\right).              \label{eq:rec-4}
\end{equation}
Since \(\mu(w_R)\le\mu(B(0,3R))\), take
\(R=\lfloor\sqrt N/(4f)\rfloor\) in \eqref{eq:rec-3}--\eqref{eq:rec-4}.  This proves \eqref{eq:rec-2},
after changing the absolute radius constant.
\end{proof}

\begin{lemma}[Localization of a selected favorite]\label{lem:plateau-bijection}
Let \(K\) be independent of the walk with
\[
       P(K=n)=(1-z)z^n,\qquad n\ge0.
\]
Attach independent uniform $[0,1]$ priorities \((U_x)_{x\in\mathbb Z^2}\) to
the sites and let \(X_K^*\) be the least-priority member of
\(\mathcal F_K\).  If \(\sigma\) and \(\ell\) are respectively its first and
last visit times, then the endpoint law of \(-X_K^*\), restricted by any event measurable with respect to the base tuple defined in the proof, is a positive
mixture satisfying \eqref{eq:rec-1}.  In particular,
\begin{equation}
P\left(\lVert X_K^*\rVert_2\le C\frac{\sqrt K}{f}\right)
       \ge c\,\frac{1+\log f}{f^2},\qquad
       2\le f\le N^{1/8}.                                     \label{eq:rec-5}
\end{equation}
\end{lemma}

\begin{proof}
We give the atomwise bijection; this is needed because a mere informal path
reversal does not control plateaux of tied favorites.

\smallskip\noindent\underline{\textit{Step 1.}}\quad Decompose and reconstruct the path.

Fix a clock atom \(K=n\), a path
\(s=(s_0,\ldots,s_n)\), and first insert
\[
       1=\sum_x{\bf1}_{\{x=X_n^*\}}.
\]
Thus, in the sector currently considered, \(x\) is deterministic while the
priorities are translated.  Put \(m=M_n\), \(r=m-1\), and let
\(\sigma,\ell\) be the first and last visits to \(x\).  Reverse and center
the prefix through \(\ell\):
\[
       A_i=s_{\ell-i}-x,\qquad 0\le i\le\ell.
\]
It has exactly \(r\) positive returns to zero by time
\(\ell-\sigma\).  Hence it decomposes uniquely as
\begin{equation}
A=E_1\oplus\cdots\oplus E_r\oplus C,                    \label{eq:rec-6}
\end{equation}
where the \(E_i\)'s are completed positive root excursions and \(C\)
starts at zero, has no positive return to zero, and ends at
\(a=-x\).  The centered forward suffix
\[
       B_i=s_{\ell+i}-x,\qquad 0\le i\le n-\ell,
\]
also has no positive return to zero.

Write \(H_r(y)\) for the number of visits to \(y\ne0\) in the completed
excursions and \(N_C(y),N_B(y)\) for the corresponding positive-time
counts in \(C,B\).  Translate priorities by the fixed sector shift,
\(V_y=U_{x+y}\).  The statement that the root is the priority-selected
favorite is exactly
\begin{equation}
H_r(y)+N_C(y)+N_B(y)
       \le r+{\bf1}_{\{V_y>V_0\}},\qquad y\ne0.                 \label{eq:rec-7}
\end{equation}
Indeed, a better-priority site must be strictly below the root's value
\(r+1\), while a worse-priority site may tie it.

Expose the base tuple
\[
       \mathcal{B}=(r,E_1,\ldots,E_r,B,V)
\]
and define its fixed capacities
\begin{equation}
c_\mathcal{B}(y):=
       r+{\bf1}_{\{V_y>V_0\}}-H_r(y)-N_B(y),\qquad y\ne0.       \label{eq:rec-8}
\end{equation}
Let \(\tau_c\) be the first positive return of a new walk \(C\) to zero or
the first time some \(N_C(y)\) exceeds \eqref{eq:rec-8}.  The possible terminal arms
are exactly the prefixes \(C_0,\ldots,C_q\) with \(q<\tau_c\).

Conversely, a feasible base tuple and one such prefix reconstruct the
original path: concatenate \eqref{eq:rec-6}, reverse it, translate by
\(-a=x\), and append the translated suffix \(B\).  Inequality \eqref{eq:rec-7} makes
\(x\) a favorite and excludes every better-priority tie; avoidance of zero
by \(B\) makes \(\ell\) its last visit.  The endpoint \(a\) recovers the
sector \(x=-a\), so both the path and the priority translation are unique.
The forward and inverse maps are therefore genuine bijections, including
\(r=0\).

\smallskip\noindent\underline{\textit{Step 2.}}\quad Identify the stopped occupation measure.

If \(\rho_r\), \(q\), and \(k\) are the lengths of the completed part,
\(C\), and \(B\), reversal and symmetry preserve the path weight and
\begin{equation}
z^{\rho_r}4^{-\rho_r}\,
 z^q4^{-q}\,(1-z)z^k4^{-k}
       =(1-z)z^n4^{-n}.                                      \label{eq:rec-9}
\end{equation}
Thus no multiplicity or normalization factor is hidden.

For a fixed capacity let
\[
       g_c(u)=E_0\sum_{q<\tau_c}z^q{\bf1}_{\{C_q=u\}}.
\]
Direct telescoping gives
\begin{equation}
(I-zP)g_c=\delta_0-E_0[z^{\tau_c}\delta_{C_{\tau_c}}].  \label{eq:rec-10}
\end{equation}
Mixing \eqref{eq:rec-10} with the positive base weights in \eqref{eq:rec-9}, and optionally
restricting those base weights by a base-measurable event, proves \eqref{eq:rec-1}.

\smallskip\noindent\underline{\textit{Step 3.}}\quad Compare the deterministic and random clock scales.

It remains to make the radius self-normalized.  Whole-path reversal,
translation to the new starting point, and transport of the i.i.d. priorities
preserve the joint law and send
\[
       (\sigma,\ell)\longmapsto(n-\ell,n-\sigma).
\]
Conditionally on $K=n$, the pair has this symmetry. Since $\sigma\le\ell$,
\begin{equation}
E[n-\sigma\mid K=n]=E[\ell\mid K=n]\ge n/2,
       \qquad P(n-\sigma\ge n/4\mid K=n)\ge1/3.                        \label{eq:rec-11}
\end{equation}
Also \(P(K\ge N/2)\ge1/4\) for all large \(N\).  Hence the restriction
\[
       \Gamma=\{n-\sigma\ge N/8\}
\]
has mixture mass \(\beta\ge1/12\).  In the tuple,
\[
       n-\sigma=(\ell-\sigma)+(n-\ell)=\rho_r+|B|,
\]
so \(\Gamma\) is base-measurable, as required above.  Apply Lemma~\ref{lem:stopped-resolvent} to this restricted mixture. It gives a disk of radius $C_1\sqrt N/f$; on $\Gamma$, this is at most $\sqrt8 C_1\sqrt K/f$. Absorb the fixed factor into the radius constant.  This proves \eqref{eq:rec-5}.
\end{proof}

\subsection{Removing ties and absorbing the past}

\begin{lemma}[Localization with a unique favorite]\label{lem:singleton-geometric-localization}
For the same geometric clock, uniformly for \(2\le f\le N^{1/8}\),
\begin{equation}
P\left(\mathcal F_K=\{X\},\
        \lVert X\rVert_2\le C\frac{\sqrt K}{f}\right)
       \ge c\,\frac{1+\log f}{f^2}.                            \label{eq:rec-12}
\end{equation}
The conclusion remains valid, up to a universal change of \(c\), after
restricting the output clock to
\begin{equation}
aN\le K\le JN,\qquad J=\log N,\qquad
       a=\frac c8\,\frac{1+\log f}{f^2}.                       \label{eq:rec-13}
\end{equation}
\end{lemma}

\begin{proof}
For every atom counted in \eqref{eq:rec-5}, insert immediately after the last visit to
its priority-selected favorite \(x\) the fixed loop
\begin{equation}
x,\ x+e_1,\ x,\ x+e_2,\ x.                             \label{eq:rec-14}
\end{equation}
The local time at \(x\) rises by two, those at the two neighbors rise by
one, and all other local times are unchanged.  Thus \(x\) is the unique
favorite of the new path.  The old suffix avoids \(x\), so in the image the
last visit to the unique favorite is the final vertex of \eqref{eq:rec-14}.  Removing
the four preceding prescribed steps is therefore the unique inverse.  On
geometrically weighted path atoms the image-to-source weight ratio is
\(z^4/4^4\), bounded below for \(z\) close to one.  This proves \eqref{eq:rec-12}.

Now \(P(K<aN)\le2a\), while
\(P(K>JN)\le\exp(-J/2)\).  With the constant in \eqref{eq:rec-13} chosen after the
constant in \eqref{eq:rec-12}, these discarded masses total at most half the right-hand
side of \eqref{eq:rec-12}.  Four added steps are absorbed by changing the universal
constants.
\end{proof}

\begin{lemma}[Adding an inherited local-time profile]\label{lem:inherited-profile-dichotomy}
Let \(H\ge0\) be any finitely supported integer profile and let
\(\ell\ge0\) be a second integer profile with a unique maximizer \(y\).
Every maximizer of \(H+\ell\) lies in
\(\operatorname{supp}H\cup\{y\}\).
\end{lemma}

\begin{proof}
If \(y\) is not a maximizer of \(H+\ell\), then
\(\max(H+\ell)>(H+\ell)(y)\ge\ell(y)\).  Outside
\(\operatorname{supp}H\), the profile \(H+\ell\) equals \(\ell\) and is at
most \(\ell(y)\), so every maximizer lies in \(\operatorname{supp}H\).
If \(y\) is a maximizer, uniqueness of the maximum of \(\ell\) shows that
the only possible maximizer outside \(\operatorname{supp}H\) is \(y\).
\end{proof}

\subsection{Infinitely many successful blocks}

\begin{proposition}[Recurrence at the critical scale]\label{prop:critical-recurrence}
Almost surely,
\begin{equation}
\liminf_{n\to\infty}
       \frac{R_n^+}{\sqrt{n/\log n}}=0.                        \label{eq:rec-15}
\end{equation}
Consequently, for every \(\gamma\le\tfrac12\),
\begin{equation}
\liminf_{n\to\infty}\frac{R_n^+}{a_\gamma(n)}=0.
                                                                    \label{eq:rec-16}
\end{equation}
\end{proposition}

\begin{proof}
\smallskip\noindent\underline{\textit{Step 1.}}\quad Choose independent fresh segments.

Fix \(\varepsilon>0\).  Choose a deterministic sequence
\[
       J_{k+1}=J_k+10\log J_k,\qquad s_k=\lfloor e^{J_k}\rfloor,
\]
and, for the fresh block beginning at \(s_k\), put
\[
       N_k=\frac{e^{J_{k+1}}}{J_{k+1}^2},\qquad
       f_k=C\varepsilon^{-1}\sqrt{J_{k+1}}.
\]
Apply Lemma~\ref{lem:singleton-geometric-localization} to the fresh increment
walk, with an independent geometric clock with parameter $1/N_k$.  It gives a
unique fresh favorite \(Y_k\) satisfying
\begin{equation}
\lVert Y_k\rVert\le\frac{\varepsilon}{8}
             \sqrt{\frac{K_k}{J_{k+1}}}                       \label{eq:rec-17}
\end{equation}
with conditional probability at least
\begin{equation}
d_k=c_\varepsilon\frac{\log J_{k+1}}{J_{k+1}},          \label{eq:rec-18}
\end{equation}
even after
\begin{equation}
\frac{d_k}{8}N_k\le K_k\le J_{k+1}N_k.                 \label{eq:rec-19}
\end{equation}
For all sufficiently large $k$, the translated clock interval in \eqref{eq:rec-19} lies in $[s_k,s_{k+1})$: its upper endpoint is $s_k+e^{J_{k+1}}/J_{k+1}<s_{k+1}$, while its lower clock endpoint is much larger than $s_k$. The singleton estimate already includes the four inserted steps.

\smallskip\noindent\underline{\textit{Step 2.}}\quad Include the inherited profile.

Let \(H_k(x)=\sum_{j=0}^{s_k-1}{\bf1}_{\{S_j=x\}}\), while the fresh
profile includes its time-zero visit at \(S_{s_k}\).  Thus their sum is
exactly the full profile at \(s_k+K_k\), with no double counting.  The
fresh profile is translated by \(S_{s_k}\).  By
Lemma~\ref{lem:inherited-profile-dichotomy}, at time
\(t_k=s_k+K_k\) every favorite of the full walk belongs either to the old
range or to \(S_{s_k}+Y_k\).

Part \eqref{eq:rec-2} of Lemma~\ref{lem:elementary-walk-estimates} and Borel--Cantelli give
\[
       \max_{j\le s_k}\lVert S_j\rVert
       \le \sqrt{s_k}\,J_k
\]
eventually.  The choices above imply, uniformly over \eqref{eq:rec-19},
\begin{equation}
\sqrt{s_k}\,J_k=o\left(\sqrt{\frac{K_k}{J_{k+1}}}\right),
 \qquad \log(s_k+K_k)\sim J_{k+1},\qquad s_k=o(K_k).            \label{eq:rec-20}
\end{equation}
Combining \eqref{eq:rec-17}--\eqref{eq:rec-20}, every full favorite at \(t_k\) is in
\begin{equation}
B\left(0,\frac{\varepsilon}{2}
                    \sqrt{\frac{t_k}{\log t_k}}\right)        \label{eq:rec-21}
\end{equation}
for all sufficiently large successful \(k\).

\smallskip\noindent\underline{\textit{Step 3.}}\quad Obtain infinitely many witness times.

Moreover \(J_k\sim10k\log k\), and hence
\begin{equation}
\sum_kd_k=\infty.                                      \label{eq:rec-22}
\end{equation}
The fresh increments occupy disjoint deterministic blocks.  On the product
space enlarged by the independent clocks and priorities, let \(E_k\)
be the fresh singleton-success event \eqref{eq:rec-17}, with clock restriction \eqref{eq:rec-19}.
Its conditional probability given the past at \(s_k\) is at least
\(d_k\).  Lévy's conditional Borel--Cantelli lemma and \eqref{eq:rec-22} imply that
\(E_k\) occurs infinitely often almost surely.  Independently, the
old-range bound and the deterministic comparisons \eqref{eq:rec-20} hold eventually
almost surely; on their intersection, every sufficiently late \(E_k\)
implies \eqref{eq:rec-21}.  Thus \eqref{eq:rec-21} occurs infinitely often.  Fubini removes the
auxiliary clocks and priorities: for almost every walk path, their
probability of producing infinitely many witness times is one, so that
path itself has infinitely many such integer witness times.

Intersecting over rational \(\varepsilon>0\) proves \eqref{eq:rec-15}.  Finally,
\[
       \frac{a_{1/2}(n)}{a_\gamma(n)}
       =(\log n)^{\gamma-1/2}\le1
\]
eventually for \(\gamma\le1/2\), which proves \eqref{eq:rec-16}.
\end{proof}


% ---- continuum_transfer ----
\section{Continuum barriers and lattice traversal counts}
\label{sec:continuum-transfer}

Traversal counts separate the large-scale accumulation of excursions from
occupation inside the last annulus.  We first record the continuum barrier
estimates in the form needed below.  We then compare rare upper events by
multiplying one-step lattice probabilities.  For the lower event, a single
strong coupling transfers the existence of a suitable Brownian center.
The two comparisons have different purposes: the first preserves probabilities
on their exponential scale, while the second preserves a positive probability
of existence.

\subsection{Spherical traversal estimates}

We use Rosen's spherical traversal construction \cite{Rosen2022}.
Its upper and lower arguments use separate fixed choices of the initial
radius.  We keep the customary notation $r_0$ in each argument and state
the change explicitly; the two choices are never used in the same pathwise
comparison.

Let $X$ be spherical Brownian motion started at $v$, stopped at its
first hit $\tau$ of $\partial B_d(v,r_*)$.  When defining source
stopping times we extend it by an independent Brownian continuation;
the safety event below will imply that every portion actually used lies
before $\tau$.  Put
\begin{equation}
h(r)=2\arctan(r/2),\qquad r_j=r_0e^{-j},\qquad h_j=h(r_j),\label{eq:ct-external-23}
\end{equation}
For the upper estimates, choose the fixed radii as in
\cite[Section~2, (2.2)]{Rosen2022}, in particular $2r_*\le h_0$.  Let $F_j$ be a
maximal $d_0h_j$-separated set, augmented if necessary to a
$d_0h_j$-cover; thus $|F_j|\asymp e^{2j}$.  For $y\in F_L$,
$T_{y,j}$ is the number, before $\tau$, of completed excursions from
$\partial B_d(y,h_{j-1})$ to $\partial B_d(y,h_j)$, with the initial
and final incomplete pieces omitted.  Its source shell is
\begin{equation}
k_y=\inf\{k:y\notin B_d(v,h_k)\}.\label{eq:ct-external-24}
\end{equation}
Pairs have separation shell $k$ when
$2h_k<d(y,y')\le2h_{k-1}$.

Set
\[
\rho_L=2-\frac{\log L}{L},\qquad l_L=l\wedge(L-l),
\]
\begin{equation}
\alpha_{z,+}(l)=\rho_Ll+z+l_L^{1/4},\qquad
 \alpha_{z,-}(l)=\rho_Ll+z-l_L^{1/4}.\label{eq:ct-external-25}
\end{equation}
Fix an admissible $z\ge z_0$ in Rosen's upper argument.  The upper
first-crossing event $U(y,l,z)$ means
\[
\sqrt{2T_{y,j}}\le\alpha_{z,+}(j)
       \quad(k_y<j<l),\qquad
 \sqrt{2T_{y,l}}>\alpha_{z,+}(l).
\]
Rosen's representative construction is deterministic.  With fixed
tie-breaking, every fine center is sent to its closest point of the
predetermined coarser net at scale
$h_{l+\frac14 l_L^{1/4}}$; equations (2.26)--(2.33) give the buffered
ball inclusions and domination of its traversal count.  Proposition 2.5,
specifically equations (2.34)--(2.35), then proves the complete estimates
\begin{equation}
\#\{\hbox{representatives with source }k\}
 \le C\exp\{2(l-k)+\tfrac12l_L^{1/4}\},\label{eq:ct-external-26}
\end{equation}
\begin{equation}
P(U(y,l,z))\le
 C\exp\{-2l-l_L^{1/4}\}e^{-2z}.\label{eq:ct-external-27}
\end{equation}
They hold pointwise for $1\le k\le l\le L-(4\log L)^4$.
Only the choice $F_L\cap B_d(v,h_{\log L})^c$, and not the proof of
\eqref{eq:ct-external-26}--\eqref{eq:ct-external-27}, imposed $k\le\log L$.  Hence, for any
$s\le C_0\log L$, summing just the representatives whose source shell
is at least $s$ gives
\begin{equation}
P\big(\exists y:k_y\ge s,\ \exists l\le L-(4\log L)^4:
              U(y,l,z)\big)
       \le C e^{-2s+o(\log L)}e^{-2z}.\label{eq:ct-external-28}
\end{equation}

\subsubsection{The endpoint estimate for complete centered excursions}

The endpoint calculation must be applied to a complete centered excursion
forest.  To specify that law, start spherical Brownian motion at $v$ with
$h_{k_y}\le d(v,y)<h_{k_y-1}$ and continue it until its first exit from
$B_d(y,h_0)$.  Let $T^{\mathrm c}_{y,j}$ count its completed traversals
at level $j$, and let $\mathcal{G}^{\mathrm c}_{y,L,z}$ require the upper
barrier in \eqref{eq:ct-external-25} from $m=k_y+1$ through $L-1$.
Put $Y_y^{\mathrm c}=\sqrt{2T^{\mathrm c}_{y,L}}$ and
$B_L=\rho_LL+z$.  Uniformly for
$k_y\le C_0\log L$, every fixed $\epsilon_0>0$,
$\epsilon_0L\le q\le B_L+1$, and fixed $z$, the source law in
\cite[(2.8)]{Rosen2022} and its Appendix Theorem~9.1(a) give
\begin{equation}
P\big(\mathcal{G}^{\mathrm c}_{y,L,z},Y_y^{\mathrm c}\in[q,q+1]\big)
 \le C(k_y+1)^C(2+B_L-q)L^{-2}e^{-q^2/(2L)}.\label{eq:ct-external-29}
\end{equation}
This assertion concerns the centered stopping rule just defined.  We do not
identify the descendants of a population stopped at the off-center boundary
$\partial B_d(v,r_*)$ with an independent Galton--Watson forest.
The later lattice endpoint estimate will be proved directly for its centered
continuation.

The following version of \cite[Appendix, Theorem~9.1(a)]{Rosen2022}
is sufficient: if $T$ is critical Galton--Watson with
$P(\xi=m)=2^{-m-1}$, $T_0=x^2/2$, and
\[
\begin{gathered}
\delta>0,\quad C\ge0,\quad\eta>1,\quad\epsilon\in(0,1/2),\\
2\le x,y\le\eta L,\quad x^2/2\in\mathbb N,\\
0\le x\le a,\quad0\le y\le b,
\end{gathered}
\]
\begin{equation}
\begin{split}
P\big(&\sqrt{2T_j}\le a+(b-a)j/L
             +C(j\wedge(L-j))^{1/2-\epsilon},\ 1\le j<L;\\
&\sqrt{2T_L}\in[y,y+\delta]\big)
\le C'\frac{(1+a-x)(1+b-y)}L\sqrt{\frac{x}{yL}}
e^{-(x-y)^2/(2L)}.\label{eq:ct-external-30}
\end{split}
\end{equation}
Write $m=k_y+1$ and $N=L-m$.  Conditional on $T^{\mathrm c}_{y,m}=t$,
the $t$ complete source excursions give the critical geometric
Galton--Watson chain in \eqref{eq:ct-external-30}, started from $t$.
Indeed each excursion from level $m$ to level $m-1$ is explored completely
before the centered path can exit at level $0$; its neighboring logarithmic
side choices have probability $1/2$ independently of the entrance angle.
The logarithmic
gambler--ruin calculation (2.8) gives
\[
P(T^{\mathrm c}_{y,m}\ge t)\le Cm e^{-t/m},
\]
and hence, after grouping the integer values of $t$ for which
$x=\sqrt{2t}$ lies in a unit band,
\begin{equation}
P(\sqrt{2T^{\mathrm c}_{y,m}}\in[x,x+1])
       \le Cm^C e^{-x^2/(2m)}.\label{eq:ct-external-30a}
\end{equation}
On $\mathcal{G}^{\mathrm c}_{y,L,z}$, $x\le a:=\alpha_{z,+}(m)=O(m)$.
Apply \eqref{eq:ct-external-30} over the remaining $N$ generations with this $a$ and
with $b=B_L+1$.  The needed deterministic barrier inclusion follows
from
\begin{equation}
((m+j)\wedge(L-m-j))^{1/4}
 \le (1-j/N)m^{1/4}+C(j\wedge(N-j))^{1/4};\label{eq:ct-external-30b}
\end{equation}
for $j\le N/2$ use
$(m+j)^{1/4}\le m^{1/4}+j^{1/4}$ and
$m^{1/4}j/N\le j^{1/4}$, while for $j>N/2$ use the
right-end distance directly.  Thus every parameter in \eqref{eq:ct-external-30} has been
specified, including $x\le a$, $q\le b$, and
$x,q\le\eta N$.

The sole positive start not covered by the hypothesis $x\ge2$ is
$t=1$.  Couple a process from two ancestors as the sum of two
independent one-ancestor processes.  On the event, of probability
$1/2$, that the second ancestor has no child in the first generation,
the total process agrees from generation $1$ onward with the first
process.  The barrier in \eqref{eq:ct-external-30} is imposed precisely at generations
$1,\ldots,N-1$, so the $t=1$ contribution is at most twice the same
event for initial population two, to which \eqref{eq:ct-external-30} applies.  The case
$t=0$ cannot have the positive endpoint $q\ge\epsilon_0L$.

Multiplication of \eqref{eq:ct-external-30a} with \eqref{eq:ct-external-30}, the identity
\begin{equation}
\frac{x^2}{k+1}+\frac{(q-x)^2}{L-k-1}
 =\frac{q^2}{L}+
   \frac{L(x-(k+1)q/L)^2}{(k+1)(L-k-1)},\label{eq:ct-external-31}
\end{equation}
proves \eqref{eq:ct-external-29}.  Indeed $q\ge\epsilon_0L$, $N\asymp L$,
$x\le Cm$, and so the square-root factor in \eqref{eq:ct-external-30}, together with its
explicit $N^{-1}$, is at most $Cm^{1/2}L^{-2}$; all remaining source
and summation factors cost only $m^C$.

No estimate of the form \eqref{eq:ct-external-29} is asserted at $q=0$: extinction there
has order-one probability.  In the lattice upper bound below,
$q<\epsilon_0L$ is instead disposed of by the saturated terminal
Chernoff factor.  We shall also use the following truncated version.
If $n=J-h$, $m\le C_0\log J$, and
\[
B_l=\rho_Jl+z+(l\wedge(J-l))^{1/4},
\]
then \eqref{eq:ct-external-29}, with $L$ replaced by $n$ and $B_L$ by $B_n$, holds
for $\epsilon_0n\le q\le B_n+1$.  To verify the only new point, put
$\phi(l)=(l\wedge(J-l))^{1/4}$.  Its two Hölder bounds from $m$
and $n$ give
\begin{equation}
\phi(m+j)\le(1-j/N)\phi(m)+(j/N)\phi(n)
                 +C(j\wedge(N-j))^{1/4}.\label{eq:ct-external-31a}
\end{equation}
The affine part $\rho_Jl+z$ interpolates exactly, so \eqref{eq:ct-external-30} applies with
the arbitrary endpoints $a=B_m,b=B_n+1$.  In particular, no
$\rho_n$-barrier theorem is being assumed.

\subsubsection{A fixed number of source excursions}

For the lower event, now choose $r_0$ sufficiently small that
$4h(r_{-1})<r_*$, as in \cite[Section~3]{Rosen2022}.  The preceding
upper estimates retain their separate choice of radii.  Fix the source
parameter $x_*>0$ from that section so that
\begin{equation}
m_0:=x_*^2\in\mathbb N.\label{eq:ct-external-32a}
\end{equation}
(This $x_*$ is not the variable called $x$ in Appendix Theorem 9.1;
the latter equals $\sqrt{2m_0}$ when that theorem is started with
$m_0$ particles.)  Put
\[
F_L^0=F_L\cap\{h_1/40\le d(v,y)\le h_1/20\}.
\]
This fixed annulus, instead of Rosen's full ball, retains a fixed
positive area fraction; the one- and two-point sums in Rosen's proof are
therefore unchanged up to constants, and every source index is $O(1)$.

We spell out the superscript in Rosen (3.3), because it is important for
the stopping-time comparison below.  Write
$\Gamma_{y,j}=\partial B_d(y,h_j)$ and set
\[
A_{y,1}=H_{\Gamma_{y,1}},\qquad
 B_{y,i}=\inf\{t>A_{y,i}:X_t\in\Gamma_{y,0}\},
\]
\begin{equation}
A_{y,i+1}=\inf\{t>B_{y,i}:X_t\in\Gamma_{y,1}\},
                    \qquad 1\le i<m_0.\label{eq:ct-external-32b}
\end{equation}
The intervals $[A_{y,i},B_{y,i}]$, and only these intervals, are the
first $m_0$ source excursions; the intervening intervals
$[B_{y,i},A_{y,i+1}]$ are connectors and are not counted.  Set
$T_{y,1}^{1,m_0}=m_0$.  For $l\ge2$, define inside each source
interval
\[
C_{i,l,1}=\inf\{t\ge A_{y,i}:X_t\in\Gamma_{y,l-1}\},
\]
\[
D_{i,l,q}=\inf\{t>C_{i,l,q}:X_t\in\Gamma_{y,l}\},\qquad
 C_{i,l,q+1}=\inf\{t>D_{i,l,q}:X_t\in\Gamma_{y,l-1}\},
\]
and put
\begin{equation}
T_{y,l}^{1,m_0}=\sum_{i=1}^{m_0}
             \#\{q:D_{i,l,q}<B_{y,i}\}.\label{eq:ct-external-32c}
\end{equation}
Thus Rosen's abbreviation is exactly
$T_{y,l}^{\,1}=T_{y,l}^{1,x_*^2}$, not a count in one source
excursion.  The strong Markov property and logarithmic gambler's ruin
make $(T_{y,1+j}^{1,m_0})_{j\ge0}$ a critical Galton--Watson chain
with $m_0$ ancestors and offspring law
$P(\xi=r)=2^{-r-1}$, $r\ge0$.

Define $\widehat I_{y,z}$ by
\begin{equation}
\sqrt{2T_{y,l}^{1,m_0}}\le\alpha_{z,-}(l),\quad1\le l<L,
 \qquad \sqrt{2T_{y,L}^{1,m_0}}\ge\rho_LL+z.\label{eq:ct-external-32}
\end{equation}
For the outward excursions at scale $k$, let $\theta_{k,i}$ be their
angular increments and let $\nu_k$ be the exit-angle law from radius
$r_k$ to $r_{k-1}$.  If
\[
N_{k,a}=\left\lfloor\frac{(\rho_Lk+z-a+1)^2}{2}\right\rfloor,
\]
 set $N_k=N_{k,a}$ on
$\sqrt{2T_{y,k}^{1,m_0}}\in[\rho_Lk+z-a,\rho_Lk+z-a+1]$, and
\begin{equation}
W_{y,k}(N_k)=\left\{
 d_{W^1}\left(N_k^{-1}\sum_{i\le N_k}\delta_{\theta_{k,i}},
                 \nu_k\right)
 \le\frac{c_0\log(L-k)}{2\sqrt{N_k}}\right\}.\label{eq:ct-external-33}
\end{equation}
With $L^+=L-(500\log L)^4$ and Rosen's fixed $d_*$, the event is
\begin{equation}
I_{y,z}=\widehat I_{y,z}
          \cap\bigcap_{k=L^+}^{L-d_*}W_{y,k}(N_k).\label{eq:ct-external-34}
\end{equation}
Rosen Lemmas 3.3--3.4 state completely that, uniformly for
$0\le z\le\log L$,
\begin{equation}
P(I_{y,z})\asymp(1+z)e^{-2L}e^{-2z},\label{eq:ct-external-35}
\end{equation}
and, for a pair in separation class $k$,
\begin{equation}
P(I_{y,z}\cap I_{y',z})
 \le C(1+z)e^{-4L+2k}e^{-2z}
       e^{-c(k\wedge(L-k))^{1/4}}.\label{eq:ct-external-36}
\end{equation}
The convention for $L-d_*\le k<L$ is the one-center bound from that
lemma.  Fix $z_0$ in Rosen's admissible range.  After imposing the
source-safety event below we will choose a fixed $C_*$ and use the
truncated events
\begin{equation}
I_{y}^{\circ}:=I_{y,z_0}\cap
 \{\rho_LL+z_0\le\sqrt{2T_{y,L}^{1,m_0}}
                         \le\rho_LL+z_0+C_*\}\label{eq:ct-external-37}
\end{equation}
The notation in \eqref{eq:ct-external-37} is provisional until $C_*$ is chosen after
safety has been proved.  This order matters: because $X$ was extended
after $\tau$, the inequality
$T_{y,L}^{1,m_0}\le T_{y,L}^{\tau}$ can fail before safety is imposed.

\subsubsection{Keeping the source excursions before the stopping time}

The lower event was defined using a Brownian continuation after $\tau$.
We must therefore force the entire source forest to finish before $\tau$
without losing its probability on the scale $e^{-2L}$.  Let $S_y$ require every connector
$[B_{y,i},A_{y,i+1}]$, $i<m_0$, to hit $\Gamma_{y,1}$ before the
larger circle $\Gamma_{y,-1}$.  In the centered logarithmic cylinder,
the unrestricted connector endpoint subdensity and the safely killed
one, relative to Haar measure, are respectively
\[
q_\infty(u)={\sinh1\over\cosh1-\cos u},
\]
\begin{equation}
q_{\rm sf}(u)=\frac12+2\sum_{n\ge1}{\sinh n\over\sinh(2n)}\cos(nu)
 =\frac\pi2\sum_{k\in\mathbb Z}
       \operatorname {sech}\!\left({\pi\over2}(u+2\pi k)\right).\label{eq:ct-external-38a}
\end{equation}
Both are continuous and strictly positive, so
$\kappa_*:=\min_uq_{\rm sf}(u)/q_\infty(u)>0$.  Thus the safe
connector kernel dominates $\kappa_*$ times the unrestricted kernel
pointwise.  Iterating the strong Markov factorization, with the whole
later path functional left in the last factor, gives
\begin{equation}
P(I_{y,z}\cap S_y)\ge\kappa_*^{m_0-1}P(I_{y,z}).\label{eq:ct-external-38b}
\end{equation}
This is finite-measure domination against the entire later functional,
not multiplication of marginal probabilities.  Before endpoint
truncation put
\[
\widetilde Z_L^{\rm sf}:=\sum_{y\in F_L^0}{\bf1}_{I_{y,z_0}\cap S_y}.
\]
Equation \eqref{eq:ct-external-38b} preserves its first moment, whereas
$I_{y,z_0}\cap S_y\subset I_{y,z_0}$ preserves the pair upper bound
\eqref{eq:ct-external-36}.  Rosen's shell summation and Paley--Zygmund therefore give
\begin{equation}
E\widetilde Z_L^{\rm sf}\le C,\qquad
 P(\widetilde Z_L^{\rm sf}>0)\ge c_1>0.\label{eq:ct-external-38c}
\end{equation}
Rosen chooses $r_0$ so that $4h(r_{-1})<r_*$.  Since
$d(v,y)\le h_1/20$, there is a fixed $\chi>0$ such that
\begin{equation}
B_d(y,h(r_{-1}))\subset B_d(v,r_*-\chi).\label{eq:ct-external-38d}
\end{equation}
Every source block and every connector on $S_y$ lies in the left
side, hence ends before $\tau$.  This is the exact stopped-source
certificate used in the coupling below.  In particular, on $S_y$,
\begin{equation}
T_{y,L}^{1,m_0}\le T_{y,L}^{\tau}.\label{eq:ct-external-38e}
\end{equation}
Rosen's Theorem~1.6 gives the following upper tail on $F_L$; extending
its bound to bounded nonnegative $z$ only enlarges the constant:
\begin{equation}
P\left(\sup_{y\in F_L}\sqrt{2T_{y,L}^{\tau}}\ge\rho_LL+z\right)
       \le C(1+z)e^{-2z},\qquad0\le z\le\log L.\label{eq:ct-external-38f}
\end{equation}
Choose $C_*$ so large that the right side at $z=z_0+C_*$
is below $c_1/2$, and now define $I_y^\circ$ by \eqref{eq:ct-external-37}.  Equations
\eqref{eq:ct-external-38c}--\eqref{eq:ct-external-38f} give, with
\[
Z_L^{\rm sf}:=\sum_{y\in F_L^0}{\bf1}_{I_y^\circ\cap S_y},
\]
\begin{equation}
EZ_L^{\rm sf}\le C,\qquad P(Z_L^{\rm sf}>0)\ge c_1/2.\label{eq:ct-external-38}
\end{equation}
The constant $C_*$ and the safe count are now fixed.

\begin{lemma}[Continuum barrier inputs]\label{lem:external-barrier-input}
In the upper choice of radii, the first-crossing estimate
\eqref{eq:ct-external-28} holds.  For the complete centered traversal process,
the endpoint estimate \eqref{eq:ct-external-29} holds, including the truncated
barrier described in \eqref{eq:ct-external-31a}.  In the lower choice of radii,
on $I_y^\circ\cap S_y$ all first $m_0$ source excursions finish before
$\tau$, and the number of centers with this event satisfies
\eqref{eq:ct-external-38}.
\end{lemma}

\begin{proof}[Proof of Lemma~\ref{lem:external-barrier-input}]
The representative estimates are the packing and one-center bounds of
\cite[Proposition~2.5, (2.34)--(2.35)]{Rosen2022}, summed over the indicated
source shells.  The complete centered source calculation above proves
\eqref{eq:ct-external-29}, and \eqref{eq:ct-external-31a} checks the change of
barrier at a truncated horizon.  The lower construction combines
\cite[Lemmas~3.3--3.4]{Rosen2022} with the connector-kernel domination
\eqref{eq:ct-external-38b} and the terminal truncation.
All the first- and second-moment assertions \eqref{eq:ct-external-35}--\eqref{eq:ct-external-38} remain in Rosen's
original spherical geometry.  In particular, we do \emph{not} infer an
event-level Euclidean two-center theorem from Rosen's occupation-measure
argument in Section 8.  The next lemma instead maps the actual spherical
interfaces into the plane, so the spherical event and its pair law are
preserved exactly until the one-sided lattice coupling is made.
\end{proof}

\subsection{Comparison of smooth interface kernels}

Fix $\alpha\in(0,1)$.  Let
$(\Gamma_R^-,\Gamma_R^0,\Gamma_R^+)$ range over triples of
disjoint Jordan curves which, after translation, rotation, and scaling
by $R^{-1}$, together with their tubular-coordinate charts, belong to
a bounded $C^{4,\alpha}$ family.  Assume consecutive curves are
separated by distances comparable to $R$, and that a specified
conformal coordinate sends each triple to three concentric logarithmic
levels.  The interfaces used below are circles or uniformly
Möbius-transformed circles, so this hypothesis entails no loss.

For a boundary function put
\begin{equation}
\|f\|_{4,R}:=\sum_{j=0}^4R^j\|\partial_\tau^jf\|_\infty
             +R^{4+\alpha}[\partial_\tau^4f]_\alpha .\label{eq:ct-smooth-39a}
\end{equation}
A lattice starting state is a directed edge whose endpoints lie on opposite
sides of $\Gamma_R^0$; the lattice query starts from its post-crossing
endpoint and the Brownian query from the normal projection of that endpoint
onto $\Gamma_R^0$.  The side of a vertex is determined by membership in
the bounded component of the complement (vertices lying on the curve are
assigned to the exterior).  Let $Q_R^B$ be the Brownian selected-side
subprobability kernel, stopped at first contact with either competing curve,
and let $Q_R^S$ be the lattice selected-side subprobability kernel, stopped
at the first directed edge whose endpoint-side indicators differ for one of
the two competing curves.  A tangency, or any intersection of an interpolated
edge whose two endpoints remain on the same side, is not a lattice crossing.
Write $\pi_0,\pi_+$ for normal projection of the digital source
and selected target states.  For a function $f$ on the selected target
curve and a function $g$ on the middle curve, define the two differently
typed lifts by
\[
U_+f=f\circ\pi_+,
                 \qquad U_0g=g\circ\pi_0.
\]
\begin{lemma}[Comparison of interface kernels]
\label{lem:smooth-interface-discretization}
As an estimate from functions on the selected target curve to functions
on the digital source-edge space,
\begin{equation}
\|Q_R^SU_+f-U_0(Q_R^Bf)\|_\infty
       \le CR^{-1}\|f\|_{4,R}.\label{eq:ct-smooth-39b}
\end{equation}
If the Brownian side probability $p_R=Q_R^B1$ is constant on the
middle curve and lies in $[c,1-c]$, then, uniformly in the starting
angle,
\begin{equation}
Q_R^S1=p_R(1+O(R^{-1}))\label{eq:ct-smooth-39c}
\end{equation}
and the corresponding row-normalized angular kernels differ by
$O(R^{-1})$ on $C^4$ tests.

Consequently, let a chronological radial word consist of $N$
ordinary decisions between neighboring concentric interfaces at
scales $R_i\ge M$.  The next queried triple may depend on the already
exposed radial prefix, but not on an unexposed angular mark.  For a fixed
word, let $E_i$ be the directed-edge state space after its first $i$
letters and let $K^S_{i,w_{i+1}}(e,de')$ be the subkernel realizing the
next prescribed side.  If $p_i$ is the corresponding Brownian side
probability, then, for every current edge,
\begin{equation}
e^{-C/R_i}p_i\le K^S_{i,w_{i+1}}1(e)
          \le e^{C/R_i}p_i .\label{eq:ct-smooth-39cprime}
\end{equation}
After summing all compatible directed-edge refinements,
\begin{equation}
e^{-C\sum_iR_i^{-1}}P^B(w)
 \le P^S(w)\le
 e^{C\sum_iR_i^{-1}}P^B(w).\label{eq:ct-smooth-39d}
\end{equation}
The same inequality holds after integrating the initial edge against any
entrance law and after summing a disjoint family of complete words.  A
radial-prefix stopping rule is padded by cemetery letters of row mass one.
No assertion is made after conditioning on an edge list generated later by
the word.  Source-root and source-death letters are handled explicitly in
Lemma \ref{lem:buffered-transfer}.

\end{lemma}

\begin{proof}[Proof of Lemma~\ref{lem:smooth-interface-discretization}]
Let $\Omega_R$ be the annulus between the two competing curves and let
$u$ be the continuum Dirichlet solution with boundary data $f$ on the
selected component and zero on the other.  Uniform rescaled boundary
regularity gives
\begin{equation}
\|D^ju\|_{L^\infty(\overline\Omega_R)}
       \le CR^{-j}\|f\|_{4,R},\qquad0\le j\le4.\label{eq:ct-smooth-39e}
\end{equation}
The compact rescaled family has a tubular radius $\rho R$.  In a signed
normal chart $\Phi(s,t)$, with $t\ge0$ on $\Omega_R$, choose fixed
numbers $c_1,\ldots,c_5$ satisfying
\[
\sum_{a=1}^5c_a(-a)^q=1,
             \qquad q=0,1,2,3,4,
\]
and, for $-\rho R/5<t<0$, set
\[
\widetilde u(\Phi(s,t))=\sum_{a=1}^5c_a u(\Phi(s,-at)).
\]
Do this in the two disjoint collars and retain $\widetilde u=u$ on
$\overline\Omega_R$.  The five matching identities make
$\widetilde u$ a $C^4$ continuation through the boundary, with the
same scaled derivative bounds as \eqref{eq:ct-smooth-39e}.

Let $\Omega_R$ be understood with the same vertex-side convention and set
\[
T=\inf\{n\ge1:S_{n-1}\in\Omega_R,\ S_n\notin\Omega_R\}.
\]
Equivalently, $T$ is the first side-changing directed edge across one of
the two competing components; same-side contacts of the interpolated edge
are ignored.  For $n<T$, $S_n$ lies in the open annulus or on its inner
boundary (with the exterior convention above); it cannot lie on the outer
boundary.  At an open-annulus point
$\Delta\widetilde u(S_n)=\Delta u(S_n)=0$, while at an inner-boundary
lattice vertex the $C^4$ matching and continuity from the annulus give
\[
\Delta\widetilde u(S_n)
       =\lim_{\Omega_R\ni x\to S_n}\Delta u(x)=0.
\]
Thus Taylor expansion is legitimate on the full four-neighbor stencil and
its quadratic term vanishes in either case.  The linear and cubic terms
cancel by symmetry.
Thus
\begin{equation}
|(P-I)\widetilde u(S_n)|\le CR^{-4}\|f\|_{4,R},
             \qquad n<T.\label{eq:ct-smooth-39f}
\end{equation}
For large $R$, the two boundary components are more than two lattice
spacings apart and each has tubular radius $\rho R$.  Hence the exit edge
crosses exactly one component, its post-crossing endpoint is within one
lattice spacing of that component and has a unique normal projection.
Moreover $T$ is no larger than the exit time of a fixed $CR$-disk
containing the annulus.  Optional stopping of
$\lVert S_n-z_*\rVert^2-n$ gives $ET\le CR^2$, so the accumulated
defect in \eqref{eq:ct-smooth-39f} is $O(R^{-2})\|f\|_{4,R}$.  The terminal endpoint is
within one lattice unit of the uniquely crossed component; its value
under $\widetilde u$ differs from the projected boundary payoff by
$O(R^{-1})\|f\|_{4,R}$.  The post-crossing starting endpoint and its
projection onto $\Gamma_R^0$ have the same error.  Optional stopping of
$\widetilde u(S_{n\wedge T})$ minus its defect sum proves \eqref{eq:ct-smooth-39b}.
Taking $f=1$, using $p_R\in[c,1-c]$, and dividing by the row mass
proves \eqref{eq:ct-smooth-39c}.  Notice that the continuation is used only to make a
boundary-crossing stencil legal; harmonicity, or its boundary-continuity
consequence, is invoked only at a live vertex.

For completeness, the logarithmic annular probability behind \eqref{eq:ct-smooth-39c} is
also the following cited lattice estimate.  Hao--Li--Okada--Zheng,
\emph{Favorite sites for simple random walk in two and more dimensions},
\cite{HLOZ2025}, Lemma 2.1,
equation (2.5), states in full that, uniformly for
$r<\lVert x\rVert<R$,
\[
P_x(H_{\partial D(0,R)}<H_{\partial D(0,r)})
 ={\log(\lVert x\rVert/r)+O(r^{-1})\over\log(R/r)},
\]
and the complementary inner-hit probability is
\[
P_x(H_{\partial D(0,r)}<H_{\partial D(0,R)})
 ={\log(R/\lVert x\rVert)+O(r^{-1})\over\log(R/r)}.
\]
At fixed logarithmic gaps the main terms are bounded away from zero.

Finally fix a prescribed radial word and let $\nu_i^S$ be the finite
measure on its current directed-edge states after the first $i$ letters,
with every earlier edge integrated out.  The radial prefix fixes the next
interface and prescribed side before the current edge is integrated.
Brownian conformal radial symmetry makes its row mass $p_i$ independent
of the angular endpoint, while \eqref{eq:ct-smooth-39cprime} gives
\begin{equation}
e^{-C/R_i}p_i\|\nu_i^S\|
 \le\|\nu_i^SK^S_{i,w_{i+1}}\|
 \le e^{C/R_i}p_i\|\nu_i^S\|.\label{eq:ct-smooth-39iprime}
\end{equation}
Forward induction proves \eqref{eq:ct-smooth-39d}.  This remains valid for an adaptive
exploration because the radial prefix determines its next query, and for a
stopped exploration after adjoining the absorbing cemetery state.  It is a
forward marginal comparison, not a conditional law given future entrance
points, and it never applies \eqref{eq:ct-smooth-39b} to a rare future indicator.
\end{proof}

\subsection{Transfer of traversal events to the lattice}

Fix $A,C_0<\infty$ and a positive barrier intercept $z$ in the
admissible range above.  The upper first-crossing estimate below is uniform
for $0\le s\le C_0\log J$.  Choose $h=h_J$ so that
\begin{equation}
{h\over\log J}\longrightarrow\infty,\qquad h=o(J),\qquad
 M=e^h,\qquad n=J-h+O(1).\label{eq:ct-transfer-39}
\end{equation}
For the lower comparison let $v$ and $r_*$ be Rosen's starting
point and outer stopping radius.  Fix stereographic coordinates $\sigma$ with
$\sigma(v)=0$, put $R_*=2\tan(r_*/2)$, and dilate by
$A_J=e^J/R_*$:
\[
\Psi_J=A_J\sigma .
\]
Thus the spherical stopped path becomes, up to an increasing time
change, planar Brownian motion started at zero and stopped on
$\partial B(0,e^J)$.  For a spherical center $y$, let $\chi_y$ be
the rotated stereographic coordinate satisfying
\[
\chi_y(y)=0,\qquad |\chi_y(x)|=2\tan(d(x,y)/2).
\]
The lower-comparison physical interfaces are
\begin{equation}
\Gamma_{y,j}=\Psi_J\{x:|\chi_y(x)|=r_0e^{-j}\},
       \qquad-1\le j\le n.\label{eq:ct-transfer-39g}
\end{equation}
They are Euclidean circles, although their Euclidean centers need not be
the same.  In the conformal coordinate
$\chi_y\circ\Psi_J^{-1}$, they are exactly the concentric logarithmic
levels used by Rosen.  Their physical scale is comparable to
$e^{J-j}$; the $O(1)$ choice in $n=J-h+O(1)$ makes the deepest
scale comparable to $M$.

The linearly interpolated lattice path identifies which smooth curve a
lattice edge crosses, but every probabilistic cut is made at an integer
edge-index stopping time.  If $D_{y,j}$ is the bounded component inside
$\Gamma_{y,j}$, then, for adjacent levels and after an outer hit
$\eta_i$, set
\begin{equation}
\begin{aligned}
\sigma_i&=\inf\{n>\eta_i:S_{n-1}\notin D_{y,j},\ S_n\in D_{y,j}\},\\
\eta_{i+1}&=\inf\{n>\sigma_i:S_{n-1}\in D_{y,j-1},\
                                  S_n\notin D_{y,j-1}\}.
\end{aligned}\label{eq:ct-transfer-39h}
\end{equation}
A state records the directed crossing edge; interpolation is used only to
assign its projected $\chi_y$-angle.
Only portions belonging to the first $m_0$ source intervals
$[A_{y,i},B_{y,i}]$ in \eqref{eq:ct-external-32b} contribute traversal counts or terminal
bridges.  Inter-source connectors contribute only their guarded coarse
boundary decisions, never their interior occupation.  The whole segment
from $\sigma_i$ to $\eta_{i+1}$, including further contacts with the
inner interface, is one outward piece.  Recursing inside each parent
excursion gives the standard nested excursion tree.  The lower
\emph{radial word} $W_y=(W_{y,1},\ldots,W_{y,m_0})$ is a fixed collection
of $m_0$ centered words, one for each source interval, with all angles
erased; the complete guard record additionally contains the finitely many
connector boundary decisions.  If
$\ell_{y,i}\in\{1,-1\}$ records whether the connector after
$B_{y,i}$ first hits $\Gamma_{y,1}$ or $\Gamma_{y,-1}$, write
\begin{equation}
\overline W_y=(W_{y,1},\ell_{y,1},W_{y,2},\ldots,
                 \ell_{y,m_0-1},W_{y,m_0})\label{eq:ct-transfer-39hprime}
\end{equation}
for the complete conditioned guard word.  Its traversal coordinate
$T^S_{y,j}$ is the sum over these
words of completed crossings from $\Gamma_{y,j-1}$ to
$\Gamma_{y,j}$.  At the last
fixed number of levels, the interfaces are replaced through one
unit-width digital collar by concentric circles about the nearest lattice
point to $\Psi_J(y)$, of radii $M_y,eM_y$ with
$M_y\in[cM,CM]$.  We relabel $M_y$ as $M$; every later estimate is
uniform under this fixed-factor change and $\log M_y=h+O(1)$.
Define $\mathcal{R}_y$ by the chronological erasure map which replaces each
completed terminal segment, from its inward directed crossing of radius
$M$ through its next outward directed crossing of radius $eM$, by the
two endpoint edges.  It retains $\overline W_y$, all radial letters and
all crossing edges and projected angles outside the erased segments, but
it does not retain an erased segment's duration, interior crossings, path,
or occupation.  Inward connector pieces outside these erased segments are
exposed.  This quotient sigma-field is the intrinsic coarse sigma-field.

The upper application is instead defined directly on the lattice and is
entirely centered.  Write $N=e^J$, $\tau_N=\tau_{D(0,N)}$, and, for
each candidate cell center with $\|y\|<2N$, set
\begin{equation}
\sigma_y=\tau_{D(y,3N)}.\label{eq:ct-transfer-39k}
\end{equation}
The pathwise inclusions
\begin{equation}
D(0,N)\subset D(y,3N),\qquad \tau_N\le\sigma_y,
 \qquad L_{\tau_N}(x)\le L_{\sigma_y}(x)\label{eq:ct-transfer-39l}
\end{equation}
hold for every $x$.  For each one-cell probability we decompose the
extended path up to $\sigma_y$ using only circles centered at $y$,
and sum those probabilities only afterwards.  Replacing $N$ by $3N$
changes every logarithmic index by $O(1)$.  Use Rosen's formal horizon
$J$, its barrier $\rho_Jj+z+(j\wedge(J-j))^{1/4}$, and retain only
interfaces through $n$.  Let $\mathcal{U}_{J,s}$ be the union of these
centered lattice upper first-crossing events, with source shell at least
$s$, and let
$\mathcal{G}^{S,J}_{y,n,z}$ be the corresponding no-crossing event through
$n-1$, and put
\[
B_n^J=\rho_Jn+z+(n\wedge(J-n))^{1/4}.
\]
On the detailed range $\|y\|>M$, its exact source index is
\begin{equation}
u_y=\log {3e^J\over\|y\|},\qquad
 k_y=\lceil u_y\rceil,\qquad m_y=k_y+1;\label{eq:ct-transfer-39m}
\end{equation}
centers with $\|y\|\le M$ belong to the separately treated near-start
range.  Thus $u_y\in(k_y-1,k_y]$ and $m_y-u_y\in[1,2)$.
For the lower application relabel the numerical depth $n$ as $L$,
use the genuine horizon-$L$ quantities $\rho_L,l_L$, and let
$I_y^{S,\circ}$ be the \emph{radial} lattice version of \eqref{eq:ct-external-37}: it retains
the traversal barrier and the bounded terminal band but drops every
angular screen.  Angular screens are used only in the spherical
Brownian certificate that supplies the second moment.

\begin{lemma}[Transfer of traversal events]\label{lem:buffered-transfer}
With the definitions above,
\begin{equation}
P(\mathcal{U}_{J,s})\le J^{-2s/\log J+o(1)}+O(J^{-A}),\label{eq:ct-transfer-40}
\end{equation}
\begin{equation}
P\big(\mathcal{G}^{S,J}_{y,n,z},
       \sqrt{2T^S_{y,n}}\in[a,a+1]\big)
 \le C(k_y+1)^C(2+B_n^J-a)n^{-2}e^{-a^2/(2n)},\label{eq:ct-transfer-41}
\end{equation}
uniformly for $k_y\le C_0\log J$ and every fixed central range
$\epsilon_0n\le a\le B_n^J+1$.  No such assertion is made in the
extinction range $a<\epsilon_0n$.  Most importantly, if
\[
Z_L^S=\sum_{y\in F_L^0}{\bf1}_{I_y^{S,\circ}},
\]
then
\begin{equation}
EZ_L^S\le C,\qquad P(Z_L^S>0)\ge c.\label{eq:ct-transfer-42}
\end{equation}
Only the Brownian proof of \eqref{eq:ct-transfer-42} uses the two-center estimate \eqref{eq:ct-external-36}.
No two-center lattice estimate is claimed or used.

\end{lemma}

\begin{proof}[Proof of Lemma~\ref{lem:buffered-transfer}]
There are two logically separate comparisons.  Rare upper and endpoint
events are compared directly, word by word, using Lemma
\ref{lem:smooth-interface-discretization}.  The global strong coupling is used
only in the safe direction ``strict Brownian lower root implies loose
lattice lower root.''  At no point is a Brownian bad-subkernel estimate
multiplied by an arbitrary future functional of the coupled lattice path.
In the geometric estimates below $L$ denotes the numerical depth $n$;
the superscript $J$ records when the upper barrier retains the larger
formal horizon $J$.

\smallskip\noindent
\underline{\textit{Step 1.}}\quad Geometry and one global coupling.

The conformal image retains the original spherical excursion structure.  Stereographic
projection is conformal and sends spherical Brownian motion to planar
Brownian motion under a continuous increasing time change.  Moreover,
by the definition of $\chi_y$,
\[
\Psi_J\bigl(\partial B_d(y,h(r_0e^{-j}))\bigr)
       =\Gamma_{y,j}.
\]
Thus hitting orders, the complete nested excursion tree, all traversal
counts, and all $\chi_y$-angular marks are literally Rosen's spherical
objects.  In particular, \eqref{eq:ct-external-36} will be applied only before any lattice
comparison.

For the strong approximation we use Zaitsev, Estimates for the Strong
Approximation in Multidimensional Central Limit Theorem,
\cite{Zaitsev2003},
Theorem 2 and Corollary 1.  The form needed here is: for fixed
dimension $d$, centered independent vectors with identity covariance
whose laws lie in the analytic finite-exponential-moment class
$\mathcal{A}_d(\tau)$ for a fixed $\tau\ge1$, one can couple them with independent standard
Gaussian vectors so that, for
$\Delta_N=\max_{k\le N}\lVert\sum_{i\le k}(X_i-Y_i)\rVert$,
\[
E\exp\{c\Delta_N/(\tau d^{7/2}\log^*d)\}
 \le \exp\{C d^{9/4+\alpha}\log^*(N/\tau^2)\},
\]
and hence
\begin{equation}
P(\Delta_N>C\log N+x)\le C e^{-cx},\qquad x\ge0,\label{eq:ct-transfer-30}
\end{equation}
for a fixed i.i.d. law.  The planar walk increment multiplied by $\sqrt2$
has identity covariance and bounded support, so the hypotheses apply.

Take the deterministic time cap
\[
T_J=e^{2J}J^{20}.
\]
Lemma \ref{lem:elementary-walk-estimates} and its Brownian analogue show that
the walk or Brownian motion remains inside the outer disk beyond $T_J$
with probability $e^{-cJ^{10}}$.  By \eqref{eq:ct-transfer-30}, after changing Brownian time
by the covariance factor $1/2$, there is a coupling for which
\begin{equation}
\sup_{t\le T_J}\lVert S^{\rm lin}_{2t}-B_t\rVert
       \le b_J:=CJ^8\label{eq:ct-transfer-31}
\end{equation}
except on an event of probability $e^{-cJ^8}$.  Brownian interpolation
between integer times has a smaller Gaussian error.

The crossing convention is an intrinsic deterministic function of the
finite lattice path.  For a curve with bounded component $D$, a lattice
crossing is recorded exactly when
$\mathbf1_{\{S_{n-1}\in D\}}\ne\mathbf1_{\{S_n\in D\}}$; a curve vertex
is assigned to the exterior, and same-side tangencies or chord contacts are
ignored.  The crossing angle is the normal projection of the post-crossing
endpoint.  At all used scales the endpoint lies in the uniform tubular
collar, so this projection is unique.  On the compact family of
centers under consideration, the maps $\chi_y\circ\Psi_J^{-1}$ and
their inverses have uniformly bounded scaled first two derivatives.
Consequently, if two physical paths are within $b$, then every
crossing of a nominal interface by one path is sandwiched between the
crossings by the other path of the two interfaces obtained by changing
its $\log|\chi_y|$-level by $Cb/R$, where $R$ is the physical
scale.  This is just the deterministic inclusion of the two
$b$-neighborhoods of a smooth separating curve; applying it
successively preserves chronological order.
The crossed directed edge retains the digital-boundary overshoot, which
costs at most two extra lattice units.  Thus
the only difference between the nominal Brownian word and the lattice
word comes from collars of relative logarithmic width
\begin{equation}
\delta_J=C\frac{b_J+1}{M}=O\left(\frac{J^8}{M}\right).\label{eq:ct-transfer-33}
\end{equation}

\smallskip\noindent
\underline{\textit{Step 2.}}\quad Preserving the lower radial word.

Uniform path closeness alone does not preserve an excursion tree: a path
may enter a narrow collar, retreat, and thereby insert or delete a traversal.
We exclude precisely these retreats and wrong-side approaches.  Their total
conditional probability will be small for every prescribed radial word.  Work in the centered cylinder coordinate
\[
U_t=-\log\bigl(|\chi_y(X_t)|/r_0\bigr),
\]
after its conformal time change, so that $U$ is one-dimensional Brownian
motion and $\Gamma_{y,j}=\{U=j\}$.  For a stopping time $s$, write
$H_a[s]=\inf\{t\ge s:U_t=a\}$, with $\inf\varnothing=\infty$.  Put
\begin{equation}
\eta_J=8\delta_J,\qquad N_0=J^4,\label{eq:ct-transfer-R1}
\end{equation}
and take $J$ large enough that $0<\delta_J<\eta_J/4$ and
$\eta_J\le1/4$.

Here is the finite depth-$L$ word belonging to \eqref{eq:ct-external-32b}.  From
$S_{i,0}=A_{y,i}$, $K_{i,0}=1$, recursively set, when $K_{i,r}<L$,
\[
S_{i,r+1}=\inf\{t>S_{i,r}:U_t\in
             \{K_{i,r}-1,K_{i,r}+1\}\},\qquad
 K_{i,r+1}=U_{S_{i,r+1}}.
\]
When $K_{i,r}=L$, instead set
\begin{equation}
S_{i,r+1}=H_{L-1}[S_{i,r}],\qquad K_{i,r+1}=L-1;\label{eq:ct-transfer-R1a}
\end{equation}
this is a forced terminal return, and the motion below $\Gamma_{y,L}$
is deliberately not resolved.  Stop at
$N_i=\inf\{r\ge1:K_{i,r}=0\}$, so $B_{y,i}=S_{i,N_i}$.
For $i<m_0$ put
\[
C_i=\inf\{t>B_{y,i}:U_t\in\{-1,1\}\}.
\]
On $S_y$, $U_{C_i}=1$, $C_i=A_{y,i+1}$, and the connector is the
single prescribed sign $0\to1$.  Concatenating the $m_0$ block words
and these $m_0-1$ connector signs gives exactly $\overline W_y$ in
\eqref{eq:ct-transfer-39hprime}.  A block letter $j-1\to j$ pushes the next Ulam--Harris child,
a letter $j\to j-1$ pops it, and \eqref{eq:ct-transfer-R1a} closes a terminal leaf.  Hence
\begin{equation}
T_{y,1}^{1,m_0}=m_0,\qquad
 T_{y,j}^{1,m_0}=\#\{j-1\to j\hbox{ block letters}\},
                  \quad2\le j\le L.\label{eq:ct-transfer-R1b}
\end{equation}
All these times are recursively defined hitting times of closed sets;
the completed $m_0$-block word and the cap $|\overline W_y|\le N_0$
are therefore stopping-line measurable.

We next define the guards.  Consider first a competitive letter beginning
at a stopping time $s$, at level $k<L$, whose prescribed sign is
$\epsilon\in\{-1,1\}$.  Its end is
\[
t=\inf\{u>s:U_u\in\{k-1,k+1\}\},\qquad U_t=k+\epsilon.
\]
Define
\[
\begin{array}{lll}
 p=H_{k+\epsilon(1-\eta_J)}[s],
 &w=H_{k-\epsilon(1-\eta_J)}[s],
 &r=H_{k+\epsilon\eta_J}[p],\\
 e=H_{k+\epsilon(1+\eta_J)}[t],
 &o=H_{k+\epsilon\eta_J}[t].
\end{array}
\]
The bad events for this letter are
\begin{equation}
D^{\rm wrong}=\{w<t\},\qquad
 D^{\rm ret}=\{r<t\},\qquad
 D^{\rm post}=\{o<e\}.\label{eq:ct-transfer-R2}
\end{equation}
The last comparison is with the old-side collar
$k+\epsilon\eta_J$, not merely with the old nominal line $k$.
For a forced terminal return $L\to L-1$ beginning at $s$, set
\[
t=H_{L-1}[s],\quad p=H_{L-1+\eta_J}[s],\quad
 r=H_{L-\eta_J}[p],\quad
 e=H_{L-1-\eta_J}[t],\quad o=H_{L-\eta_J}[t],
\]
and use $D^{\rm ret}=\{r<t\}$ and
$D^{\rm post}=\{o<e\}$.  There is no wrong-side guard because the
path below level $L$ is not resolved.

The first source opening needs its own two guards.  The fixed annulus
defining $F_L^0$ gives a uniform initial margin.  Indeed,
$d(v,y)\le h_1/20$, convexity of $\tan$ on the relevant fixed small
interval, and $r_1=r_0e^{-1}$ imply
\[
|\chi_y(v)|=2\tan\!\left(\frac{d(v,y)}2\right)
 \le 2\tan\!\left(\frac{h_1}{40}\right)
 \le \frac{r_1}{20},
 \qquad U_0\ge1+\log20.
\]
(The lower annulus bound keeps $U_0$ finite.)  Thus, with
$c_{\rm src}:=\log20$, one has $U_0\ge1+c_{\rm src}$, and for
$\eta_J<c_{\rm src}$ the stopping time
$c_1=H_{1+\eta_J}[0]$ satisfies $c_1<A_{y,1}$.  Since the first
opening approaches level $1$ from the inner side, define
\begin{equation}
D^{A_1,{\rm pre}}=\{H_{2-\eta_J}[c_1]<A_{y,1}\},\qquad
 D^{A_1,{\rm post}}=
 \{H_{2-\eta_J}[A_{y,1}]<H_{1-\eta_J}[A_{y,1}]\}.\label{eq:ct-transfer-R2a}
\end{equation}
For $i\ge2$, the pre-opening guard at $A_{y,i}$ is exactly
$D^{\rm ret}$ for connector $i-1$, and its post-opening guard is
exactly $D^{\rm post}$ for that connector; they are named both ways
but counted once.  At every $B_{y,i}$, the last block symbol is the
competitive closing letter $1\to0$; its wrong, retreat, and post events
in \eqref{eq:ct-transfer-R2} are respectively the pre-closing and closing-recognition guards.
For $i<m_0$ the following symbol is the connector $0\to1$.  At the
last close $B_{y,m_0}$, retain its post-guard and expose it without any
following letter.  Let $D_y$ be the union of these events through the
last post-closing collar hit.  Every member of this union has the form
$\{H_a[S]<H_b[S]\}$ for a named stopping time $S$, so $D_y$ is
measurable at that finite stopping line.

We record the exact conditional ratios.  Fix an admissible complete word
$w$.  Translation, reflection, gambler's ruin, and strong Markov at the
integer-line hits give, for every competitive letter,
\begin{equation}
\frac{P(D^{\rm wrong},\overline W_y=w)}{P(\overline W_y=w)}
       =\frac{\eta_J}{2-\eta_J},\qquad
 \frac{P(D^{\rm ret},\overline W_y=w)}{P(\overline W_y=w)}
       =\frac{\eta_J(1+\eta_J)}
          {(1-\eta_J)(2-\eta_J)}.\label{eq:ct-transfer-R2b}
\end{equation}
If the next symbol is competitive, then
\begin{equation}
\frac{P(D^{\rm post},\overline W_y=w)}{P(\overline W_y=w)}
 =\begin{cases}
    \eta_J^2,&\text{if the next sign continues the present direction},\\
    2\eta_J-\eta_J^2,&\text{if it reverses it}.
   \end{cases}\label{eq:ct-transfer-R2c}
\end{equation}
Indeed, after translating the new level to $1$,
$P_1(H_{\eta_J}<H_{1+\eta_J})=\eta_J$; from $\eta_J$, the
probabilities of the next exits at $2$ and $0$ are respectively
$\eta_J/2$ and $1-\eta_J/2$, while either prescribed next sign has
probability $1/2$.  If the successor is the forced return in \eqref{eq:ct-transfer-R1a}, or
there is no successor, the post ratio is simply $\eta_J$.  For a forced
return $L\to L-1$,
\begin{equation}
\frac{P(D^{\rm ret},\overline W_y=w)}{P(\overline W_y=w)}
       =\frac{\eta_J}{1-\eta_J}.\label{eq:ct-transfer-R2d}
\end{equation}
The first-opening pre ratio in \eqref{eq:ct-transfer-R2a} is also
$\eta_J/(1-\eta_J)$.  If $d_1=-1$ is its crossing direction and
$\epsilon_1$ is the first block sign, its post ratio is
$\eta_J^2$ when $\epsilon_1=d_1$, and
$2\eta_J-\eta_J^2$ otherwise.  The same table applies at every later
opening through its incoming connector.  In particular, at $B_{y,i}$,
$i<m_0$, the close has direction $-1$ and the next connector has sign
$+1$, so its post ratio is $2\eta_J-\eta_J^2$.  At the final close
$B_{y,m_0}$ it is exactly $\eta_J$, with no fictitious next sign.
All later letters cancel by strong Markov; no arbitrary angular or future
functional is being inserted into these ratios.  A union bound, without
any independence assertion among guards, gives
\begin{equation}
P(D_y\mid\overline W_y=w)
 \le8(|w|+2m_0+1)\eta_J
 \le\varepsilon_J,
 \qquad \varepsilon_J:=CJ^4\delta_J=o(J^{-A})\label{eq:ct-transfer-R3}
\end{equation}
for every $|w|\le N_0$ and every fixed $A$.  The estimate remains
valid after conditioning on barriers, traversal counts, or any other
function of $w$, but is not claimed after conditioning on Rosen's
angular screens.

For completeness, here is the deterministic implication used below.  Let
$\mathcal A_J$ be the finite set of all nominal integer levels and all
guard-collar levels occurring in \eqref{eq:ct-transfer-R2}--\eqref{eq:ct-transfer-R2a}, including the forced-return
collars.  Let $V$ be the corresponding log-radius of the linearly
interpolated lattice path, in the coupling time parametrization.  Its
nominal source blocks and truncated word are defined at integer edge indices
by the endpoint-side-change convention above: a letter is recorded only
when an edge's endpoints change sides, and interpolation along such an edge
is used only to locate its crossing and evaluate collar separation.
Same-side tangencies and chord contacts record no letter.  Assume the
following symmetric collar condition: whenever either $U_t$ or $V_t$ is
within $2\delta_J$ of $\mathcal A_J$, both log-radii are finite and
$|U_t-V_t|\le\delta_J$.  No comparison is required while both paths lie
strictly on the unresolved inner side of the deepest collars.  Suppose $S_y$,
$|\overline W_y|\le N_0$, and $D_y^c$ occur.  Before the target collar
of a competitive Brownian letter, $V$ cannot make a side-changing crossing
of the target nominal line, while $D^{{\rm wrong},c}$ prevents it from hitting the other
neighbor.  From the first target-collar hit until the nominal target,
$D^{{\rm ret},c}$ keeps $U$ beyond the old-side collar; hence an early
side-changing target crossing by $V$ cannot be followed by a reversal.  From the nominal
hit to the far collar, $D^{{\rm post},c}$ gives the same separation,
and at the far collar
\[
\epsilon\bigl(V-(k+\epsilon)\bigr)
       \ge\eta_J-\delta_J>0.
\]
Thus $V$'s next nominal neighboring line is exactly $k+\epsilon$,
and it has made no intervening push or pop.  The forced return \eqref{eq:ct-transfer-R1a} uses
the retreat and post guards only; deeper visits are ignored by both
truncated words.

At $A_{y,1}$, \eqref{eq:ct-transfer-R2a} first prevents a premature completed neighboring
letter and then forces $V$ across the source line.  The preceding
letter argument initializes every later $A_{y,i}$ through its connector.
Induction in chronological order on the active Ulam--Harris stack now
gives the same push/pop operation at every symbol.  At $B_{y,i}$ the
same stack empties; for $i<m_0$ the guarded connector has sign $0\to1$
and starts the next stack, while the unconditional final post-guard forces
$V$ across level $0$ at $B_{y,m_0}$.  Hence the two stopped ordered
forests are canonically isomorphic and
\begin{equation}
T_{y,j}^{1,m_0}(V)=T_{y,j}^{1,m_0}(U),\qquad1\le j\le L.\label{eq:ct-transfer-R4}
\end{equation}
Moreover the exact word length is
\begin{equation}
|\overline W_y|=2\sum_{j=2}^{L}T_{y,j}^{1,m_0}+2m_0-1
       \le CL^3<N_0,\label{eq:ct-transfer-R4a}
\end{equation}
on the lower barrier and terminal band.  On the single global coupling
event \eqref{eq:ct-transfer-31}, the symmetric collar condition holds.  Indeed, every
$2\delta_J$-neighborhood of a nominal or guard collar used by the
truncated word has physical radius at least $cM$.  If either path is in
such a neighborhood, physical $b_J$-closeness, the radial Lipschitz
inequality, chart distortion, and the $O(1)$ terminal-interface rounding
give
\[
|U_t-V_t|\le C(b_J+1)/M\le\delta_J
\]
after fixing the constant in \eqref{eq:ct-transfer-33}.  No log-radius comparison is asserted or
used while both paths are deeper than the terminal collars.  Source safety
keeps every collar test before global killing.  This condition suffices:
any premature target, wrong-neighbor, or reversal crossing occurs at a
member of $\mathcal A_J$, where collar closeness and
$\delta_J<\eta_J/4$ contradict the corresponding safe guard.  This also
covers a premature $L-1$ hit during a forced $L\to L-1$ return.
Consequently
\begin{equation}
I_y^{\circ}\cap S_y\cap D_y^c\cap\{\text{global coupling \eqref{eq:ct-transfer-31}}\}
       \subseteq I_y^{S,\circ}.\label{eq:ct-transfer-R4b}
\end{equation}
Only the radial consequence is transferred: Rosen's original angular
screens remain solely on the Brownian certificate.

\smallskip\noindent
\underline{\textit{Step 3.}}\quad The upper first-crossing probability.

We now prove the rare upper estimates by direct lattice comparison.  Fix a target cell
and use the centered continuation \eqref{eq:ct-transfer-39k}--\eqref{eq:ct-transfer-39l}.  Put
\[
R_j=R_0e^{-j},\qquad R_0=3e^J,
\]
rounding every radius by at most two lattice units.  On the detailed-source
range $\|y\|>M$, define from the actual starting radius
\begin{equation}
u_y=\log {R_0\over\|y\|},\qquad
 k_y=\lceil u_y\rceil,\qquad m=k_y+1.\label{eq:ct-transfer-45a0}
\end{equation}
Then $u_y\in(k_y-1,k_y]$, $m-u_y\in[1,2)$, and the
deterministic start lies outside $D(y,R_{m-1})$.  This index differs by
only $O(1)$ from $\lfloor\log(e^J/\|y\|)\rfloor$, so it changes no
shell exponent.  Write
\begin{equation}
\begin{split}
 E_j^{\rm in}&=\{(a,b):a\notin D(y,R_j),\ b\in D(y,R_j),\ a\sim b\},\\
 E_j^{\rm out}&=\{(a,b):a\in D(y,R_j),\ b\notin D(y,R_j),\ a\sim b\}.
\end{split}\label{eq:ct-transfer-45a1}
\end{equation}
Every radial query retains its phase, level, Ulam--Harris stack and child
counters, population vector, and current directed crossing edge.

Potential-kernel optional stopping, uniformly in that edge, gives for
$k<l$
\begin{equation}
q^S_{k,l}\le C{k+1\over l},\qquad
 p_l^S(e)={l-1+O(R_l^{-1})\over l+O(R_l^{-1})},\qquad
 d_l^S(e)={1+O(R_l^{-1})\over l+O(R_l^{-1})}.\label{eq:ct-transfer-45a}
\end{equation}
Here $q^S_{k,l}$ includes the first completion at level $l$, while
$p_l^S$ and $d_l^S=1-p_l^S$ are the continuation and death rows
after return to level $l-1$.  Indeed, the two complete formulas from
Hao--Li--Okada--Zheng, \cite{HLOZ2025}, Lemma 2.1, equation (2.5), quoted in
Lemma \ref{lem:smooth-interface-discretization}, have respective logarithmic
gaps $l-1+O(R_l^{-1})$, $1+O(R_l^{-1})$, and
$l+O(R_l^{-1})$.  Thus both the continuation and the order-$1/l$
death probability are relative estimates.  Strong Markov at every
successive directed hit yields, for $N\ge1$,
\begin{equation}
P(T_{y,l}^{S,k\to0}\ge N\mid\hbox{initial directed edge})
 \le C{k+1\over l}\exp\left\{-{N\over l}+{CN\over lM}\right\}.\label{eq:ct-transfer-45b}
\end{equation}
In particular, for $X=\sqrt{2T_{y,m}}$ and $x\le Cm$,
\begin{equation}
P(X\in[x,x+1])
       \le Cm^C e^{-x^2/(2m)}.\label{eq:ct-transfer-45c}
\end{equation}
No independence of directed entrance edges is used in these scalar bounds.

A union bound over all fine centers would count many copies of the same
crossing.  We replace each center by a deterministic nearby representative
and enlarge its target disk so that the crossing count can only increase.  Put
\[
\phi(l)=\exp\{\tfrac14(l\wedge(J-l))^{1/4}\}.
\]
Fix constants
\[
0<c_{\rm sep}\le c_{\rm mesh}<C_\delta<\frac{e-1}{e+1}.
\]
For each $l$, choose a deterministic maximal
$c_{\rm sep}R_l/\phi(l)$-separated subset of $\mathbb Z^2$, with
$c_{\rm sep}$ chosen so that its covering radius (for lattice points) is
at most $c_{\rm mesh}R_l/\phi(l)$.  Given a candidate lattice center
$y$ in source shell $k$, choose a mesh point
$\widehat y=\widehat y(y,l)\in\mathbb Z^2$ with
\[
|y-\widehat y|\le c_{\rm mesh}R_l/\phi(l),\qquad
 \delta_l=C_\delta/\phi(l),\qquad R_j^\pm=(1\pm\delta_l)R_j.
\]
For all sufficiently large $J$, the gap
$(C_\delta-c_{\rm mesh})R_l/\phi(l)$ absorbs the at-most-two-unit
rounding of every radius on the used range.  Moreover
$\delta_l<(e-1)/(e+1)$, and hence
$R_{l-1}^->R_l^+$, so every buffered annulus below is nonempty.
Starting with the first entrance into
$D(\widehat y,R_k^+)$, and stopping on exiting
$D(\widehat y,R_0^+)$, let
$\widetilde T_{\widehat y,l}^{S,k\to0}$ count completed crossings from
$\partial D(\widehat y,R_{l-1}^-)$ to
$\partial D(\widehat y,R_l^+)$.  The inclusions
\begin{equation}
\begin{split}
 D(y,R_k)&\subset D(\widehat y,R_k^+),\\
 D(y,R_l)&\subset D(\widehat y,R_l^+),\\
 D(\widehat y,R_{l-1}^-)&\subset D(y,R_{l-1}),\\
 D(y,R_0)&\subset D(\widehat y,R_0^+)
\end{split}\label{eq:ct-transfer-45d0}
\end{equation}
show pathwise that
\begin{equation}
T_{y,l}^{S,k\to0}
       \le\widetilde T_{\widehat y,l}^{S,k\to0}.\label{eq:ct-transfer-45d1}
\end{equation}
Indeed every $y$-crossing begins outside the representative's buffered
outer circle, ends inside its buffered inner circle, and occurs after the
representative start and before its stop.  The same inclusions cover
$k=l-1$.  Area comparison for the mesh in a shell of diameter
$O(R_k)$ gives, for every $k\le l-1$,
\begin{equation}
\#\{\hbox{representatives at }(k,l)\}
 \le C e^{2(l-k)}\phi(l)^2.\label{eq:ct-transfer-45d}
\end{equation}
The buffered inclusions in Rosen (2.26)--(2.33) do not impose
$k\le\log J$.  For the four modified radii
$(R_k^+,R_{l-1}^-,R_l^+,R_0^+)$, the same potential-kernel
calculation gives, uniformly in all directed entrance edges,
\begin{equation}
\widetilde q_{k,l}^S\le C{k+1\over l},\qquad
 \widetilde p_l^S
 \le1-\frac{1-C/\phi(l)-C/M}{l}.\label{eq:ct-transfer-45d2}
\end{equation}
Consequently \eqref{eq:ct-transfer-45b} holds with $T$ replaced by
$\widetilde T_{\widehat y,l}^{S,k\to0}$, the additional exponent being
$CN/(l\phi(l))+CN/(lM)$.  Taking
$N=\alpha_{z,+}(l)^2/2$ there gives
\begin{equation}
P(\widetilde T_{\widehat y,l}^{S,k\to0}\ge N)
 \le C e^{-2l-(l\wedge(J-l))^{1/4}}e^{-2z}.\label{eq:ct-transfer-45e}
\end{equation}
On the used range $l\le J-(4\log J)^4$, the elementary exponent
calculation gives $l/\phi(l)=O(1)$ (the bounded early values are
absorbed into $C$), and $N/(lM)=o(1)$.  Multiplying \eqref{eq:ct-transfer-45d} and \eqref{eq:ct-transfer-45e}
and summing first in $l$, then in $k\ge s$, proves \eqref{eq:ct-transfer-40}, including
the range $\log J<k\le C_0\log J$.

\smallskip\noindent
\underline{\textit{Step 4.}}\quad The centered endpoint probability.

The source population and its descendants must be compared together.
Conditioning on a later list of entrance edges would change the law of the
current radial decision.  We therefore compare the entire chronological
word, integrating each edge when it is first produced, and only then
condition under the ideal Galton--Watson law.  In addition to
the directed-edge/stack states in \eqref{eq:ct-transfer-45a1}, introduce a special deterministic
initial atom $\star_y$, representing $S_0=0$.  The first query is
\texttt{ROOT} (hit $D_m$ before leaving $D_0$), whose kernel maps
$\star_y$ to $E_m^{\rm in}$, or \texttt{FAILURE}, whose kernel maps
$\star_y$ to the cemetery state.  After every root, its descendants are
explored depth first: at an active node of level $j<n$, \texttt{BIRTH} is the
first hit of $D_{j+1}$ and pushes the next child, while \texttt{CLOSE} is the
first exit from $D_{j-1}$ and pops the node.  A level-$n$ leaf has a
forced return to level $n-1$, of row mass one.  After a root closes, the
current state lies in $E_{m-1}^{\rm out}$; the next source query is
another \texttt{ROOT}, mapping to $E_m^{\rm in}$, or the final \texttt{DEATH}, mapping
through $E_0^{\rm out}$ to the cemetery state.  Thus a word with
$t\ge1$ roots has the
physical order
\begin{equation}
\mathtt{ROOT},\operatorname{DFS}(1),\mathtt{ROOT},
 \operatorname{DFS}(2),\ldots,\mathtt{ROOT},
 \operatorname{DFS}(t),\mathtt{DEATH}.\label{eq:ct-transfer-45f0}
\end{equation}
The word stops at the first barrier violation, endpoint decision, death,
or a fixed deterministic cap $Q_*=C_*J^3$, chosen above
$3+2\sum_{j=m}^n(T_{y,j}+1)$ on the coefficient-two barrier,
and is then padded by cemetery letters.  It contains every terminating
close and the final source death, so there is no residual future condition.

Let $P_*$ be the angle-free ideal law.  Its first root and failure
probabilities are $u_y/m$ and $(m-u_y)/m$; subsequent root and death
probabilities are $(m-1)/m$ and $1/m$; every descendant
birth/close decision has probability $1/2$.  Potential-kernel optional
stopping, with every output edge summed but every current directed edge
retained, gives the four lattice source rows
\begin{equation}
\begin{aligned}
 r_{0,y}^S
   &=\frac{u_y+O(R_m^{-1})}{m+O(R_m^{-1})},
 & f_{0,y}^S
   &=\frac{m-u_y+O(R_m^{-1})}{m+O(R_m^{-1})},\\
 r_m^S(e)
   &=\frac{m-1+O(R_m^{-1})}{m+O(R_m^{-1})},
 & d_m^S(e)
   &=\frac{1+O(R_m^{-1})}{m+O(R_m^{-1})}.
\end{aligned}\label{eq:ct-transfer-45f-src}
\end{equation}
Here the first pair starts from $\star_y$, while the second pair is
uniform in $e\in E_{m-1}^{\rm out}$.  Since every candidate satisfies
$\|y\|<2e^J$,
\[
u_y\ge\log(3/2),\qquad m-u_y\in[1,2),\qquad m-1\ge1.
\]
Also $R_m=e^{u_y-m}\|y\|\ge e^{-2}\|y\|\ge e^{-2}M$ on the detailed
range.  Consequently each of the four displayed rows, including the two
order-$1/m$ FAILURE and DEATH rows, equals its ideal row times
$1+O(M^{-1})$.  Ordinary birth/close rows are
$(1/2)(1+O(R_{j+1}^{-1}))$.

Fix a complete radial word $w$, but no directed-edge refinement.  If
$\nu_r$ is the edge-valued finite measure after its first $r$ letters,
the next ideal row $p_r$ is determined by the radial prefix, and the
uniform row estimates give
\[
e^{-C/M}p_r\|\nu_r\|\le\|\nu_{r+1}\|
       \le e^{C/M}p_r\|\nu_r\|.
\]
Forward induction, with all compatible entrance and exit edges integrated,
therefore gives
\begin{equation}
e^{-C|w|/M}P_*(w)\le P_S(w)
       \le e^{C|w|/M}P_*(w).\label{eq:ct-transfer-45f}
\end{equation}
This is not a conditional comparison given a future edge list.

Under $P_*$, the source population has the exact law
\begin{equation}
\pi_{m,u_y}(0)=1-{u_y\over m},\qquad
 \pi_{m,u_y}(t)={u_y\over m}
       \left({m-1\over m}\right)^{t-1}{1\over m},
       \quad t\ge1,\label{eq:ct-transfer-45g0}
\end{equation}
and, conditional on $t$, the $t$ descendant trees are independent
critical Galton--Watson trees with
$P(\xi=r)=2^{-r-1}$.  Hence, for the disjoint union
$\mathscr U(q)$ of complete barrier-respecting words ending in
$\sqrt{2T_n}\in[q,q+1]$,
\begin{equation}
P_S(\mathscr U(q))
 \le e^{CJ^3/M}\sum_{t\ge1}\pi_{m,u_y}(t)
       P_t^{\rm GW}(\mathscr A_t^{\rm GW}(q)).\label{eq:ct-transfer-45g}
\end{equation}
Only the ideal law on the right has been disintegrated.  Moreover, for
$x\le Cm$,
\begin{equation}
\sum_{\sqrt{2t}\in[x,x+1]}\pi_{m,u_y}(t)
 \le C(1+x)m^{-1}e^{-x^2/(2m)+O(1)}
 \le Ce^{-x^2/(2m)}.\label{eq:ct-transfer-45gprime}
\end{equation}
Apply Rosen Appendix Theorem 9.1(a), quoted in \eqref{eq:ct-external-30}, to the
ideal Galton--Watson factor in \eqref{eq:ct-transfer-45g}, group $\sqrt{2t}$ into unit bands,
and use \eqref{eq:ct-external-31}--\eqref{eq:ct-external-31a}.  This proves \eqref{eq:ct-transfer-41}.  The $t=1$ term is covered by
the extra-ancestor argument following \eqref{eq:ct-external-30}, while $t=0$ cannot reach a
positive central endpoint.  Since $J^3/M=o(J^{-A})$, the comparison is
relative even for the rare endpoint event.

\smallskip\noindent
\underline{\textit{Step 5.}}\quad The lower count and the terminal interfaces.

We apply the Brownian second moment before transferring any event to the lattice.  We retain Rosen's original event
$I_y^\circ$ and its original angular screens exactly as written in
\eqref{eq:ct-external-33}--\eqref{eq:ct-external-37}, and intersect it with the source-safety event $S_y$.
No tightened or halved screen is used.  Put
$I_y^{\rm sf}=I_y^\circ\cap S_y$.  By \eqref{eq:ct-external-38},
\begin{equation}
P\left(\sum_y1_{I_y^{\rm sf}}>0\right)\ge c_* .\label{eq:ct-transfer-46}
\end{equation}
The screens occur only in this Brownian second-moment certificate; the
lattice target discards all angular screens.

Let $\widehat I_y^B$ be the Brownian event obtained from
$I_y^\circ$ by deleting every angular screen and retaining only the
radial barrier, the bounded endpoint band, and the fixed $m_0$-source
definition.  Rosen Appendix Theorem 9.1(a), followed by the same source
and barrier calculation as \eqref{eq:ct-external-29}, gives
\begin{equation}
P(\widehat I_y^B)\le Ce^{-2L},
             \qquad \sum_yP(\widehat I_y^B)\le C.\label{eq:ct-transfer-47}
\end{equation}
Since
$I_y^{\rm sf}\cap D_y\subseteq
  \widehat I_y^B\cap S_y\cap D_y$, estimate \eqref{eq:ct-transfer-R3} may be summed over
complete safe radial words before any angular screen is imposed.  Hence
\begin{equation}
\sum_yP(I_y^{\rm sf}\cap D_y)
 \le\varepsilon_J\sum_yP(\widehat I_y^B\cap S_y)
 \le\varepsilon_J\sum_yP(\widehat I_y^B)=o(1).\label{eq:ct-transfer-47a}
\end{equation}
Let $\mathcal{K}_J$ be the single global coupling/time-cap event \eqref{eq:ct-transfer-31}.
The pathwise implication \eqref{eq:ct-transfer-R4b} gives
\[
\{Z_L^{\rm sf}>0\}\cap\mathcal{K}_J
       \setminus\bigcup_y(I_y^{\rm sf}\cap D_y)
       \subseteq\{Z_L^S>0\}.
\]
Consequently, paying the coupling failure only once,
\begin{equation}
P(Z_L^S>0)\ge P(Z_L^{\rm sf}>0)
   -\sum_yP(I_y^{\rm sf}\cap D_y)-P(\mathcal{K}_J^c)
 \ge c_*-o(1).\label{eq:ct-transfer-47b}
\end{equation}
Source safety and \eqref{eq:ct-external-38d} keep every block, connector, and final exposed
post-closing collar a fixed positive distance inside the stopping boundary,
so the close lattice path is not killed before the transferred blocks
finish.  No Brownian angular screen is transferred.
This is the only use of the two-center estimate \eqref{eq:ct-external-36}, and it occurs
entirely on the spherical Brownian side.

For the lattice first moment, extend the path if necessary and decompose
its first $m_0$ source blocks, treating every connector as an integrated
subkernel of row mass at most one.  For a fixed collection of block words,
apply the edge-marginal induction \eqref{eq:ct-smooth-39iprime} to all competitive letters,
uniformly over the entrance law supplied by each preceding connector.
Forced terminal returns have row mass one.  Thus, if $q$ is the number
of competitive decisions,
\[
P_S(\hbox{these block words})\le e^{CJ^3/M}2^{-q}.
\]
Summing the disjoint words obeying the lower barrier and terminal band gives
the critical geometric Galton--Watson probability with exactly $m_0$
ancestors.  Rosen Appendix Theorem 9.1(a), followed by the fixed-source
calculation already used in \eqref{eq:ct-transfer-47}, gives
\begin{equation}
P(I_y^{S,\circ})\le Ce^{-2L},\qquad EZ_L^S\le C.\label{eq:ct-transfer-47cprime}
\end{equation}
This proves \eqref{eq:ct-transfer-42}; in particular, the lower first moment does not use the
upper source-population estimate \eqref{eq:ct-transfer-45c}.

Finally we justify the terminal switch used in the definition of
$\mathcal{R}_y$.  On the compact center family, with
$F_y=\sigma\circ\chi_y^{-1}$,
\[
F_y(w)=F_y(0)+a_yw+b_yw^2+O(|w|^3),
       \qquad0<c\le|a_y|\le C.
\]
At $r_L\asymp e^{-J+h}$, dilation by $A_J\asymp e^J$ shows that
the actual deepest image circle is at Hausdorff distance
\begin{equation}
O(A_Jr_L^2)=O(e^{-J+2h})=o(1)\label{eq:ct-transfer-47c}
\end{equation}
from its affine concentric circle.  Rounding the center and the digital
boundary costs $O(1)$, or relative width $O(M^{-1})$, and is one
additional collar in the lower guard construction \eqref{eq:ct-transfer-R1}--\eqref{eq:ct-transfer-R4}.  The induced
angular conjugacy, after removing the fixed rotation
$a_y/|a_y|$, is $C^4$-close to the identity by
$O(e^{-J+h})$; its Haar-density defect is therefore absorbed by
$M^{-7/24}\log M+M^{-1}$.  Thus the outer $\chi_y$-rotation kernels
and the terminal Euclidean rotation kernels use the same common center
up to the stated error.  All lattice
events were defined before introducing the auxiliary coupling and are
measurable in $\mathcal{R}_y$.  We may therefore discard the Brownian
path after \eqref{eq:ct-transfer-47b}; every later conditional law is the original
simple-random-walk law given $\mathcal{R}_y$.
\end{proof}


% ---- terminal_upper ----
\section{Terminal occupation and the localized upper bound}
\label{sec:terminal-upper}

The radial traversal estimates leave the occupation accumulated inside each terminal disk to be controlled. We first compare the entrance--exit chain in a centered block, then derive its exponential occupation moment. The last subsection combines that moment with the traversal barrier to prove the localized upper estimate.

\subsection{Endpoint products in one centered block}
\begin{lemma}\label{lem:typed-centered-endpoint-product}
Let \(D=D(y,eM)\), let \(A\) be the inner vertex boundary of
\(D(y,M)\), and let \(E\) be the directed edges leaving \(D\).  Put
\[
 H^S(a,e)=P_a((S_{\tau_D-1},S_{\tau_D})=e),\qquad
 (\mathcal{H}^Sf)(a)=\sum_eH^S(a,e)f(a,e).
\]
Let \(G^S=G_D\).  For \(x,x'\in V:=D(y,M/2)\), define on
\(A\times E\)
\[
 m_x(a,e)={G^S(a,x)H^S(x,e)\over H^S(a,e)},
\]
\[
\begin{aligned}
q_{x,x'}(a,e)=\frac{1}{H^S(a,e)}\big(&G^S(a,x)G^S(x,x')H^S(x',e)\\
  &+G^S(a,x')G^S(x',x)H^S(x,e)\\
  &-\mathbf1_{\{x=x'\}}G^S(a,x)H^S(x,e)\big),
\end{aligned}
\]
and let
\begin{equation}
 \mathcal{F}_M^{\rm G}=\{m_x/C:x\in V\}\cup
 \{q_{x,x'}/[C(1+G^S(x,x'))]:x,x'\in V\},               \label{eq:tu-TP0}
\end{equation}
where the fixed \(C\) is large enough that all functions have sup norm
at most one.
Let \(R_i^S(e,da')\) be a fixed radial-only centered return word from a
previous exit edge to the next inward entrance, let
\(\sigma_i=R_i^S1\), \(J_i^S=R_i^S/\sigma_i\), and set
\begin{equation}
 P_i^S((a,e),d(a',e'))=J_i^S(e,da')H^S(a',de').           \label{eq:tu-TP1}
\end{equation}
The word may be selected adaptively from its radial prefix, but not from an
angular endpoint.  If it contains \(q_i\) elementary letters, all at scales
at least \(M\), and \(\lambda_A^S\) is the center-to-\(A\) first-hit law,
put \(\Lambda^S=\lambda_A^SH^S\).  For \(f\in\mathcal{F}_M^{\rm G}\) and
every consecutive portion \(r\le s\) of
one centered block,
\begin{equation}
 \left|\Lambda^SP_r^S\cdots P_s^Sf-\Lambda^Sf\right|
 \le C\left\{M^{-7/24}\log M+{1\over M}
                   \sum_{i=r}^s(q_i+1)\right\}.            \label{eq:tu-TP2}
\end{equation}
The bound is uniform over the fixed realized radial word.  An arbitrary
inter-block reset terminates the estimate; the next block restarts from an
arbitrary law.
\end{lemma}

\begin{proof}[Proof of Lemma~\ref{lem:typed-centered-endpoint-product}]
The comparison is applied after averaging the exit edge, which removes the
roughness of the digital boundary. We then normalize each complete return
word and compare the resulting product with rotation-invariant Brownian
kernels.

\smallskip\noindent
\underline{\textit{Step 1.}}\quad Smooth the endpoint observables.
A discrete pair is
\((a,e)\in A\times E\).  Every exterior kernel starts at the head
\(\operatorname{out}(e)\) of the preceding exit edge.  With
\(P_i^S=J_i^SH^S\), exact operator composition is
\begin{equation}
 \Lambda^SP_r^S\cdots P_s^Sf
 =\lambda_A^SH^SJ_r^SH^S\cdots J_s^S(\mathcal{H}^Sf).         \label{eq:tu-TP3}
\end{equation}
Thus the first operation on the terminal observable is the exact average
\(\mathcal{H}^S\).

The continuum entrance and exit state spaces are separate copies of
\(\mathbb T\), in the common angular coordinate of the centered block.
Let \(\bar H\) be Brownian exit from radii \(M\) to \(eM\), and let
\(\bar J_i\) be the normalized Brownian return word.  Their densities are
rotation convolutions.  Explicitly,
\begin{equation}
 {d\bar H(\alpha,\cdot)\over d\theta/(2\pi)}(\beta)
 ={e^2-1\over e^2+1-2e\cos(\beta-\alpha)},                \label{eq:tu-TP4}
\end{equation}
and an elementary side kernel from logarithmic level \(x\) to \(b\) before
\(a\) has unnormalized density
\begin{equation}
 {1\over2\pi}\left\{{x-a\over b-a}
 +2\sum_{n\ge1}{\sinh(n(x-a))\over\sinh(n(b-a))}
       \cos n(\beta-\alpha)\right\}.                      \label{eq:tu-TP5}
\end{equation}
Consequently Haar measure \(m(d\alpha)=d\alpha/(2\pi)\) is preserved
by \(\bar H\) and every \(\bar J_i\). The continuum pair transition
\(\bar J_i\bar H\) therefore preserves
\(\bar\Lambda(d\alpha,d\beta)=m(d\alpha)\bar H(\alpha,d\beta)\).

We next construct the continuum test functions.  Write \(G^S=G_D\).
For \(x,x'\in D(y,M/2)\), the exact bridge identities are
\begin{equation}
 \mathcal{H}^Sm_x(a)=G^S(a,x),                               \label{eq:tu-TP6}
\end{equation}
\begin{equation}
 \begin{split}
 \mathcal{H}^Sq_{x,x'}(a)
 &=G^S(a,x)G^S(x,x')+G^S(a,x')G^S(x',x)\\
 &\quad-\mathbf1_{\{x=x'\}}G^S(a,x).
 \end{split}                                               \label{eq:tu-TP7}
\end{equation}
Let \(a_\alpha=y+Me^{i\alpha}\) and
\[
       \mathscr G^B(a_\alpha,x)={2\over\pi}
              g_{B(y,eM)}^B(a_\alpha,x).
\]
For \(f=m_x/C\), define
\(\varphi_f(\alpha)=\mathscr G^B(a_\alpha,x)/C\).  For
\(f=q_{x,x'}/[C(1+G^S(x,x'))]\), define \(\varphi_f\) by the right side
of \eqref{eq:tu-TP7}, divided by \(C(1+G^S(x,x'))\), and replace only the boundary-variable factors
\(G^S(a,x),G^S(a,x')\) by \(\mathscr G^B(a_\alpha,x)\) and
\(\mathscr G^B(a_\alpha,x')\).  The short-distance lattice coefficients
\(G^S(x,x')\) are retained, including on the diagonal.

After scaling by \(M\),
\begin{equation}
 g_{B(0,e)}^B(z,\xi)=\log\left|{e^2-z\bar\xi\over e(z-\xi)}\right|,
       \qquad |z|=1,\quad|\xi|\le\tfrac12,               \label{eq:tu-TP8}
\end{equation}
where the logarithmic Green function is normalized as in the preceding
display. The denominators stay uniformly
away from zero, and symmetry of \(G^S\) bounds all retained normalized
coefficients.  Hence
\begin{equation}
                  \sup_f\|\varphi_f\|_{C^4(\mathbb T)}\le C.          \label{eq:tu-TP9}
\end{equation}

We use the Green-function comparison of Kozdron--Lawler
\cite[Corollary~3.5, equation~(30)]{KozdronLawler2005}.
For a simply connected lattice domain \(A\) of inradius in \([n,2n]\),
let \(\widetilde A\) be its union-of-squares domain. If
\(u\in A_n^*:=\{z:g_{\widetilde A}(0,z)\ge n^{-1/16}\}\), the
comparison, away from the diagonal, is
\begin{equation}
 G_A(u,z)={2\over\pi}g_{\widetilde A}(u,z)+k_{z-u}
                 +O(n^{-7/24}\log n),                    \label{eq:tu-TP10}
\end{equation}
where \(a\) is the lattice potential kernel and
\(k_w=k_0+(2/\pi)\log|w|-a(w)\) for \(w\ne0\).
Apply this with the core pole in \(D(y,M/2)\). This pole is a fixed
fraction of the radius from the boundary and belongs to \(A_n^*\).
The boundary-variable point is at distance at least \(M/2-O(1)\) from
the pole, so \(k_{a-x}=O(M^{-2})\). Sandwiching the union-of-squares
disk between disks whose radii differ by \(O(1)\), and projecting
radially to the continuum circle, each costs \(O(M^{-1})\).
Equations
\eqref{eq:tu-TP6}--\eqref{eq:tu-TP10} therefore give, with
\(U_AF(a)=F(\arg(a-y))\),
\begin{equation}
 \|\mathcal{H}^Sf-U_A\varphi_f\|_\infty
       \le C\epsilon_G(M),\qquad
 \epsilon_G(M)=M^{-7/24}\log M.                           \label{eq:tu-TP11}
\end{equation}

For an exit edge put \(U_EF(e)=F(\arg(\operatorname{out}(e)-y))\).
Lemma~\ref{lem:smooth-interface-discretization}, with the same tubular
continuation for a disk having only its exit boundary, gives for every
smooth angular test \(F\)
\begin{equation}
 \|H^SU_EF-U_A\bar HF\|_\infty
       +|\lambda_A^SU_AF-mF|
       \le {C\over M}\|F\|_{C^4}.                         \label{eq:tu-TP12}
\end{equation}
For the second term, solve the disk Dirichlet problem with boundary data
\(F\), start at its center, and apply the tubular continuation argument
of \eqref{eq:ct-smooth-39e}--\eqref{eq:ct-smooth-39f}. For the first term, solve the exit-disk problem and
start on the inner circle. This proves the comparison for the digital
circles directly.

\smallskip\noindent
\underline{\textit{Step 2.}}\quad Normalize each complete return word.
A prescribed word may have very small probability. The elementary
relative estimates must therefore be multiplied before the complete word
is normalized. Write its elementary kernels as
\[
 T_{i,l}^S(z,dz')=p_{i,l}^S(z)K_{i,l}^S(z,dz'),\qquad
 \bar T_{i,l}=\bar p_{i,l}\bar K_{i,l},\qquad1\le l\le q_i,
\]
and let \(Q_i^S=K_{i,1}^S\cdots K_{i,q_i}^S\).  Under this normalized
digital chain put
\[
 W_i=\prod_{l=1}^{q_i}{p_{i,l}^S(Z_{l-1})\over\bar p_{i,l}}.
\]
By \eqref{eq:ct-smooth-39c}, if \(\eta_i=C\sum_lR_{i,l}^{-1}\le Cq_i/M\), then
\(e^{-\eta_i}\le W_i\le e^{\eta_i}\).  With
\(\bar s_i=\prod_l\bar p_{i,l}\), the identities
\begin{equation}
 {R_i^S1(e)\over\bar s_i}=E_{Q_{i,e}^S}W_i,\qquad
 J_i^SF(e)={E_{Q_{i,e}^S}[W_iF(Z_{q_i})]
                 \over E_{Q_{i,e}^S}W_i}                  \label{eq:tu-TP13}
\end{equation}
are exact.  Hence
\begin{equation}
             \|J_i^SF-Q_i^SF\|_\infty
             \le C\eta_i\|F\|_\infty.                    \label{eq:tu-TP14}
\end{equation}

Apply the normalized smooth-test estimate following \eqref{eq:ct-smooth-39c} to the
\(K_{i,l}^S\) and telescope from the right.  If
\(\Phi_{l,F}=\bar K_{i,l+1}\cdots\bar K_{i,q_i}F\), convolution gives
\(\|\Phi_{l,F}\|_{C^4}\le\|F\|_{C^4}\).  With the appropriate angular
pullback at every state,
\[
 \|Q_i^SU_AF-U_E\bar J_iF\|_\infty
       \le C\eta_i\|F\|_{C^4}.
\]
Together with \eqref{eq:tu-TP14},
\begin{equation}
 \|J_i^SU_AF-U_E\bar J_iF\|_\infty
       \le C\eta_i\|F\|_{C^4}.                            \label{eq:tu-TP15}
\end{equation}
The weight in \eqref{eq:tu-TP13} contains the row masses of every letter of the fixed
word, including letters after the state currently being compared.

\smallskip\noindent
\underline{\textit{Step 3.}}\quad Compare the complete product.

Write \eqref{eq:tu-TP3} as alternating operators
\(H^S,J_r^S,H^S,\ldots,J_s^S\), with terminal test
\(g_f^S=\mathcal{H}^Sf\).  Let \(L_j^S,\bar L_j\) be the corresponding
discrete and Brownian operators. Write \(U_j\) for the angular pullback
on the state space after the \(j\)-th operator, and put
\(\Phi_{j,f}=\bar L_{j+1}\cdots\bar L_N\varphi_f\).  The operator identity, with every pullback on its corresponding state
space, is
\begin{equation}
\begin{split}
 L_1^S\cdots L_N^Sg_f^S-U_0\bar L_1\cdots\bar L_N\varphi_f
 &=L_1^S\cdots L_N^S(g_f^S-U_N\varphi_f)\\
 &\quad+\sum_{j=1}^NL_1^S\cdots L_{j-1}^S
 [L_j^SU_j\Phi_{j,f}-U_{j-1}\bar L_j\Phi_{j,f}].          \end{split}
\label{eq:tu-TP16}
\end{equation}
All \(\Phi_{j,f}\) have uniformly bounded \(C^4\)-norm by \eqref{eq:tu-TP9} and
convolution contraction.  Equations \eqref{eq:tu-TP11}, \eqref{eq:tu-TP12}, and \eqref{eq:tu-TP15} bound
\eqref{eq:tu-TP16} by the right side of \eqref{eq:tu-TP2}.  In the continuum the alternating
product preserves Haar flux, so it has integral \(m\varphi_f\).  A
separate zero-length use of \eqref{eq:tu-TP11}--\eqref{eq:tu-TP12} gives
\(|\Lambda^Sf-m\varphi_f|\le C(\epsilon_G+M^{-1})\).
Subtracting these two comparisons proves \eqref{eq:tu-TP2}.

For a lower conformal block the same proof uses its one common
\(\chi_y\)-angle.  Each occurrence of \(H^S\) has angular conjugacy error
\(O(M^{-1}+e^{-J+h})\).  Since \(2h\le J\),
\(e^{-J+h}\le M^{-1}\), and paying this error once per occurrence is
already contained in \(M^{-1}\sum_i(q_i+1)\). A decision depending on
an angular endpoint can change Haar measure. Such decisions are therefore
included in the reset density constructed below, at the boundary between
centered blocks.
\end{proof}

\subsection{Occupation of completed terminal excursions}
Fix a center \(y\) and set \(o_y:=y\).  Let
\[
\begin{aligned}
 D_y&=\{z\in\mathbb Z^2:\|z-y\|<eM\},\\
 A_y&=\partial_{\rm in}\{z:\|z-y\|<M\},\\
 V_y&=\{z:\|z-y\|<M/2\},
\end{aligned}
\]
where \(\partial_{\rm in}D=\{a\in D:\exists a'\notin D,\ a'\sim a\}\),
and let \(E_y=\{(u,v):u\in D_y,v\notin D_y,u\sim v\}\) be the directed
exit edges.  A completed
terminal piece begins at the first lattice vertex \(a\) after an inward
crossing of radius \(M\), and is killed on the directed edge \(e\in E_y\)
that first crosses radius \(eM\).  Erasing the interiors of all such
outward pieces gives exactly the intrinsic sigma-field \(\mathcal{R}_y\) of
Lemma~\ref{lem:buffered-transfer}.  Conditional on their endpoint pairs
\(\omega_i=(a_i,e_i)\), their occupation vectors
\(X_i=(X_i(x):x\in V_y)\) are independent killed lattice bridges.

An \emph{admissible centered block word} \(w\) is a positive-probability
radial-only cylinder.  It has \(b=O(J^2)\) completed pieces, at most
\(Q\le J^3\) adjacent-shell decisions, and a fixed
number \(G_0\) of blocks.  Within each block every decision is between
circles centered at \(y\), or, in the lower spherical transfer, between successive level
sets of the single conformal coordinate \(\chi_y\).  In either case the
normalized Brownian kernels are rotation convolutions in one common
angular coordinate; this is what ``centered block'' means here.  Between
two lower blocks the safe connector contributes no observable and its
normalized angular kernel is treated only as an arbitrary reset.  Its
scalar prescribed-side row mass is retained in the exact word
factorization, but no comparison of its angular density or invariant law
is made.  Before factorization, retain only source intervals containing
at least one completed terminal piece as blocks.  Fold any initial empty
intervals and connectors into the first chronological law, every
intervening string of empty intervals and connectors into the preceding
reset, and every terminal such string into the final reset.  Hence
\(b_g\ge1\), and in the lower application \(G_0\le m_0\); the
inter-source connectors and empty intervals contribute no core
occupation.  In the upper application \(G_0=1\) when \(b>0\), after the
centered extension \eqref{eq:ct-transfer-39k}.  The case \(b=0\) is trivial and is henceforth
excluded.  The reset construction below assigns every initial/final
suffix and connector exactly once, so there is no unrecorded future
weight.  Put
\[
 H(a,e)=P_a\big((S_{\tau_{D_y^c}-1},S_{\tau_{D_y^c}})=e\big).
\]
Write \(b_g\) for the number of completed pieces in block \(g\), so
\(\sum_gb_g=b\).  For \(i<b_g\), let
\(R_{g,i}^w(e,da')\) be the sub-Markov kernel of the prescribed exterior
return word, started at \(\operatorname{out}(e_{g,i})\), from
\(e_{g,i}=e\) to the next entrance
\(a_{g,i+1}=a'\).  The pair subkernel factors exactly as
\begin{equation}
 Q_{g,i}^w((a,e),d(a',e'))=
       R_{g,i}^w(e,da')H(a',de').                           \label{eq:tu-47c}
\end{equation}
Write
\begin{equation}
 \sigma_{g,i}(e)=R_{g,i}^w1(e),\qquad
 J_{g,i}^w={R_{g,i}^w\over\sigma_{g,i}},
 \qquad\widetilde K_{g,i}^w=J_{g,i}^wH.                    \label{eq:tu-47d}
\end{equation}
For \(g<G_0\), let
\begin{equation}
 C_g^w(e,d\zeta,d\omega')                                  \label{eq:tu-47d-prime}
\end{equation}
be the genuine sub-Markov reset kernel which, after the last endpoint
of block \(g\), contains the prescribed final no-more-piece suffix, the
safe connector and its data \(\zeta\), and the prefix of block \(g+1\)
through its first endpoint pair \(\omega'\).  For \(g=G_0\), the same
notation denotes the final suffix to a cemetery state.  Put
\begin{equation}
 \sigma_{g,b_g}(e)=C_g^w1(e),\qquad
 J_{g,b_g}^w=C_g^w/\sigma_{g,b_g}.                          \label{eq:tu-47d-doubleprime}
\end{equation}
In a centered suffix every adjacent-shell letter is included in \(Q\).
A lower inter-source connector is not declared centered; its normalized
kernel \(J_{g,b_g}^w\) is used as an arbitrary reset.

Let \(\alpha_{1,w}\) be the actual subprobability law of the first endpoint
pair after imposing only the word prefix through that piece.  Admissibility
means \(\alpha_{1,w}1>0\); put
\(\mu_{1,w}=\alpha_{1,w}/(\alpha_{1,w}1)\).  The
\emph{sequentially normalized} law \(\widetilde P_w\) is the
joint law of all endpoint and reset variables that starts with
\(\mu_{1,w}\), uses \(\widetilde K_{g,i}^w\) internally, and uses the
normalized reset kernels \(J_{g,b_g}^w\) between blocks and at the final
cemetery transition.  Thus the first law of every later block is
generated chronologically by the preceding reset, not conditioned on
its future output.  Write \(P_w\) for the true endpoint-and-reset law
conditioned on the complete word \(w\).  All estimates are uniform over
the arbitrary normalized reset kernels, which are marginalized at the
end.

Let \(\lambda_A\) be the law of the first hit of \(A_y\) by a walk started
at \(o_y\), and set
\begin{equation}
                 \Lambda_M(da,de)=\lambda_A(da)H(a,de).     \label{eq:tu-47e}
\end{equation}
For a pair observable \(f\), write
\begin{equation}
                 (\mathcal{H}f)(a)=\sum_eH(a,e)f(a,e).         \label{eq:tu-47f}
\end{equation}
For \(x,x'\in V_y\), define
\begin{equation}
 m_x(a,e)={G(a,x)H(x,e)\over H(a,e)},                       \label{eq:tu-47g}
\end{equation}
\begin{equation}
\begin{aligned}
q_{x,x'}(a,e)=\frac{1}{H(a,e)}\big(&G(a,x)G(x,x')H(x',e)\\
&+G(a,x')G(x',x)H(x,e)\\
&-{\bf1}_{\{x=x'\}}G(a,x)H(x,e)\big).
\end{aligned}                                             \label{eq:tu-47h}
\end{equation}
Uniformly,
\begin{equation}
 0\le m_x\le C,\qquad0\le q_{x,x'}\le C(1+G(x,x')),       \label{eq:tu-47h-prime}
\end{equation}
and the exit-edge averages are exactly
\[
 \mathcal{H}m_x(a)=G(a,x),
\]
\begin{equation}
 \mathcal{H}q_{x,x'}(a)=G(a,x)G(x,x')+G(a,x')G(x',x)
             -{\bf1}_{\{x=x'\}}G(a,x).                    \label{eq:tu-47h-doubleprime}
\end{equation}
The finite normalized observable family is
\begin{equation}
 \mathcal{F}_M=\{m_x/C:x\in V_y\}\cup
 \left\{{q_{x,x'}\over C(1+G(x,x'))}:x,x'\in V_y\right\}; \label{eq:tu-47i}
\end{equation}
it has at most \(CM^4\) members.  Put
\begin{equation}
 \varepsilon_M=CM^{-7/24}\log M,
 \qquad
 \eta_w^*:=CQ/M+CG_{\rm nc},                                \label{eq:tu-48a}
\end{equation}
where \(G_{\rm nc}\le G_0\le m_0\) counts noncentered whole resets; each
is factored once through the separator \eqref{eq:tu-59a0} below.  Thus
\(G_{\rm nc}=0\) in the upper application and is fixed in the lower one.
Take the deterministic centers \(c_f=\Lambda_Mf\).  Before the
normalizations in \eqref{eq:tu-47i}, they satisfy
\begin{equation}
 c_{m_x}=g+O(\varepsilon_M),                                \label{eq:tu-48b}
\end{equation}
\begin{equation}
 c_{q_{x,x'}}=gG(x,x')+gG(x',x)-{\bf1}_{\{x=x'\}}g
       +O(\varepsilon_M(1+G(x,x'))).                        \label{eq:tu-48c}
\end{equation}
\begin{lemma}\label{lem:terminal-feynman-kac}
With the preceding notation, the exact density of the complete word is
\begin{equation}
 {dP_w\over d\widetilde P_w}
 ={\displaystyle\prod_{g=1}^{G_0}\prod_{i=1}^{b_g}
          \sigma_{g,i}(e_{g,i})\over
   \displaystyle E_{\widetilde P_w}
       \prod_{g=1}^{G_0}\prod_{i=1}^{b_g}
          \sigma_{g,i}(e_{g,i})},
\qquad
 e^{-2\eta_w^*}\le
 {dP_w\over d\widetilde P_w}
 \le e^{2\eta_w^*}.                                        \label{eq:tu-48d}
\end{equation}
The product visibly includes the final no-more-piece suffix and every
connector likelihood.  For every \(t>0\), conditionally on \(w\),
\begin{equation}
\begin{split}
P_w\Bigg(&\max_{f\in\mathcal{F}_M}\left|
 b^{-1}\sum_{i\le b}f(\omega_i)-c_f\right|\\
&>t+C\{G_0/b+\varepsilon_M+Q/M\}\Bigg)
 \le2|\mathcal{F}_M|\exp\{2\eta_w^*-cbt^2\}.                   \end{split}
\label{eq:tu-49}
\end{equation}

For \(x\in V_y\), \(G=G_{D_y}\), and
\begin{equation}
 r_\theta(x):=
 \frac{e^\theta-1}
 {1-(G(x,x)-1)(e^\theta-1)},                                 \label{eq:tu-39}
\end{equation}
uniformly for \(0\le r_\theta\le c/h\), \(h=\log M\),
\begin{equation}
 \log E\left[\exp\left\{\theta\sum_{i=1}^bX_i(x)\right\}
             \,\middle|\,w\right]
       =bg\,r_\theta+O\left(br_\theta^2+
       r_\theta\{G_0+b\varepsilon_M+bQ/M\}
                    +\eta_w^*\right).                       \label{eq:tu-50}
\end{equation}
The estimate is uniform in \(x,w,b\) and concerns completed pieces.
In both applications below, every visit to \(V_y\) lies in a completed
piece, so no conditional incomplete-piece estimate is needed.

For the upper-bound choice from Lemma~\ref{lem:buffered-transfer},
\[
 h=(C_h\log J)^4,\qquad n=J-h+O(1),\qquad
 U=u_{J,c},\quad c>2,
\]
put \(Y=\sqrt{2b}\).  On a coefficient-two coarse barrier, uniformly over
all \(0\le Y\le 3n\),
\begin{equation}
\begin{split}
 P\left(\max_{x\in V_y}\sum_iX_i(x)\ge U\mid w,b\right)
 &\le C e^{-2n}J^c e^{b/n}\\
 &\quad{}\times
 \exp\left\{-\kappa\min\left(
 \frac{(p-p_*)^2}{h},\frac{J^2}{h}\right)\right\},            \end{split}
\label{eq:tu-51}
\end{equation}
whenever \(Y\) is in the central range.  Here
\[
 p=B_n+1-Y,\qquad B_n=\rho_Jn+z_0+h^{1/4},\qquad
 p_*=O(\log J).
\]
Here the coefficient-two barrier means
\begin{equation}
 \sqrt{2T^S_{y,j}}\le
 B_j:=\rho_Jj+z_0+(j\wedge(J-j))^{1/4}
 \quad(k_y<j\le n),                                         \label{eq:tu-52}
\end{equation}
after the fixed \(O(1)\) relaxation of Lemma~\ref{lem:buffered-transfer}.  The central range is
\(\epsilon n\le Y\le B_n+1\), for one fixed small
\(\epsilon>0\).  For \(Y<\epsilon n\), \eqref{eq:tu-51} holds with its last factor
replaced by \(\exp(-\kappa J^2/h)\).
\end{lemma}

\begin{proof}[Proof of Lemma~\ref{lem:terminal-feynman-kac}]
The bridge formula reduces the occupation moment to a product of bounded
endpoint observables. We control that product under the sequential law,
where centered transitions contract, and then restore the likelihood of
the complete word. The likelihood comparison must include the last
suffix and every reset, because conditioning on their future outcomes
can otherwise change the endpoint distribution.

\smallskip\noindent
\underline{\textit{Step 1.}}\quad Compute bridge moments and their centers.
Let
\[
 H(a,e)=P_a\big((S_{\tau-1},S_\tau)=e\big),\qquad
 \tau=\tau_{D_y^c}.
\]
For the bridge from \(a\) to \(e\), write \(X(x)\) for its occupation of
\(x\).  Splitting at the first visit to \(x\), and then summing the
geometric number of returns to \(x\), gives the exact rank-one identity
\begin{equation}
\begin{split}
 H_\theta(a,e)
 &:=E_a[e^{\theta X(x)};(S_{\tau-1},S_\tau)=e]\\
 &=H(a,e)+r_\theta G(a,x)H(x,e).                              \end{split}
\label{eq:tu-42}
\end{equation}
Thus
\begin{equation}
       \frac{H_\theta(a,e)}{H(a,e)}
       =1+r_\theta m_x(a,e),\qquad
       m_x(a,e):=\frac{G(a,x)H(x,e)}{H(a,e)}.                 \label{eq:tu-43}
\end{equation}
A Harnack chain inside the fixed-ratio disk gives \(0\le m_x\le C\),
uniformly in \(x\in V_y\). Although \(m_x(a,e)\) is indexed by the
rough digital exit edges, the transition \eqref{eq:tu-47c} first averages over that
edge. From \eqref{eq:tu-43},
\begin{equation}
 (\mathcal{H}m_x)(a)=\sum_eH(a,e)m_x(a,e)=G(a,x).               \label{eq:tu-43a}
\end{equation}
More generally, if \(q_{x,x'}(a,e)=E^{a,e}[X(x)X(x')]\), the
two ordered visit-time decompositions give
\begin{equation}
\begin{split}
 (\mathcal{H}q_{x,x'})(a)
 &=G(a,x)G(x,x')+G(a,x')G(x',x)\\
 &\quad-{\bf1}_{\{x=x'\}}G(a,x).                            \end{split}
\label{eq:tu-43b}
\end{equation}
Thus the only functions subsequently fed into a return-word kernel are
Green functions with their singularities a distance at least \(M/2\)
from \(A_y\).

We next identify the common center and the only observables that require
comparison.  With \(\lambda_A\) defined as the first-hit law of the
inner vertex boundary, the strong Markov property gives
\begin{equation}
 d_x:=\sum_a\lambda_A(a)G_{D_y}(a,x)
 =G_{D(y,eM)}(o_y,x)-G_{D(y,M)}(o_y,x)
                  +O(M^{-1}\log M).                         \label{eq:tu-44}
\end{equation}
Indeed the one-edge collar can return from \(A_y\) to \(D(y,M/2)\)
before exiting \(D(y,M)\) with probability \(O(M^{-1})\), while the
Green function is at most \(C\log M\).  This remainder is
\(o(\varepsilon_M)\).

Apply the Green comparison \eqref{eq:tu-TP10} to both disks in \eqref{eq:tu-44}. For distinct
arguments it reads
\begin{equation}
 G_A(u,z)={2\over\pi}g_A(u,z)+k_{z-u}
                  +O(n^{-7/24}\log n),\label{eq:tu-45}
\end{equation}
with \(k_w=k_0+(2/\pi)\log|w|-a(w)\) for \(w\ne0\).
The lattice correction is identical in the two domains and cancels.
When \(x=o_y\), use instead the diagonal formula in
\cite[Theorem~1.2]{KozdronLawler2005}; its additive potential-kernel
constant also cancels. The core points lie a fixed fraction of the radius
from both boundaries, and the union-of-squares domains are between disks
whose radii differ by \(O(1)\). The Brownian Green-function difference
for the concentric disks is exactly \(\log e=1\). Consequently, uniformly
in \(x\in V_y\),
\begin{equation}
                     d_x=g+O(\varepsilon_M),
 \qquad \varepsilon_M=CM^{-7/24}\log M.                    \label{eq:tu-46}
\end{equation}
The identities \eqref{eq:tu-43a}--\eqref{eq:tu-43b} now give the centers \eqref{eq:tu-48b}--\eqref{eq:tu-48c}.

\smallskip\noindent
\underline{\textit{Step 2.}}\quad Compare and contract centered transitions.
The state of
\(R_{g,i}^w\) is the preceding directed exit edge, and the walk restarts at
its head; its target is the next inward directed edge.  A word is called
admissible only when its cylinder has positive probability, so every row
normalization below is defined.  Brownian row masses of its elementary
letters are scalar, while Lemma~\ref{lem:smooth-interface-discretization} gives
their relative lattice comparison.  The exact whole-word weighting formula
\eqref{eq:tu-TP13} therefore yields, for a scalar \(s_{g,i}>0\),
\begin{equation}
 \sup_e\left|{\sigma_{g,i}(e)\over s_{g,i}}-1\right|
       \le {Cq_i\over M}.                                  \label{eq:tu-58c}
\end{equation}
This estimate holds for the complete word, since \eqref{eq:tu-TP13} retains the row
masses of all its letters.

In the common logarithmic conformal coordinate, every normalized Brownian
letter is the rotation convolution \eqref{eq:tu-TP5}, and the terminal disk exit is
\eqref{eq:tu-TP4}.  Apply Lemma~\ref{lem:typed-centered-endpoint-product} to the exact
exit-edge averages \eqref{eq:tu-43a}--\eqref{eq:tu-43b}.  For consecutive internal transitions
\(r\le s\) in one block and every \(f\in\mathcal{F}_M\), it gives
\begin{equation}
 \left|\Lambda_M
   \widetilde K_{g,r}^w\widetilde K_{g,r+1}^w\cdots
   \widetilde K_{g,s}^wf-c_f\right|
 \le C\left\{\varepsilon_M+
       {1\over M}\sum_{i=r}^s(q_i+1)\right\}.              \label{eq:tu-58e}
\end{equation}
The two terms are compared separately with the same Haar integral of the
disk-Green traces \eqref{eq:tu-TP6}--\eqref{eq:tu-TP9}. The estimate therefore controls complete
prefixes even though \(\Lambda_M\) need not be invariant for a single
discrete transition. In a lower source interval the common coordinate is
\(\chi_y\); its terminal chart error is included in the displayed sum.
An inter-source connector ends this comparison.

To remove dependence on the initial law of each block, we also need a
uniform contraction in total variation. The elliptic Harnack inequality
\cite[Theorem~1.1]{AthreyaEtAl2022} states that for every \(d\ge1\)
there is \(C_d\) such that a nonnegative function harmonic in
\(B(x_0,R)\) satisfies
\begin{equation}
       \max_{B(x_0,R/2)}u\le C_d\min_{B(x_0,R/2)}u.           \label{eq:tu-58a}
\end{equation}
Starting one step inside the outer boundary is too close to use \eqref{eq:tu-58a}
directly.  Insert instead the deterministic intermediate circle
\[
                 C_M=\partial D(y,e^{1/2}M),
\]
which every successful outer-to-inner crossing must hit.  For a target
atom \(a\in A_y\), put
\[
 h_a(z)=P_z(S_{H_{A_y}}=a,\ H_{A_y}<\tau_{D(y,eM)}),
 \qquad s(z)=\sum_{a\in A_y}h_a(z),\qquad z\in C_M.
\]
A bounded chain of lattice balls of radius \(cM\), all a fixed positive
fraction from both boundary components, applies \eqref{eq:tu-58a} to \(h_a\) and
\(s\).  Uniformly in \(a,z,z_0\in C_M\),
\[
 ch_a(z_0)\le h_a(z)\le Ch_a(z_0),\qquad
 cs(z_0)\le s(z)\le Cs(z_0).
\]
After division by the total successful-leg probability this gives
\begin{equation}
       J_{\rm in}(z,\cdot)\ge\alpha
                     J_{\rm in}(z_0,\cdot)                  \label{eq:tu-58b}
\end{equation}
where \(J_{\rm in}(z,\cdot)=h_{\cdot}(z)/s(z)\) is the entrance
law conditional on reaching the inner boundary, and \(\alpha>0\) is
uniform in \(M,z,z_0\). Precomposition by all earlier
within-block word factors and postcomposition by \(H\) preserve this row-independent
minorization. Writing \(\operatorname{osc}(f)=\sup f-\inf f\), we obtain
\(\operatorname {osc}(\widetilde K_{g,i}^wf)\le(1-\alpha)
\operatorname {osc}(f)\) for every bounded \(f\) at every within-block
transition.  Write \(\rho=1-\alpha<1\).  A reset transition is used
exactly and has Dobrushin coefficient at most one; there are only
\(G_0-1\) of them.

\smallskip\noindent
\underline{\textit{Step 3.}}\quad Retain the full likelihood at resets.
A connector can have an arbitrary normalized angular kernel. What is
needed is a bound on the oscillation of its row mass, uniform even when
the prescribed future event is rare. To obtain it, factor each reset
through a deterministic separator. Fix
\begin{equation}
        \lambda_*=e^{3/2}\in(e,e^2),\qquad
        \Sigma_y=\partial D(o_y,\lambda_*M).                \label{eq:tu-59a0}
\end{equation}
The next active outer interface after the terminal \(eM\)-circle has scale
\(e^2M\).  The fixed-gap geometry and \eqref{eq:tu-47c} put \(\Sigma_y\) strictly
between them and keep \(D(o_y,\lambda_*M)\) inside the applicable source
interval and global killing domain.

Let \(e=(u,v)\in E_y\) be the last exit edge of a nonempty block, so the
exterior kernel starts at \(v=\operatorname{out}(e)\) and
\(|v-o_y|=eM+O(1)\).  Put
\[
 \theta_-=H_{D(o_y,M)},\qquad
 \theta_+=\tau_{D(o_y,\lambda_*M)},
\]
and on \(\{\theta_+<\theta_-\}\) let \(F_+\) be the directed outward
edge at \(\theta_+\).  For \(g<G_0\), let
\(B_g^w(f,d\zeta,d\omega')\) be the genuine post-separator subkernel
which realizes the rest of the no-more-piece suffix, all folded empty
source intervals and safe connector signs, and the next nonempty block
prefix through its first endpoint pair.  For \(g=G_0\) it instead contains
the remaining suffix and final cemetery transition.  Durations, parities,
angles, and already visible edge data are outputs of this kernel, never
extra scalar conditions.  Set \(r_g^w(f)=B_g^w(f,\hbox{all outputs})\).
The strong Markov property gives the exact factorization
\begin{equation}
 C_g^w(e,d\gamma)=\sum_f
 P_v(\theta_+<\theta_-,F_+=f)B_g^w(f,d\gamma).              \label{eq:tu-59a1}
\end{equation}
Indeed, a return to \(D(o_y,M)\) before \(\Sigma_y\) after the alleged
last completed piece would, on the subsequent trip to the block's outer
endpoint, force another exit of \(D_y\), hence another completed terminal
piece.  Conversely, no reset restriction other than avoiding the inner
circle is tested before \(\Sigma_y\); all remaining restrictions occur in
\(B_g^w\).

Define on the digital annulus
\[
 \Omega_y=D(o_y,\lambda_*M)\setminus D(o_y,M)
\]
\begin{equation}
 h_g^w(z)=E_z[r_g^w(F_+);\theta_+<\theta_-].                \label{eq:tu-59a2}
\end{equation}
This is a nonnegative discrete harmonic function at every interior vertex,
with exact Dirichlet value zero on the inner boundary and the possibly rough
edge data \(r_g^w\) on the outer boundary.  Global killing and the entire
later connector/prefix occur only in that outer boundary datum.  Moreover,
\begin{equation}
                  \sigma_{g,b_g}(e)=C_g^w1(e)=h_g^w(v).     \label{eq:tu-59a3}
\end{equation}
All such \(v\)'s lie in the band
\(\bigl|\|v-o_y\|-eM\bigr|\le2\), a distance
comparable to \(M\) from both components of \(\partial\Omega_y\).
A fixed number of overlapping lattice balls of radius \(cM\), whose doubled
balls remain in \(\Omega_y\), connects any two points of this band.
Applying \eqref{eq:tu-58a} along this chain gives
\begin{equation}
       \sup_e\sigma_{g,b_g}(e)
              \le C_H\inf_e\sigma_{g,b_g}(e).              \label{eq:tu-59a}
\end{equation}
If \(h_g^w\) vanishes at one interior start, the mean-value property and
connectedness make it identically zero, and the complete word cylinder has
zero mass and is discarded.  Otherwise every row is positive and the
normalization \(J_{g,b_g}^w=C_g^w/\sigma_{g,b_g}\) is defined.  This proves
\eqref{eq:tu-59a} independently of how rare or rough the later event is.

Start the first block with its actual chronological law through the first
endpoint.  Use every normalized internal kernel and every normalized reset
in chronological order, including the final cemetery output, to define
\(\widetilde P_w\).  Direct disintegration of the unnormalized cylinder
measure gives
\begin{equation}
 {dP_w\over d\widetilde P_w}
 ={\displaystyle\prod_{g=1}^{G_0}\prod_{i=1}^{b_g}
          \sigma_{g,i}(e_{g,i})\over
   \displaystyle E_{\widetilde P_w}
       \prod_{g=1}^{G_0}\prod_{i=1}^{b_g}
          \sigma_{g,i}(e_{g,i})}.                           \label{eq:tu-60a}
\end{equation}
The initial prefix is absorbed into \(\mu_{1,w}\), and the product
includes the final suffix. Equation \eqref{eq:tu-58c}, applied also to a centered final suffix, bounds the
logarithmic row oscillation of centered factors by \(Cq_i/M\); \eqref{eq:tu-59a} costs
\(\log C_H\) once per genuinely noncentered reset.  Thus in the upper
centered extension \(G_{\rm nc}=0\), while in the lower construction
\(G_{\rm nc}\le G_0\le m_0\).  Scalar
constants cancel between numerator and denominator, so, after absorbing
\(\log C_H\) into \(C\),
\begin{equation}
 \left|\log{dP_w\over d\widetilde P_w}\right|
 \le2\{CQ/M+CG_{\rm nc}\}=2\eta_w^*.                       \label{eq:tu-60b}
\end{equation}
\smallskip\noindent
\underline{\textit{Step 4.}}\quad Concentrate the endpoint averages.
Fix \(w\) and a normalized
\(f\in\mathcal{F}_M\), expose the sequential joint chain under
\(\widetilde P_w\), and use the Doob martingale of
the whole sum \(\sum_{i\le b}f(\omega_i)\).  Changing \(\omega_i\) changes the
conditional expectation of the future sum by at most
\[
       2G_0\sum_{j\ge0}\rho^j\le {2G_0\over1-\rho};
\]
the factor \(G_0\) accounts for the possibility that a reset preserves
all the remaining variation.
Azuma--Hoeffding consequently gives
\begin{equation}
 \widetilde P_w\left(\left|\sum_{i\le b}f(\omega_i)
       -E_{\widetilde P_w}\sum_{i\le b}f(\omega_i)\right|>bt\right)
 \le2e^{-cbt^2}.                                             \label{eq:tu-61}
\end{equation}
More explicitly, if the \(k\)-th endpoint of a block has actual initial
law \(\mu\), contraction and the full prefix comparison \eqref{eq:tu-58e} give
\begin{equation}
\begin{split}
\left|\mu\widetilde K_{g,1}^w\cdots
 \widetilde K_{g,k-1}^wf-c_f\right|
 &\le2\rho^{k-1}\\
 &\quad+C\left\{\varepsilon_M+
 {1\over M}\sum_{j<k}(q_j+1)\right\}.                      \end{split}
\label{eq:tu-61a0}
\end{equation}
For \(k=1\) this just uses \(\|f\|_\infty\le1\).
Summing \eqref{eq:tu-61a0} over all endpoints and blocks, and using
\(\sum_j(q_j+1)\le C Q\), gives
\begin{equation}
 \left|E_{\widetilde P_w}\sum_{i\le b}f(\omega_i)-bc_f\right|
 \le C\{G_0+b\varepsilon_M+bQ/M\}.                         \label{eq:tu-61a}
\end{equation}
A union bound under \(\widetilde P_w\), followed by \eqref{eq:tu-48d}, proves
\eqref{eq:tu-49}.

\smallskip\noindent
\underline{\textit{Step 5.}}\quad Bound the exponential occupation moment.
Conditional on the endpoints, \eqref{eq:tu-42}--\eqref{eq:tu-43} makes the exponential moment
\(\prod_i(1+rm_x(\omega_i))\).  Under \(\widetilde P_w\), put
\begin{equation}
                 \psi_i=\log(1+rm_x(\omega_i)).              \label{eq:tu-62}
\end{equation}
Since \(0\le m_x\le C\),
\(\psi_i=rm_x(\omega_i)+O(r^2)\) and
\(\operatorname {osc}(\psi_i)\le Cr\).  The same influence calculation
as for \eqref{eq:tu-61} shows that every Doob martingale difference for
\(\sum_i\psi_i\) is at most \(CG_0r/(1-\rho)\).  Conditional Hoeffding and
Jensen therefore give
\begin{equation}
 0\le\log E_{\widetilde P_w}e^{\sum_i\psi_i}
       -E_{\widetilde P_w}\sum_i\psi_i\le Cbr^2.             \label{eq:tu-62a}
\end{equation}
The exact exit-edge average and its Taylor remainder are
\begin{equation}
 \sum_eH(a,e)\log(1+rm_x(a,e))=rG(a,x)+O(r^2),             \label{eq:tu-62b}
\end{equation}
because \(m_x\le C\) and \(\sum_eH(a,e)m_x(a,e)=G(a,x)\le C\).
The product telescoping \eqref{eq:tu-61a}, now for \(m_x/C\), and \eqref{eq:tu-46} give
\begin{equation}
 \sum_{i\le b}E_{\widetilde P_w}m_x(\omega_i)
       =bg+O\{G_0+b\varepsilon_M+bQ/M\}.                    \label{eq:tu-62c}
\end{equation}
Combining \eqref{eq:tu-62}--\eqref{eq:tu-62c} proves the completed-piece version of \eqref{eq:tu-50} under
\(\widetilde P_w\).  Finally,
\eqref{eq:tu-48d} puts the exact full-word density \eqref{eq:tu-60a} between
\(e^{-2\eta_w^*}\) and \(e^{2\eta_w^*}\).  Multiplying the nonnegative
exponential by this global density and taking logarithms changes the
answer by at most \(2\eta_w^*\), exactly the last error in \eqref{eq:tu-50}.
This proves \eqref{eq:tu-50} under \(P_w\).

\smallskip\noindent
\underline{\textit{Step 6.}}\quad Optimize the terminal upper tail.
The first choice of tilt will cancel the endpoint Gaussian cost. A further
adjustment supplies decay in the distance from the preferred endpoint,
which is needed when summing all endpoint bands. The inverse of \eqref{eq:tu-39} is
\begin{equation}
       \theta(r)=
       \log\frac{1+G(x,x)r}{1+(G(x,x)-1)r}.                   \label{eq:tu-49-2}
\end{equation}
Take \(r_0=(gn)^{-1}\).  Since
\(G(x,x)=gh+O(1)\) and \(h^2/J=o(1)\),
\begin{equation}
 \theta(r_0)=\frac1{gJ}+O(J^{-2}+h^2J^{-3}),\qquad
 \theta(r_0)U=2J-c\log J+O(1),\qquad bg r_0=b/n.              \label{eq:tu-50-2}
\end{equation}
Since \(G_0=1\), \(b=O(J^2)\), \(Q\le J^3\), and
\(M=e^{(C_h\log J)^4}\), every error term in \eqref{eq:tu-50}
is \(O(1)\) at \(r_0\), uniformly in the word.
A union bound over \(O(e^{2h})\) sites, \eqref{eq:tu-50}, and Chernoff's inequality
give the first line of \eqref{eq:tu-51}.

We retain the tilt freedom to sum every endpoint band.  Set
\[
       F_b(r)=\theta(r)U-bgr-C_0br^2.
\]
Here are the derivative and feasibility details.  Uniformly on every fixed
interval \(0\le r\le a/h\),
\[
 \theta'(r)=\frac1{(1+Gr)(1+(G-1)r)},
\]
\[
 -\theta''(r)=\theta'(r)
 \left\{\frac G{1+Gr}+\frac{G-1}{1+(G-1)r}\right\}.
\]
Thus \(c_a\le\theta'(r)\le C_a\),
\(c_ah\le-\theta''(r)\le C_ah\), and
\(|\theta'''(r)|\le C_ah^2\).  Since \(b=O(J^2)\), after increasing
\(J\),
\begin{equation}
       -C_1J^2h\le F_b''(r)\le-c_1J^2h.                     \label{eq:tu-51a}
\end{equation}
Also
\begin{equation}
 \theta'(r_0)=1-\frac{2h}{n}
                  +O(n^{-1}+h^2n^{-2}).                     \label{eq:tu-51b}
\end{equation}
Define \(b_*\) by \(F_{b_*}'(r_0)=0\).  Since
\[
       b_*={U\theta'(r_0)\over g+2C_0r_0},
\]
the relations \(J=n+h+O(1)\) and \(h^2=o(J)\) give
\[
 b_*=2n^2-cn\log J+O(n+h^2+h\log J)
     =2n^2-cn\log J+O(n).
\]
Consequently
\begin{equation}
       Y_*=\sqrt{2b_*}
          =2n-\frac c2\log J+O(1),                            \label{eq:tu-51-2}
\end{equation}
and hence \(p_*=B_n+1-Y_*=O(\log J)\).  More exactly,
\begin{equation}
 F_b'(r_0)=(g+2C_0r_0)(b_*-b)
          =(g+o(1)){(Y_*-Y)(Y_*+Y)\over2}.                  \label{eq:tu-51c}
\end{equation}

Fix a sufficiently small constant \(\delta>0\).  If
\(\epsilon n\le Y\le B_n+1\) and \(|Y-Y_*|\le\delta J\), Taylor's
formula, \eqref{eq:tu-51a}, and the displacement
\[
       s={F_b'(r_0)\over C_1J^2h}
\]
(with \(C_1\) enlarged once) give
\begin{equation}
 \sup_rF_b(r)-F_b(r_0)
 \ge\kappa\frac{(Y-Y_*)^2}{h}.                              \label{eq:tu-51d}
\end{equation}
This displacement is feasible.  If it is negative, the barrier and
\(h^{1/4}=O(\log J)\) imply \(Y-Y_*=O(\log J)\), whence
\(|s|=O(\log J/(Jh))=o(r_0)\).  If it is positive, choosing \(\delta\)
small ensures \(s<a/(2h)\).  When
\(\epsilon n\le Y\le Y_*-\delta J\), choose instead a fixed
\(s=a_0/h>0\), with \(a_0\) small compared with \(\delta/C_1\).
Equations \eqref{eq:tu-51a}--\eqref{eq:tu-51c} then give
\begin{equation}
       F_b(r_0+s)-F_b(r_0)\ge\kappa J^2/h.                  \label{eq:tu-51e}
\end{equation}
There is no range \(Y\ge Y_*+\delta J\) because the upper barrier permits
only an \(O(\log J)\) overshoot above \(Y_*\).

Finally, if \(Y<\epsilon n\), take \(r=a_0/h\) directly.  The integral
formula \(\theta(r)=\int_0^r\theta'(t)\,dt\) gives
\(\theta(a_0/h)\ge c_aa_0/h\).  First choosing \(\epsilon\) and then
\(a_0\) small yields
\begin{equation}
       F_b(a_0/h)\ge\kappa_0J^2/h,
       \qquad F_b(r_0)=O(J)=o(J^2/h).                        \label{eq:tu-51f}
\end{equation}
Equations \eqref{eq:tu-51d}--\eqref{eq:tu-51f}, and
\(p-p_*=Y_*-Y\), prove precisely the two alternatives in \eqref{eq:tu-51}.

\smallskip\noindent
\underline{\textit{Step 7.}}\quad Account for the initial and final portions.
In the applications, all core occupation belongs to completed pieces.  A lower source block begins and ends on interfaces outside
\(D_y\).  An upper block with \(k\le C_0\log J\) starts at the origin,
which is outside \(D_y\), and ends at \(\sigma_y\), also outside
\(D_y\).  Before the first inward crossing of radius \(M\) the path
cannot visit \(V_y\).  After any visit to \(V_y\), reaching the block's
outer endpoint forces a later crossing of radius \(eM\), so that inward
crossing and all intervening core occupation form a completed piece.
Thus the initial and final incomplete portions contribute exactly zero
to every \(X(x)\) used here.  Upper cells with \(k>C_0\log J\), including
the cells near the walk's starting point, are treated separately by the
one-site Green estimate in \eqref{eq:tu-62-2}.
\end{proof}

\subsection{A local upper bound for the maximum}
\begin{proposition}\label{prop:localized-disk-upper}
For every fixed \(\beta>0,c>2\) and \(A<\infty\),
\begin{equation}
\begin{split}
 P\left(\max_{\lVert x\rVert\le e^JJ^{-\beta}}
             L_{\tau_{e^J}}(x)\ge u_{J,c}\right)
 \le J^{-2\beta+c-2+o(1)}+O(J^{-A}).                         \end{split}
\label{eq:tu-53}
\end{equation}
The \(o(1)\) is uniform when \(\beta,c\) range over compact subsets of
\((0,\infty)\) and \((2,\infty)\), respectively.
\end{proposition}

\begin{proof}[Proof of Proposition~\ref{prop:localized-disk-upper}]
We work with one terminal cell at a time and extend the walk to a disk
centered at that cell. This puts the source law, traversal barrier, and
terminal word in the same centered geometry. The traversal estimate and
the terminal exponential moment then have opposite Gaussian factors;
after their cancellation, the ballot prefactor supplies the required
power \(J^{-2}\).

\smallskip\noindent
\underline{\textit{Step 1.}}\quad Choose cells and exclude barrier crossings.
Use the upper buffer
\[
       h=(C_h\log J)^4,\qquad M=e^h,\qquad n=J-h,
\]
with \(C_h\) fixed large.  Cover the target disk by original \(M\)-cells
obtained from the dilated image of Rosen's predetermined net \(F_n\),
rounding each image center to its nearest lattice point.  The chart has
uniformly bounded distortion on the relevant compact patch, so these
centers are \(cM\)-separated and form a \(CM\)-cover.  After altering the
fixed core and terminal radius factors, every lattice point lies in a
core, the overlap multiplicity is bounded, and Lemmas
\ref{lem:buffered-transfer} and \ref{lem:terminal-feynman-kac}
apply to the same cells.  We continue to denote their radii by \(M/2\) and \(eM\).
For a cell center with \(\|y\|>M\), use the exact index from \eqref{eq:ct-transfer-45a0},
\begin{equation}
 u_y=\log {3e^J\over\|y\|},\qquad
 k=k(y)=\lceil u_y\rceil,\qquad m=k+1.                     \label{eq:tu-54a}
\end{equation}
Cells with \(\|y\|\le M\) are assigned to the final sitewise range
\(k>C_0\log J\).  Uniform distortion of the fixed chart makes \eqref{eq:tu-54a}
equal to the spherical source index \eqref{eq:ct-external-24} up to a bounded additive error,
which is absorbed by the barrier slack.
Every target cell has
\begin{equation}
       k\ge s:=\beta\log J+O(1).                              \label{eq:tu-54}
\end{equation}
For this fixed cell continue the walk from \(\tau_{e^J}\) to
\(\sigma_y=\tau_{D(y,3e^J)}\).  By \eqref{eq:ct-transfer-39l}, a high local time in the cell
before \(\tau_{e^J}\) is also high before \(\sigma_y\).  All counts,
words, terminal bridges, and the source index below refer to this single
\(y\)-centered extended path.  This domination is applied separately to
each cell before the final union bound.

First exclude a crossing of the fixed positive-intercept
coefficient-two barrier \eqref{eq:tu-52} at or before level \(n\).  Since
\(h>(4\log J)^4\) after choosing \(C_h\) large, all these levels lie in
the range of \eqref{eq:ct-transfer-45d}--\eqref{eq:ct-transfer-45e}.  Restricting the source shells to \(k\ge s\)
and using \eqref{eq:ct-transfer-40} gives
\begin{equation}
\begin{split}
 P(\hbox{some bad prefix in the target})
 &\le C\sum_{k\ge s}\sum_{l=k+1}^{n}
 e^{2(l-k)+\frac12 l_J^{1/4}}
 e^{-2l-l_J^{1/4}}\\
 &\le J^{-2\beta+o(1)},\qquad
 l_J:=l\wedge(J-l).                                          \end{split}
\label{eq:tu-55}
\end{equation}
The direct lattice estimates \eqref{eq:ct-transfer-45a}--\eqref{eq:ct-transfer-45e}, for both the representative
count and the probability of a crossing at one center, apply through the
source-shell range \(k\le C_0\log J\).

\smallskip\noindent
\underline{\textit{Step 2.}}\quad Include the source law in the ballot estimate.
Fix one original cell in a shell
\[
       s\le k\le C_0\log J.
\]
Sum the complete words that include the source and whose radial prefix
stays below the barrier. The comparison is applied to this joint event
before conditioning on a particular source height.  Put \(m=k+1\) and
\[
       X=\sqrt{2T_m},\qquad Y=\sqrt{2T_n}.
\]
The directed-edge source estimate \eqref{eq:ct-transfer-45c} gives, for each unit band \(I_x\),
\begin{equation}
       P(X\in I_x)\le C(k+1)^C\exp\{-x^2/(2m)\}.             \label{eq:tu-56}
\end{equation}
For a unit endpoint band \(I_q=[q,q+1]\), let \(\mathscr A(q)\) denote
the disjoint union of complete words that retain the good barrier and end
in \(Y\in I_q\).  The joint source-inclusive comparison \eqref{eq:ct-transfer-45g}, before any
disintegration of the actual lattice law, gives for every source band
\(I_x\)
\begin{equation}
\begin{split}
 &P_S\bigl(X\in I_x,\mathscr A(q)\bigr)\\
 &\qquad\le e^{CJ^3/M}
   \sum_{\sqrt{2t}\in I_x}\pi_{m,u_y}(t)
       P_t^{\rm GW}(\mathscr A_t^{\rm GW}(q)).              \end{split}
\label{eq:tu-56a}
\end{equation}
Only the ideal law on the right is disintegrated by \(t\).  The remaining
traversal chain has length \(L=n-m\).  Its affine barrier has endpoints
\(a=B_m\) and \(b=B_n+1\), and \eqref{eq:ct-external-31a} shows that it dominates the truncated
\(J\)-barrier without replacing \(\rho_J\) by \(\rho_n\).

Fix \(\epsilon_0>0\) small and
\(\epsilon_0n\le q\le B_n+1\), and put \(p=B_n+1-q\).  For each
\(t\ge2\), set \(x_t=\sqrt{2t}\) and apply \eqref{eq:ct-external-30} pointwise with initial
height \(x_t\), endpoint band \([q,q+1]\), and duration \(L\).  On the
good prefix, \(x_t\le B_m=O(m)\), while the exact geometric source law
satisfies
\[
       \pi_{m,u_y}(t)\le Cm^{-1}\exp\{-x_t^2/(2m)\}.
\]
Because \(m+L=n\), the two exponential costs combine by the exact identity
\begin{equation}
 \frac{x_t^2}{m}+\frac{(q-x_t)^2}{L}
 =\frac{q^2}{n}+\frac{n(x_t-mq/n)^2}{mL}.                   \label{eq:tu-57}
\end{equation}
The identity uses the actual height \(x_t\). Grouping those heights into
unit bands changes the exponential by only \(e^{O(1)}\).  The Gaussian sum in the
last term and the remaining source factors cost only a polynomial in
\(m\).  Since \(q\ge\epsilon_0n\) and \(L\asymp n\), the ballot prefactor
in \eqref{eq:ct-external-30}, including \(\sqrt{x_t/(qL)}\), gives
\begin{equation}
 P(\hbox{prefix good},\,Y\in I_q)
 \le C(k+1)^C(1+p)n^{-2}\exp\{-q^2/(2n)\}.                  \label{eq:tu-58}
\end{equation}
For the only positive start below two, namely \(T_m=1\), add an
independent ancestor and require that ancestor to die in its first
generation.  This has probability \(1/2\), so the desired probability is
at most twice the bound from initial population two: from generation
one onward the two paths agree, and these are exactly the generations on
which the barrier in \eqref{eq:ct-external-30} is imposed.  Initial population zero cannot
produce the present positive endpoint.  Thus
\eqref{eq:tu-58} proves the central-range estimate \eqref{eq:ct-transfer-41}. Its relative lattice
comparison remains valid after summing over the cells.

\smallskip\noindent
\underline{\textit{Step 3.}}\quad Combine the traversal and terminal costs.
The complementary endpoint range is controlled directly by the terminal
tilt.  If
\(Y<\epsilon_0n\), then \(b\le\epsilon_0^2n^2/2+O(n)\), and the saturated
line of \eqref{eq:tu-51}, uniformly in the radial word, gives
\begin{equation}
 P(\hbox{terminal mark}\mid w)
 \le C e^{-2n}J^c
       \exp\{\epsilon_0^2n/2-\kappa J^2/h+O(1)\}.            \label{eq:tu-58a-2}
\end{equation}
After summing this over at most \(Ce^{2n}\) original cells, the result is
\(o(J^{-A})\) for every fixed \(A\), because \(h=o(J)\), after fixing
\(\epsilon_0\) sufficiently small.  This is why neither \eqref{eq:ct-external-29} nor \eqref{eq:ct-transfer-41}
was claimed in the extinction range.

For the number \(b=T_n+O(1)\) of completed outward pieces in the band
\(I_q\),
\[
       b=q^2/2+O(n).
\]
Conditional on the radial word, Lemma~\ref{lem:terminal-feynman-kac} applies;
the start/end observation at the end of that lemma shows that all visits
to \(V_y\) belong to its completed pieces.  Multiply
\eqref{eq:tu-58} by \eqref{eq:tu-51}.  The factors
\(\exp\{-q^2/(2n)\}\) and \(\exp\{b/n\}\) cancel up to a constant.
Moreover,
\begin{equation}
 \sum_{0\le p\le B_n+1}(1+p)
 \exp\left\{-\kappa\min\left(
          \frac{(p-p_*)^2}{h},\frac{J^2}{h}\right)\right\}
       =J^{o(1)}.                                             \label{eq:tu-59}
\end{equation}
Indeed the nonsaturated Gaussian sum is
\[
 O\bigl((1+|p_*|)\sqrt h+h\bigr)=J^{o(1)},
\]
while its saturated complement contributes
\(O(J^2)e^{-\kappa J^2/h}=o(1)\).  The small-\(Y\) part has the
saturated factor in \eqref{eq:tu-51}.  Thus, uniformly
over all endpoint bands and all radial words,
\begin{equation}
 P(\hbox{this original cell is good and marked})
       \le (k+1)^Ce^{-2n}J^{c-2+o(1)}.                        \label{eq:tu-60}
\end{equation}
The terminal maximum has already included all sites of the original
cell, so the remaining union bound is over the original cells.

There are \(O(e^{2(n-k)})\) original cells in shell \(k\).  Summing \eqref{eq:tu-60}
and using \eqref{eq:tu-54} yields
\begin{equation}
 \sum_{k=s}^{C_0\log J}e^{2(n-k)}
       (k+1)^Ce^{-2n}J^{c-2+o(1)}
       \le J^{-2\beta+c-2+o(1)}.                              \label{eq:tu-61-2}
\end{equation}

\smallskip\noindent
\underline{\textit{Step 4.}}\quad Treat cells near the starting point.
It remains to treat \(k>C_0\log J\). For a lattice site in shell \(k\),
Green-function identities and Lemma~\ref{lem:elementary-walk-estimates} give
\[
 P_0(H_x<\sigma_y)\le C(k+1)/J
\]
and, conditionally on the hit,
\[
 P_x(L_{\sigma_y}(x)\ge u_{J,c})
       \le Ce^{-2J}J^c.
\]
There are \(O(e^{2J-2k})\) sites in that shell.  Their total contribution
is
\begin{equation}
       J^{c-1-2C_0+o(1)}.                                    \label{eq:tu-62-2}
\end{equation}
Choose \(C_0\) after \(A,c\).  Combining \eqref{eq:tu-55}, \eqref{eq:tu-61-2}, and \eqref{eq:tu-62-2} proves
\eqref{eq:tu-53}.  The word-cap and digital row errors are superpolynomial by
Lemma~\ref{lem:buffered-transfer} and are included in \(O(J^{-A})\).
\end{proof}


% ---- terminal_lower ----
\section{The terminal maximum and the lower bound}
\label{sec:terminal-lower}

The coarse construction produces a cell with about
$2L^2-2L\log L$ completed terminal pieces.  Its mean occupation is
still below the desired level because the construction stops at depth
$L=J-h$.  We now recover the missing contribution by maximizing over
the sites inside the terminal disk.  Conditional on the entrance and
exit data, the pieces are independent, and the covariance of their
occupation field is close to a multiple of the killed Green function.
A Gaussian approximation therefore reduces the remaining estimate to
the maximum of the discrete Gaussian free field.

\subsection{Conditional occupation in one terminal disk}

Use the terminal domains $D_y$ and $V_y$, the endpoint law $\Lambda_M$,
and the completed-piece occupation vectors $X_i$ defined above.  Within
this subsection, write $G=G_{D_y}$.  For an endpoint pair $\omega=(a,e)$,
set
\[
 \mu_x(\omega)=E^\omega X(x),\qquad
 q_{x,x'}(\omega)=E^\omega[X(x)X(x')].
\]
The first task is to control these conditional moments simultaneously
over the terminal disk.  This is the reason for taking a smaller
terminal logarithmic scale here than in the upper bound.

\begin{lemma}[The terminal maximum]
\label{lem:terminal-maximum-lower}
Fix $1<p<2$ and put
\[
 h=(\log J)^p,\qquad M=e^h,\qquad L=J-h.
\]
Suppose an admissible centered block word has $G_0\le m_0$ blocks and
\begin{equation}
 b=2L^2-2L\log L+O(L)
 \label{eq:terminal-piece-count}
\end{equation}
completed terminal pieces, with a uniform constant in the remainder.
Put $t_J=\log J/(J\log\log J)$.  Call the endpoint realization
\emph{moment regular} if, for every $x,x'\in V_y$,
\begin{equation}
 \begin{split}
 \left|\frac1b\sum_{i\le b}\mu_x(\omega_i)-\Lambda_M\mu_x\right|
 &\le 2t_J,\\
 \left|\frac1b\sum_{i\le b}q_{x,x'}(\omega_i)
                    -\Lambda_Mq_{x,x'}\right|
 &\le 2t_J\bigl(1+G(x,x')\bigr).
 \end{split}
 \label{eq:terminal-moment-regularity}
\end{equation}
Uniformly over these words,
\begin{equation}
 P(\text{moment regularity fails}\mid w)
 \le CM^4e^{-cbt_J^2}=o(J^{-A})
 \label{eq:terminal-moment-failure}
\end{equation}
for every fixed $A>0$.

Let $\mathcal C_y$ be the intrinsic coarse sigma-field, which records
the word and its endpoint pairs and erases the bridge interiors.
For every moment-regular realization, under the conditional law given
$\mathcal C_y$,
\begin{equation}
 \max_{x\in V_y}\sum_{i\le b}X_i(x)
 =gb+2\sqrt2\,g\sqrt b\,h
 +O_P\left(J\sqrt h+\frac{J\log J}{\log\log J}\right).
 \label{eq:terminal-maximum-expansion}
\end{equation}
The remainder is uniform over such realizations.  In particular, for
every fixed $\eta>0$, the conditional probability that this maximum
is at least $u_{J,2+\eta}$ tends to one uniformly.
\end{lemma}

\begin{proof}
We first identify the conditional means and covariances.  We then
approximate the centered maximum by a Gaussian maximum and compare
the latter with the discrete Gaussian free field.  The scale
$1<p<2$ makes both the simultaneous moment screen and the final
error estimate compatible with the $J\log J$ margin.

\smallskip\noindent
\underline{\textit{Step 1.}}\quad Conditional moments and their common center.
Fix the entrance and directed-exit pairs $(a_i,e_i)_{i\le b}$.
The strong Markov property makes the killed bridge occupation vectors
independent under this conditioning.  Ordering the two visit times
gives the exact identities
\begin{equation}
 E^{a,e}X(x)=\frac{G(a,x)H(x,e)}{H(a,e)}
 \label{eq:terminal-exact-mean}
\end{equation}
and
\begin{equation}
 \begin{split}
 E_a[X(x)X(x');E=e]
 &=G(a,x)G(x,x')H(x',e)\\
 &\quad+G(a,x')G(x',x)H(x,e)
 -\mathbf1_{\{x=x'\}}G(a,x)H(x,e).
 \end{split}
 \label{eq:terminal-exact-second-moment}
\end{equation}
The subtraction accounts for the visit counted in both orders when
$x=x'$.  Dividing \eqref{eq:terminal-exact-second-moment} by $H(a,e)$
gives $q_{x,x'}(a,e)$.

The fixed-ratio Harnack estimates yield
\begin{equation}
 \|\mu_x/C\|_\infty\le1,\qquad
 \left\|\frac{q_{x,x'}}{C(1+G(x,x'))}\right\|_\infty\le1
 \label{eq:terminal-normalized-moments}
\end{equation}
uniformly in $x,x'$.  The exact exit-edge averages \eqref{eq:tu-43a}--\eqref{eq:tu-43b}
identify these functions with the finite Green class \eqref{eq:tu-47i}, so the
word concentration estimate \eqref{eq:tu-49} applies.  There are at most $CM^4$
functions.  Moreover, admissibility and the present scale imply
\begin{equation}
 G_0/b+\varepsilon_M+Q/M=o(t_J),\qquad \eta_w^*=O(1).
 \label{eq:terminal-screen-errors}
\end{equation}
Apply \eqref{eq:tu-49} with $t=t_J/C_1$, where $C_1$ is a sufficiently large fixed
multiple of the normalization constant in \eqref{eq:tu-47i}.  After restoring
the normalizations in \eqref{eq:terminal-normalized-moments}, the
tolerances are those in \eqref{eq:terminal-moment-regularity}.  Thus
\begin{equation}
 \begin{split}
 P(\text{\eqref{eq:terminal-moment-regularity} fails}\mid w)
 &\le CM^4\exp\left\{-c\frac{\log^2J}{(\log\log J)^2}\right\}\\
 &=o(J^{-A})
 \end{split}
 \label{eq:terminal-screen-union-bound}
\end{equation}
for every fixed $A>0$, because
$h=o(\log^2J/(\log\log J)^2)$.  The estimate is uniform over the
normalized reset kernels; these are included in the joint law and
then integrated out.  This proves \eqref{eq:terminal-moment-failure}.
Only the lower bound uses this simultaneous $M^4$-function screen.
At the upper-bound scale $h=(C_h\log J)^4$, Lemma~\ref{lem:terminal-feynman-kac}
uses the one-point exponential estimate \eqref{eq:tu-50} before taking its
spatial union bound.

To compute the center, multiply
\eqref{eq:terminal-exact-mean}--\eqref{eq:terminal-exact-second-moment}
by $\Lambda_M(a,e)=\lambda_A(a)H(a,e)$ and sum first over $e$.
The Green identities \eqref{eq:tu-44}--\eqref{eq:tu-45} give
\begin{equation}
 \begin{split}
 \Lambda_M\mu_x&=d_x,\\
 \Lambda_Mq_{x,x'}
 &=d_xG(x,x')+d_{x'}G(x',x)-\mathbf1_{\{x=x'\}}d_x,\\
 \sup_{x\in V_y}|d_x-g|&\le C\varepsilon_M.
 \end{split}
 \label{eq:terminal-common-moments}
\end{equation}
In particular, a moment-regular realization satisfies
\begin{equation}
 \max_{x\in V_y}
 \left|\frac1b\sum_{i\le b}E^{\omega_i}X_i(x)-g\right|
 \le Ct_J+O(\varepsilon_M).
 \label{eq:terminal-conditional-mean}
\end{equation}
Write
\[
 \widetilde X_i(x)=X_i(x)-E^{\omega_i}X_i(x),\qquad
 U_x=b^{-1/2}\sum_{i\le b}\widetilde X_i(x).
\]
Since $\sup_{x,x'}G(x,x')\le Ch$, the second-moment error in the
screen is $O(t_J(1+h))=o(1)$.  The error from replacing $d_x$ by $g$
in \eqref{eq:terminal-common-moments} is $O(\varepsilon_Mh)=o(1)$.
The diagonal correction is bounded, as is
\[
 \frac1b\sum_{i\le b}
 E^{\omega_i}X_i(x)E^{\omega_i}X_i(x')
\]
by \eqref{eq:terminal-normalized-moments}.  Consequently
\begin{equation}
 \max_{x,x'\in V_y}
 \left|\operatorname{Cov}(U_x,U_{x'})-2gG_{D_y}(x,x')\right|
 \le C.
 \label{eq:terminal-covariance-comparison}
\end{equation}
The exact exit-edge averaging is what supplies the common center;
no angular interpolation of a directed-edge observable is required.

\smallskip\noindent
\underline{\textit{Step 2.}}\quad Approximation of the maximum by a Gaussian maximum.
The number of coordinates is $P=|V_y|=e^{2h+O(1)}$, whereas there
are only $b\asymp J^2$ independent vectors.  We therefore need an
approximation whose dimension dependence is logarithmic.

We use Chernozhukov--Chetverikov--Kato
\cite[Theorem~3.1]{CCK2016}.  In the centered form needed here, let
$Z_i\in\mathbb R^P$ be independent centered vectors with finite
absolute third moments, and let $Y_i$ be independent centered
Gaussian vectors with the same covariances.  Set
\[
 Z=\max_{j\le P}b^{-1/2}\sum_{i\le b}Z_{ij},\qquad
 \widetilde Z=\max_{j\le P}b^{-1/2}\sum_{i\le b}Y_{ij}.
\]
For every Borel set $A\subset\mathbb R$ and $\delta>0$,
\begin{equation}
 P(Z\in A)\le P(\widetilde Z\in A^{C_7\delta})
 +\frac{C_8\log^2P}{\delta^3\sqrt b}
   \{L_b+M_{b,Z}(\delta)+M_{b,Y}(\delta)\},
 \label{eq:terminal-cck}
\end{equation}
where $A^r=\{s:\operatorname{dist}(s,A)\le r\}$, the constants
are universal, and
\[
 \begin{split}
 L_b&=\max_{j\le P}\frac1b\sum_{i\le b}E|Z_{ij}|^3,\\
 M_{b,Z}(\delta)
 &=\frac1b\sum_{i\le b}E\left[
   \max_{j\le P}|Z_{ij}|^3;
   \max_{j\le P}|Z_{ij}|>\frac{\delta\sqrt b}{\log P}\right].
 \end{split}
\]
The definition of $M_{b,Y}$ is the same with $Y$ in place of $Z$.
We apply this statement conditionally, with $Z_i=\widetilde X_i$.

The rank-one identity \eqref{eq:tu-42} gives the needed tail estimate with a
factor $h^{-1}$.  Indeed, at $\theta=c_0/h$, with $c_0>0$ fixed
and sufficiently small, \eqref{eq:tu-39} gives $r_\theta\le C/h$, and hence
\[
 E^\omega(e^{\theta X(x)}-1)
 =r_\theta\mu_x(\omega)\le C/h.
\]
For $s\ge h$, exponential Markov's inequality gives
$P^\omega(X(x)>s)\le Ch^{-1}e^{-cs/h}$.  To see the same scale
for smaller $s$, the generating function in \eqref{eq:tu-42} identifies
$X(x)$ as zero with probability $1-\mu_x(\omega)/G(x,x)$ and,
otherwise, as a geometric random variable of mean $G(x,x)$.
Since $G(x,x)\asymp h$ on $V_y$, we obtain, uniformly in the
endpoint pair,
\begin{equation}
 P^\omega(X(x)>s)\le Ch^{-1}e^{-cs/h},\qquad
 E^\omega|X(x)-E^\omega X(x)|^3\le Ch^2.
 \label{eq:terminal-bridge-tails}
\end{equation}
Here and below $s\ge0$.  Integrating the tail proves the moment
bound, and therefore $L_b=O(h^2)$.

For $\delta=1$, the truncation threshold is
\[
 T=\frac{\sqrt b}{\log P}\asymp\frac Jh.
\]
A union bound over $P$ coordinates, followed by integration of
\eqref{eq:terminal-bridge-tails}, gives
\begin{equation}
 \frac1b\sum_{i\le b}E^{\omega_i}\left[
   \max_{x\in V_y}|\widetilde X_i(x)|^3;
   \max_{x\in V_y}|\widetilde X_i(x)|>T\right]
 \le \exp\{-cJ/h^2+Ch\}.
 \label{eq:terminal-truncated-third-moment}
\end{equation}
The polynomial factors from integration are absorbed by the negative
exponent.  The matching Gaussian vectors have coordinate variances
at most $Ch$, by the second derivative of \eqref{eq:tu-42}.  Their contribution
is bounded by $\exp\{-cJ^2/h^3+Ch\}$ after the same union bound.
Thus the error in \eqref{eq:terminal-cck} is at most
\begin{equation}
 O(h^4/J)+\exp\{-cJ/h^2+Ch\}=o(1).
 \label{eq:terminal-cck-error}
\end{equation}
All estimates depend on the endpoint data only through the uniform
bounds already established.

\smallskip\noindent
\underline{\textit{Step 3.}}\quad The Gaussian maximum.
Let $\Gamma=(\Gamma_x:x\in V_y)$ be centered Gaussian with the
covariance of $U$, and let $\phi$ be the zero-boundary discrete
Gaussian free field in $D_y$, so that
\[
 E[\phi_x\phi_{x'}]=G_{D_y}(x,x').
\]
The extremal-process result of Biskup--Louidor
\cite[(1.2), (1.6), and Theorem~2.1]{BiskupLouidor2018}, together
with its stated consequence for maxima on open subsets following
(1.8), gives tightness of the maximum on a fixed inner subdisk after
centering by
\[
 m_M=2\sqrt g\log M-\frac34\sqrt g\log\log M.
\]
The digital disks here approximate a disk, and $V_y$ is a fixed
nonempty subdisk strictly inside it.  In particular,
\begin{equation}
 \max_{x\in V_y}\phi_x=2\sqrt g\,h+O_P(\log h).
 \label{eq:terminal-dgff-maximum}
\end{equation}
Only this localization, rather than the full limiting point process,
is used below.

To pass to expectations, convergence in distribution alone is
insufficient.  The tightness just stated places every median of the
maximum at $m_M+O(1)$.  Its maximal coordinate variance is $O(h)$,
so Gaussian concentration places the mean within $O(\sqrt h)$ of
a median.  Therefore
\begin{equation}
 E\max_{x\in V_y}\phi_x=2\sqrt g\,h+O(\sqrt h).
 \label{eq:terminal-dgff-mean}
\end{equation}

We next check that the bounded covariance error in
\eqref{eq:terminal-covariance-comparison} changes the expected
maximum by at most $O(\sqrt h)$.  For $\lambda>0$, define
\[
 F_\lambda(z)=\lambda^{-1}\log\sum_{x\in V_y}e^{\lambda z_x}.
\]
If $p_x=e^{\lambda z_x}/\sum_{x'}e^{\lambda z_{x'}}$, then
\[
 \partial_{xx'}F_\lambda
 =\lambda(p_x\mathbf1_{\{x=x'\}}-p_xp_{x'}),\qquad
 \sum_{x,x'}|\partial_{xx'}F_\lambda|\le2\lambda.
\]
Gaussian interpolation between $\Gamma$ and $\sqrt{2g}\,\phi$
therefore bounds the difference of their expected $F_\lambda$
values by $C\lambda$, using
\eqref{eq:terminal-covariance-comparison}.  Since
\[
 0\le F_\lambda(z)-\max_xz_x\le\frac{\log P}{\lambda},
\]
choosing $\lambda\asymp\sqrt{\log P}$ yields
\begin{equation}
 \left|E\max_x\Gamma_x-E\max_x\sqrt{2g}\,\phi_x\right|
 \le C\sqrt{\log P}.
 \label{eq:terminal-gaussian-mean-comparison}
\end{equation}
Both fields have maximal coordinate variance $O(h)$.  Combining
Gaussian concentration with
\eqref{eq:terminal-dgff-mean}--\eqref{eq:terminal-gaussian-mean-comparison}
gives
\begin{equation}
 \max_x\Gamma_x=2\sqrt2\,g h+O_P(\sqrt h+\log h)
             =2\sqrt2\,g h+O_P(\sqrt h).
 \label{eq:terminal-gaussian-maximum}
\end{equation}

Apply \eqref{eq:terminal-cck} to each of the two half-lines.  For
any $K>0$, the probability that $\max_xU_x$ is below
$2\sqrt2\,gh-K\sqrt h$ is bounded by the corresponding Gaussian
lower-tail probability with its threshold increased by $C_7$, plus
the error in \eqref{eq:terminal-cck-error}.  The upper tail is bounded
by the Gaussian upper tail with its threshold decreased by $C_7$,
plus the same error.  Equation~\eqref{eq:terminal-gaussian-maximum}
thus transfers to $\max_xU_x$, uniformly in the regular endpoint
data.  This argument compares distributions; it requires no
pathwise coupling of the maxima.

\smallskip\noindent
\underline{\textit{Step 4.}}\quad Restoring the means and the terminal depth.
Equation~\eqref{eq:terminal-conditional-mean} implies
\[
 \sum_{i\le b}E^{\omega_i}X_i(x)
 =gb+O\left(\frac{J\log J}{\log\log J}\right)
\]
uniformly in $x$.  Multiplying the centered maximum estimate by
$\sqrt b=O(J)$ proves \eqref{eq:terminal-maximum-expansion}.

Finally, \eqref{eq:terminal-piece-count} gives
\begin{equation}
 \begin{split}
 gb+2\sqrt2\,g\sqrt b\,h
 &=2gJ^2-2gJ\log J+O(h^2+h\log J+J).
 \end{split}
 \label{eq:terminal-restored-depth}
\end{equation}
For example, this follows by first using
$\sqrt b=\sqrt2L-(\log L)/\sqrt2+O(1+(\log L)^2/L)$ and
then substituting $L=J-h$.  The $4gLh$ term from the Gaussian
maximum restores the missing terminal depth in $2gL^2$.
Since $p<2$,
\[
 J\sqrt h+\frac{J\log J}{\log\log J}
 +h^2+h\log J+J=o(J\log J).
\]
The level in \eqref{eq:terminal-restored-depth} exceeds
$u_{J,2+\eta}$ by $\eta gJ\log J+o(J\log J)$, which proves
the final assertion.
\end{proof}

\subsection{A positive probability of a large disk maximum}
\label{sec:disk-lower}

The terminal estimate is conditional on the coarse data of one
fixed cell.  To combine it with the existence of a coarse root, we
sum the failures over all candidates.  Selecting a successful coarse
root first would introduce additional conditioning on bridge
interiors from other cells, which the conditional estimate does not
cover.

\begin{proposition}[The disk maximum]
\label{prop:disk-maximum-lower}
For every fixed $\eta>0$, there is $p_\eta>0$ such that, for all
sufficiently large $J$,
\begin{equation}
 P\left(M_{\tau_{e^J}}\ge u_{J,2+\eta}\right)\ge p_\eta.
 \label{eq:disk-maximum-positive-probability}
\end{equation}
\end{proposition}

\begin{proof}
Take $h=(\log J)^p$, $1<p<2$, and $L=J-h$.  Apply the lower-root
construction \eqref{eq:ct-external-32}--\eqref{eq:ct-external-38} at depth $L$, with terminal physical radius
$M=e^h$.  Restrict the centers to a fixed compact subdisk so that
their enlarged terminal domains are contained in $D(0,e^J)$.
This keeps a fixed positive fraction of the available area, and the
estimates are uniform on the restricted set.  The endpoint band in
\eqref{eq:ct-external-37} is
\[
 \sqrt{2T_L}\in[\rho_LL+z_0,\rho_LL+z_0+C_*],
\]
which gives the completed-piece count
\eqref{eq:terminal-piece-count}.

Lemma~\ref{lem:buffered-transfer} transfers the probability of
existence of a coarse root directly from Brownian motion; the pair
estimate is used on the Brownian side.  At each root, the terminal
pieces are ordered separately within the $m_0$ source intervals.
The inter-source connectors are exposed and omitted from the
occupation sums.  Their occupation is nonnegative, so omitting it
preserves a lower bound for the full-walk maximum.
For the lattice coarse-root count $Z_L^S$, the count assertion of
that lemma gives
\begin{equation}
 P(Z_L^S>0)\ge p_0,\qquad EZ_L^S\le C.
 \label{eq:disk-root-count}
\end{equation}

We must also show that the moment screen does not remove every
root.  Let $\widehat Z_L^S$ count the same radial barriers and
endpoint bands with all angular screens dropped.  The scalar
radial-word comparison \eqref{eq:ct-smooth-39d} and Rosen's Appendix Theorem~9.1(a)
give
\begin{equation}
 P(\widehat I_y^{S,\circ})\le Ce^{-2L},\qquad
 E\widehat Z_L^S\le C,
 \label{eq:disk-unscreened-count}
\end{equation}
because the deterministic net has $O(e^{2L})$ points.  Conditional
on each complete radial word counted by $\widehat Z_L^S$,
\eqref{eq:terminal-moment-failure} bounds failure of moment
regularity by the same $\epsilon_J=o(J^{-A})$.  Sum first over
words and then over centers to obtain
\begin{equation}
 E\sum_y\mathbf1_{I_y^{S,\circ}}
       \mathbf1_{\{\text{the endpoints at }y\text{ are not moment regular}\}}
 \le\epsilon_J E\widehat Z_L^S=o(1).
 \label{eq:disk-irregular-roots}
\end{equation}
Using the unscreened count here is essential: the estimate for the
endpoints is conditional on a radial word, so its failure probability
must be integrated against the radial-word mass.  Equations
\eqref{eq:disk-root-count} and \eqref{eq:disk-irregular-roots}
show that a moment-regular lattice root exists with probability
at least $p_0-o(1)$.

For a fixed root $y$, erase its completed pieces from the path and
let $\mathcal C_y$ record the exterior path and every entrance and
directed-exit pair.  The event $G_y$ that $y$ is a moment-regular
coarse root is $\mathcal C_y$-measurable.  Given this sigma-field,
the erased pieces have the independent killed bridge laws used in
Lemma~\ref{lem:terminal-maximum-lower}.  Let
\[
 A_y=\left\{\max_{x\in V_y}\sum_{i\le b}X_i(x)
                       \ge u_{J,2+\eta}\right\}.
\]
The uniform conclusion of that lemma gives a deterministic
$\widetilde\epsilon_J\downarrow0$ such that
\begin{equation}
 P(G_y\cap A_y^c)
 =E[\mathbf1_{G_y}P(A_y^c\mid\mathcal C_y)]
 \le\widetilde\epsilon_JP(G_y).
 \label{eq:disk-fixed-root-failure}
\end{equation}
Consequently
\begin{equation}
 \begin{split}
 P(\exists y:G_y\cap A_y)
 &\ge P(\exists y:G_y)-\sum_yP(G_y\cap A_y^c)\\
 &\ge p_0-o(1)-\widetilde\epsilon_J EZ_L^S
 \ge p_0/2
 \end{split}
 \label{eq:disk-successful-root}
\end{equation}
for all sufficiently large $J$.  Each summand uses the intrinsic
bridge disintegration for its own cell; the union bound does not
require disjoint terminal disks.  On the event in
\eqref{eq:disk-successful-root}, the full local-time maximum is at
least the terminal maximum of a successful cell.  This proves
\eqref{eq:disk-maximum-positive-probability}.
\end{proof}

\subsection{An almost-sure lower bound at deterministic times}
\label{sec:deterministic-lower}

The disk estimate has a probability bounded away from zero.  We
amplify it using independent segments of increments.  The number of
segments may grow as a small power of $\log N$; shortening each
segment then costs only a controlled amount in the second-order
term of the maximum.

\begin{proposition}[The deterministic-time maximum]
\label{prop:deterministic-maximum-lower}
Let $N_j=\lfloor e^{2j}\rfloor$.  For every $\eta>0$, almost surely,
\[
 M_{N_j}\ge u_{j,2+\eta}
\]
for all sufficiently large $j$.
\end{proposition}

\begin{proof}
Fix $0<\delta<\eta$, and then choose $a>0$ so that
$\delta+2a<\eta$.  Proposition~\ref{prop:disk-maximum-lower}
gives a disk-success probability $p_\delta>0$.  By
Lemma~\ref{lem:elementary-walk-estimates}, choose a fixed $C$
large enough that
\[
 \sup_R P_0(\tau_R>CR^2)\le p_\delta/2.
\]
It follows by translation invariance that a walk segment of length
$CR^2$, from any starting site, has within-segment maximum local
time at least $u_{\log R,2+\delta}$ with probability at least
$p=p_\delta/2$, uniformly for all sufficiently large $R$.

Partition the first $N_j$ increments into
\[
 q_j=\lfloor(\log N_j)^a\rfloor
\]
disjoint blocks of length $m_j=\lfloor N_j/q_j\rfloor$, ignoring
the remainder.  Use the sites immediately after the increments in
each block to define its within-block local times.  The translated
block paths are independent, and each block maximum is bounded
above by the full local-time maximum.  Bounded endpoint changes
are absorbed by the strict margin in $\delta+2a<\eta$.
Choose $R_j$ by
\[
 \log R_j=\frac12\log(m_j/C)
         =j-\frac a2\log j+O(1).
\]
The loss in logarithmic radius is therefore $(a/2)\log j+O(1)$.
Substituting this into the quadratic leading term gives
\[
 2g(\log R_j)^2
 =2gj^2-2agj\log j+O(j+(\log j)^2),
\]
while
$g(\log R_j)\log\log R_j=gj\log j+O(j+(\log j)^2)$.
Hence
\begin{equation}
 u_{\log R_j,2+\delta}
 =u_{j,\,2+\delta+2a+o(1)}
 \ge u_{j,2+\eta}
 \label{eq:deterministic-block-level}
\end{equation}
for all sufficiently large $j$.  Independence of the blocks now
implies
\begin{equation}
 P(M_{N_j}<u_{j,2+\eta})
 \le(1-p)^{q_j}\le e^{-cj^a}.
 \label{eq:deterministic-failure-summable}
\end{equation}
The right-hand side is summable, so Borel--Cantelli proves the
assertion.
\end{proof}


% ---- escape ----
\section{Escape below the critical scale}\label{sec:escape}
The localized disk estimate must exclude a favorite at every time in a block,
not just at its endpoints. Monotonicity of local times makes this possible:
use the global maximum at the left endpoint and the local time in the target
disk at the right endpoint.

\begin{proposition}[Escape below the critical scale]\label{prop:supercritical-escape}
For every \(\gamma>\tfrac12\), almost surely,
\begin{equation}
\liminf_{n\to\infty}
       \frac{R_n^-}{\sqrt n/(\log n)^\gamma}=\infty.          \label{eq:escape-80}
\end{equation}
\end{proposition}

\begin{proof}
Fix \(\gamma>1/2\) and a positive integer spatial multiplier \(A_0\).
Choose
\[
       0<\delta<\frac{2\gamma-1}{4}.
\]
For \(N_j=e^{2j}\), Proposition~\ref{prop:deterministic-maximum-lower} gives,
eventually,
\begin{equation}
M_{N_j}\ge u_{j,2+\delta}.                             \label{eq:escape-81}
\end{equation}

Consider any \(n\in[N_j,N_{j+1}]\).  If a favorite at time \(n\) lies in
\[
       B\left(0,A_0\frac{\sqrt n}{(\log n)^\gamma}\right),
\]
then, by monotonicity of local times and \eqref{eq:escape-81}, at time \(N_{j+1}\) some
site in
\begin{equation}
B\left(0,C_{A_0}e^{j+1}j^{-\gamma}\right)             \label{eq:escape-82}
\end{equation}
has local time at least \(u_{j,2+\delta}\).

Take
\[
       K_j=C_*e^{j+1}\sqrt{\log j},\qquad J'_j=\log K_j.
\]
By the maximal inequality in Lemma~\ref{lem:elementary-walk-estimates}, \(C_*\)
can be chosen so that
\begin{equation}
\sum_jP(\tau_{K_j}\le N_{j+1})<\infty.                 \label{eq:escape-83}
\end{equation}
For large \(j\), the ball \eqref{eq:escape-82} is contained in
\[
       B(0,K_j(J'_j)^{-\gamma}).
\]
Moreover \(J'_j=j+O(\log\log j)\), and the fixed extra coefficient
\(\delta\) gives
\begin{equation}
u_{j,2+\delta}\ge u_{J'_j,2+2\delta}                  \label{eq:escape-84}
\end{equation}
for all large \(j\): the negative term
\(-\delta g j\log j\) dominates the \(O(j\log\log j)\) change in the
leading term.

Choose independently an upper-error exponent \(A_{\rm err}>2\).
On the complement of the summable event in \eqref{eq:escape-83}, the implication
\eqref{eq:escape-81}--\eqref{eq:escape-84} and Proposition~\ref{prop:localized-disk-upper}, with
\(\beta=\gamma\) and \(c=2+2\delta\), show that the probability of any
such \(n\) in the whole block is at most
\begin{equation}
(J'_j)^{-2\gamma+2\delta+o(1)}+O(j^{-A_{\rm err}}).     \label{eq:escape-85}
\end{equation}
Our choice of \(\delta\) makes the power in \eqref{eq:escape-85} strictly smaller than
\(-1\).  Hence \eqref{eq:escape-85} is summable.  Borel--Cantelli proves that eventually
no favorite in the \(j\)-th block lies in the ball with multiplier \(A_0\).
Intersecting over \(A_0\in\mathbb N\) proves \eqref{eq:escape-80}.
\end{proof}


\begin{proof}[Proof of Theorem~\ref{main:precise-escape-rate}]
Propositions~\ref{prop:critical-recurrence} and
\ref{prop:supercritical-escape} give the two assertions for each fixed
exponent. Intersect the probability-one events for rational
$\gamma>1/2$ and for the single critical exponent $1/2$.
If $\gamma>1/2$ is real, choose a rational $q\in(1/2,\gamma)$; then
\[
 \frac{R_n^-}{a_\gamma(n)}
 =\frac{R_n^-}{a_q(n)}(\log n)^{\gamma-q}\longrightarrow\infty.
\]
For $\gamma\le1/2$, the inequality
$R_n^+/a_\gamma(n)\le R_n^+/a_{1/2}(n)$ gives the other assertion on the
same event. This proves the simultaneous formulation.
\end{proof}


% ---- notation ----
\clearpage
\appendix
\section{Notation guide}\label{sec:notation}

The same letter $n$ denotes physical time in the main theorem and a
truncated logarithmic depth inside the traversal argument; the relation
$n=J-h+O(1)$ is stated whenever the latter convention is used. Likewise,
the function $h(r)=2\arctan(r/2)$ converts conformal to spherical radii,
whereas the terminal depth $h$ satisfies $M=e^h$.

\small
\begin{longtable}{@{}p{.19\textwidth}p{.49\textwidth}p{.25\textwidth}@{}}
\toprule
Symbol & Meaning & Where used \\
\midrule
\endhead
$L_n(x),M_n$ & Site local time and its spatial maximum. & Section~\ref{sec:introduction} \\
$\mathcal F_n,R_n^-,R_n^+$ & Favorite set and its nearest and farthest distances from the origin. & Section~\ref{sec:introduction} \\
$a_\gamma(n)$ & $\sqrt n/(\log n)^\gamma$. & Theorem~\ref{main:precise-escape-rate} \\
$g,u_{J,c}$ & $g=2/\pi$ and $u_{J,c}=2gJ^2-cgJ\log J$. & Both disk estimates \\
$D(x,r),\tau_r$ & Lattice disk and exit time from $D(0,r)$. & Section~\ref{sec:elementary} \\
$G_D$ & Green function for the walk killed on leaving $D$. & Sections~\ref{sec:elementary} and \ref{sec:terminal-upper} \\
$z,N,K$ & Geometric weight, $N=(1-z)^{-1}$ and clock with $P(K=n)=(1-z)z^n$. & Section~\ref{sec:recurrence} \\
$X_K^*,U_x$ & Priority-selected favorite and independent site priorities. & Lemma~\ref{lem:plateau-bijection} \\
$\sigma,\ell$ & First and last visits to the selected favorite. & Lemma~\ref{lem:plateau-bijection} \\
$\mathcal B,\tau_c$ & Exposed excursion/suffix tuple and terminal-arm killing time. & Lemma~\ref{lem:plateau-bijection} \\
$J,h,M,L$ & Outer logarithmic radius, terminal logarithmic radius, $M=e^h$, and lower traversal depth $L=J-h$. & Sections~\ref{sec:continuum-transfer}--\ref{sec:terminal-lower} \\
$\rho_L,\alpha_{z,\pm}$ & Barrier slope $2-(\log L)/L$ and its upper/lower curved barriers. & Section~\ref{sec:continuum-transfer} \\
$F_L,k_y$ & Deterministic net of centers and a center's source shell. & Section~\ref{sec:continuum-transfer} \\
$T_{y,j}^{1,m_0}$ & Traversal count within the first $m_0$ complete source excursions. & Section~\ref{sec:continuum-transfer} \\
$I_y^\circ,S_y$ & Truncated Brownian root event and connector safety event. & Section~\ref{sec:continuum-transfer} \\
$\chi_y,\Psi_J$ & Centered conformal coordinate and dilated stereographic chart. & Section~\ref{sec:continuum-transfer} \\
$w,Q,G_0$ & Complete radial word, number of shell decisions, and number of nonempty source blocks. & Section~\ref{sec:terminal-upper} \\
$D_y,A_y,V_y,E_y$ & Terminal killing disk, entrance interface, core, and directed exit edges. & Section~\ref{sec:terminal-upper} \\
$\omega_i,X_i(x)$ & Entrance/exit pair and occupation at $x$ of its killed bridge. & Section~\ref{sec:terminal-upper} \\
$P_w,\widetilde P_w$ & True complete-word endpoint law and its sequentially normalized comparison law. & Section~\ref{sec:terminal-upper} \\
$\Lambda_M$ & Reference entrance/exit measure obtained from harmonic measure at the center. & Section~\ref{sec:terminal-upper} \\
$\mathcal R_y,\mathcal C_y$ & Intrinsic information exposing exterior paths and bridge endpoints while erasing bridge interiors. & Sections~\ref{sec:continuum-transfer} and \ref{sec:terminal-lower} \\
$Z_L^S,\widehat Z_L^S$ & Lattice root counts used before conditional terminal success. & Section~\ref{sec:terminal-lower} \\
\bottomrule
\end{longtable}
\normalsize


% ---- references ----
\clearpage
\begin{thebibliography}{9}
\raggedright

\bibitem{AthreyaEtAl2022}
S. Athreya, N. Gadhiwala and R. R. Radhakrishnan,
\emph{Elliptic Harnack inequality for $\mathbb Z^d$},
Involve \textbf{15} (2022), 687--708.
\href{https://doi.org/10.2140/involve.2022.15.687}{doi:10.2140/involve.2022.15.687}.
\href{https://arxiv.org/abs/2201.11947}{arXiv version}.

\bibitem{Bass2023}
R. F. Bass,
\emph{The rate of escape of the most visited site of Brownian motion},
Electron. J. Probab. \textbf{28} (2023), 1--12.
\href{https://doi.org/10.1214/23-EJP916}{doi:10.1214/23-EJP916}.
See also \href{https://arxiv.org/abs/1303.2040v5}{arXiv:1303.2040v5},
Theorem~2.2.

\bibitem{BiskupLouidor2018}
M. Biskup and O. Louidor,
\emph{Full extremal process, cluster law and freezing for the two-dimensional
discrete Gaussian free field},
Adv. Math. \textbf{330} (2018), 589--687.
\href{https://doi.org/10.1016/j.aim.2018.02.018}{doi:10.1016/j.aim.2018.02.018}.
\href{https://arxiv.org/abs/1606.00510}{arXiv:1606.00510}.

\bibitem{CCK2016}
V. Chernozhukov, D. Chetverikov and K. Kato,
\emph{Empirical and multiplier bootstraps for suprema of empirical processes
of increasing complexity, and related Gaussian couplings},
Stochastic Process. Appl. \textbf{126} (2016), 3632--3651.
\href{https://doi.org/10.1016/j.spa.2016.04.009}{doi:10.1016/j.spa.2016.04.009}.
\href{https://arxiv.org/abs/1502.00352}{arXiv:1502.00352}.

\bibitem{HLOZ2025}
C. Hao, X. Li, I. Okada and Y. Zheng,
\emph{Favorite sites for simple random walk in two and more dimensions},
Probab. Theory Related Fields (2025).
\href{https://doi.org/10.1007/s00440-025-01441-1}{doi:10.1007/s00440-025-01441-1}.
\href{https://arxiv.org/abs/2409.00995}{arXiv:2409.00995}.
The escape-rate question cited here is Remark~1.4(1) in
\href{https://arxiv.org/abs/2409.00995v1}{version~1}; the elementary
walk estimates are in Lemma~2.1, (2.5), also in version~2.

\bibitem{KozdronLawler2005}
M. J. Kozdron and G. F. Lawler,
\emph{Estimates of random walk exit probabilities and application to loop-erased
random walk},
Electron. J. Probab. \textbf{10} (2005), 1442--1467.
\href{https://arxiv.org/abs/math/0501189}{arXiv version}.

\bibitem{Rosen2022}
J. Rosen,
\emph{Tightness for thick points in two dimensions},
arXiv:2203.14394 (2022).
\href{https://arxiv.org/abs/2203.14394}{arXiv version}.

\bibitem{Zaitsev2003}
A. Yu. Zaitsev,
\emph{Estimates for the strong approximation in multidimensional central
limit theorem},
arXiv:math/0304373 (2003).
\href{https://arxiv.org/abs/math/0304373}{arXiv version}.

\end{thebibliography}
\end{document}
